% Leçon 8 — Expression écrite : la délinquance juvénile — Chinois 1ère A — Cours signé OnBuch+
% Programme officiel MINESEC. Compiler avec : tectonic ZHA08-ecrit-la-delinquance-juvenile.tex
\newcommand{\DOCMATIERE}{Chinois}
\newcommand{\DOCNIVEAU}{1ère A}
\newcommand{\DOCMODULE}{Module 1 : Vie familiale et intégration sociale}
\newcommand{\DOCLECON}{ZHA08}
\newcommand{\DOCTITRE}{Leçon 8 — Expression écrite : la délinquance juvénile}
% OnBuch+ — préambule « cours signés » ENRICHI (compilé avec tectonic / XeLaTeX)
% Système visuel « Pop » d'OnBuch+ : fond crème (jamais blanc), bordures encre
% épaisses, ombres « dures » décalées, titres Archivo Black en capitales.
% Version MATHÉMATIQUES : graphes (pgfplots/tikz), boîtes pédagogiques
% (définition, propriété, méthode, attention, le savais-tu, exercice résolu,
% exercice, TP, bilan), page de garde et signature OnBuch+ ancrée en bas.
%
% Macros attendues dans main.tex AVANT \input{preamble} :
%   \DOCMATIERE  \DOCNIVEAU  \DOCMODULE  \DOCLECON (numéro)  \DOCTITRE
\documentclass[11pt]{article}

\usepackage{fontspec}
\usepackage[french]{babel} % français = langue principale ; le chinois via \zh{..} / textebox (xeCJK)
\usepackage{xeCJK}
\usepackage{pifont} % \ding{51} \ding{55} utilisés par les modèles
\usepackage[a4paper,margin=2cm,top=3.4cm,bottom=2.8cm,headheight=1.4cm,footskip=1.3cm]{geometry}
\usepackage{amsmath,amssymb,amsfonts}
\usepackage{mathtools} % \xmapsto, \coloneqq... souvent utilisés par les modèles
\usepackage{cancel} % simplifications barrées (bilans de forces)
% \abs{z} (module, valeur absolue) et \norm{u} (norme), utilisés par les modèles
\providecommand{\abs}{}\let\abs\relax\DeclarePairedDelimiter{\abs}{\lvert}{\rvert}
\providecommand{\norm}{}\let\norm\relax\DeclarePairedDelimiter{\norm}{\lVert}{\rVert}
\usepackage[table]{xcolor} % [table] : fournit aussi \rowcolors (alternance de lignes)
\usepackage{pagecolor}
\usepackage{tikz}
\usetikzlibrary{arrows.meta,positioning,calc,shapes.geometric,shapes.misc,shapes.arrows,decorations.pathmorphing,decorations.pathreplacing,decorations.markings,patterns,shadows,mindmap,trees,fit,backgrounds,matrix,intersections,angles,quotes}
\usepackage{pgfplots}
\usepackage[version=4]{mhchem} % équations nucléaires \ce{^{235}_{92}U}
\usepackage[european,siunitx]{circuitikz} % schémas électriques
\pgfplotsset{compat=1.18}
\usepgfplotslibrary{fillbetween,groupplots} % aires entre courbes (calcul intégral)
\usepackage{enumitem}
\usepackage{fancyhdr}
% [most] charge toutes les bibliothèques tcolorbox (listings, minted, raster,
% documentation...) et gonfle inutilement la mémoire TeX sur un document long
% (~300 boîtes) : on ne charge que ce qui est réellement utilisé.
\usepackage[skins,breakable]{tcolorbox}
\usepackage{graphicx}
\usepackage{colortbl}
\usepackage{booktabs}
\usepackage{tabularx}
\usepackage{diagbox} % case d'en-tête diagonale des tableaux à double entrée
% Colonnes extensibles centrée / gauche / droite (C, L, R), souvent utilisées par les modèles
\newcolumntype{C}{>{\centering\arraybackslash}X}
\newcolumntype{L}{>{\raggedright\arraybackslash}X}
\newcolumntype{R}{>{\raggedleft\arraybackslash}X}
\usepackage{array}
\usepackage{multicol}
\usepackage{multirow}
\usepackage{siunitx}
% parse-numbers=false : accepte aussi \SI{2,5 \times 10^{-3}}{...}
\sisetup{locale=FR,output-decimal-marker={,},group-minimum-digits=4,parse-numbers=false}
\DeclareSIUnit\ohm{\ensuremath{\symup{\Omega}}} % U+2126 absent de la police
\DeclareSIUnit\division{div} % réglages d'oscilloscope (ms/div, V/div)
\DeclareSIUnit\tr{tr} % tours (vitesse de rotation, ex. \SI{1500}{\tr\per\minute})
\DeclareSIUnit\molar{M} % concentration molaire (ex. \SI{3}{\molar}, \SI{1}{\micro\molar})
\DeclareSIUnit\an{an} % année (le modèle confond parfois avec \year, un compteur TeX) — ex. \SI{5}{\kilo\meter\per\an}
\DeclareSIUnit\FCFA{FCFA} % franc CFA (ex. \SI{500}{\FCFA\per\kilo\gram})
\DeclareSIUnit\degreeBrix{\textdegree Brix} % teneur en sucre d'un sirop/jus (réfractométrie)
\DeclareSIUnit\month{mois} % le modèle utilise parfois \month, un compteur TeX, comme unité de durée
\usepackage{qrcode}
% \degree : commande fréquemment utilisée par les modèles pour un angle ou une
% température en dehors de \SI{}{} (non fournie par les paquets chargés).
\providecommand{\degree}{^{\circ}}
% \qed (fin de démonstration), utilisé par les modèles sans amsthm
\providecommand{\qed}{\hfill$\square$}
% \barre : double barre des valeurs interdites dans un tableau de variations (tkz-tab)
\providecommand{\barre}{\ensuremath{\|}}
% environnement proof (amsthm non chargé) utilisé par les modèles
\ifdefined\proof\else\newenvironment{proof}[1][Démonstration]{\par\noindent\textit{#1.}\ }{\qed\par}\fi
\usepackage{lastpage}
\usepackage{hyperref}
\hypersetup{hidelinks}

% ============================================================
% Palette « Pop » OnBuch+ — mêmes hex que la classe P (plus_ui.dart)
% ============================================================
\definecolor{poporange}{HTML}{F59321}
\definecolor{popdark}{HTML}{E07A0C}
\definecolor{popcream}{HTML}{FFF6E8}
\definecolor{popink}{HTML}{1C1714}
\definecolor{poppurple}{HTML}{7A5AE0}
\definecolor{popgreen}{HTML}{2E9E6A}
\definecolor{poppink}{HTML}{E0457B}
\definecolor{popblue}{HTML}{2F7DE1}
\definecolor{popgold}{HTML}{C98A12}
\definecolor{popmuted}{HTML}{8A7F76}
\definecolor{popline}{HTML}{EADFD2}
\definecolor{popgray}{HTML}{8A7F76}
\colorlet{popgrey}{popgray}
% Alias tolérants (le modèle nomme parfois les couleurs différemment)
\colorlet{popwhite}{white}
\colorlet{popblack}{popink}
\colorlet{popred}{poppink}
\colorlet{popyellow}{popgold}
% Teintes claires (fonds de tableaux/encadrés) — le modèle nomme parfois
% les couleurs avec un suffixe L (ex. popblueL) sans les avoir définies.
\colorlet{poporangeL}{poporange!20}
\colorlet{popdarkL}{popdark!20}
\colorlet{poppurpleL}{poppurple!20}
\colorlet{popgreenL}{popgreen!20}
\colorlet{poppinkL}{poppink!20}
\colorlet{popblueL}{popblue!20}
\colorlet{popgoldL}{popgold!20}
\colorlet{popredL}{popred!20}
\colorlet{popgrayL}{popgray!20}
% Teintes claires pour fonds de tableaux / zones
\colorlet{popblueL}{popblue!10!white}
\colorlet{poporangeL}{poporange!14!white}
\colorlet{popgreenL}{popgreen!12!white}
\colorlet{poppurpleL}{poppurple!12!white}
\colorlet{poppinkL}{poppink!12!white}
\colorlet{popgoldL}{popgold!14!white}

\pagecolor{popcream}

% ============================================================
% Typographie
% ============================================================
\defaultfontfeatures{Ligatures=TeX}
\setmainfont{PlusJakartaSans-Medium.ttf}[Path=../../../fonts/,BoldFont=PlusJakartaSans-Bold.ttf]
% Symboles, API et emojis absents de la police principale : repli sur DejaVu Sans (caractères actifs)
\newfontface\symfont{DejaVuSans.ttf}[Path=../../../fonts/]
\makeatletter
\newcommand\symactive[1]{\catcode#1=13 \begingroup\lccode`\~=#1 \lowercase{\endgroup\def~}{{\symfont\char#1}}}
\newcommand\symblank[1]{\catcode#1=13 \begingroup\lccode`\~=#1 \lowercase{\endgroup\def~}{}}
\newcommand\symhyph[1]{\catcode#1=13 \begingroup\lccode`\~=#1 \lowercase{\endgroup\def~}{\mbox{-}}}
\newcount\symc
\newcommand\symrange[2]{\symc=#1 \@whilenum\symc<#2 \do{\expandafter\symactive\expandafter{\the\symc}\advance\symc by 1 }}
\newcommand\symblankrange[2]{\symc=#1 \@whilenum\symc<#2 \do{\expandafter\symblank\expandafter{\the\symc}\advance\symc by 1 }}
\makeatother
\symrange{"0250}{"02FF}\symactive{"2713}\symactive{"2714}\symactive{"2717}\symactive{"2718}\symactive{"2610}\symactive{"2611}\symactive{"2612}
\symrange{"2600}{"27BF}
\catcode"2705=13 \begingroup\lccode`\~="2705 \lowercase{\endgroup\def~}{{\symfont\char"2713}}\symblank{"26BD}\symblank{"26A1}
\symrange{"2460}{"257F}\symactive{"223C}\symactive{"2248}\symactive{"2260}
\symhyph{"2011}\symhyph{"2E1A}\symblankrange{"1F300}{"1FAFF}\symblank{"FE0F}
% Caractères chinois : WenQuanYi Zen Hei (sous-ensemble GB2312 + pinyin), gras/italique simulés
\setCJKmainfont{WenQuanYiZenHei-subset.ttf}[Path=../../../fonts/,AutoFakeBold=3,AutoFakeSlant=0.2]
\xeCJKsetup{CJKspace=false}
% Pinyin : voyelles à caron (ǎ ǐ ǒ ǔ ǖ ǘ ǚ ǜ) absentes de la police principale -> caractères actifs
\newfontface\pyfont{WenQuanYiZenHei-subset.ttf}[Path=../../../fonts/]
\newcommand\pyactive[1]{\catcode#1=13 \begingroup\lccode`\~=#1 \lowercase{\endgroup\def~}{{\pyfont\char#1}}}
\pyactive{"01CD}\pyactive{"01CE}\pyactive{"01CF}\pyactive{"01D0}\pyactive{"01D1}\pyactive{"01D2}\pyactive{"01D3}\pyactive{"01D4}\pyactive{"01D5}\pyactive{"01D6}\pyactive{"01D7}\pyactive{"01D8}\pyactive{"01D9}\pyactive{"01DA}\pyactive{"01DB}\pyactive{"01DC}
% Maths en sans-serif (Fira Math) avec chiffres et lettres droites de Plus
% Jakarta, pour que les formules se fondent dans le texte.
\usepackage[math-style=ISO,bold-style=ISO]{unicode-math}
\setmathfont{FiraMath-Regular.otf}[Path=../../../fonts/]
\setmathfont{PlusJakartaSans-Medium.ttf}[Path=../../../fonts/,range={up/num,up/{Latin,latin},\mathup}]
\usepackage{icomma} % virgule décimale sans espace en mode math
% Caractères mathématiques Unicode absents des polices de texte (Plus Jakarta,
% Archivo Black) : rendus par leur équivalent mathématique, en texte comme en
% formule — sinon ils s'affichent en carré vide (ex. « Arithmétique dans ℤ »).
\usepackage{newunicodechar}
\newunicodechar{ℝ}{\ensuremath{\mathbb{R}}}
\newunicodechar{ℤ}{\ensuremath{\mathbb{Z}}}
\newunicodechar{ℂ}{\ensuremath{\mathbb{C}}}
\newunicodechar{ℕ}{\ensuremath{\mathbb{N}}}
\newunicodechar{ℚ}{\ensuremath{\mathbb{Q}}}
\newunicodechar{𝔻}{\ensuremath{\mathbb{D}}}
\newunicodechar{α}{\ensuremath{\alpha}}
\newunicodechar{β}{\ensuremath{\beta}}
\newunicodechar{γ}{\ensuremath{\gamma}}
\newunicodechar{δ}{\ensuremath{\delta}}
\newunicodechar{ε}{\ensuremath{\varepsilon}}
\newunicodechar{θ}{\ensuremath{\theta}}
\newunicodechar{λ}{\ensuremath{\lambda}}
\newunicodechar{μ}{\ensuremath{\mu}}
\newunicodechar{σ}{\ensuremath{\sigma}}
\newunicodechar{τ}{\ensuremath{\tau}}
\newunicodechar{φ}{\ensuremath{\varphi}}
\newunicodechar{ψ}{\ensuremath{\psi}}
\newunicodechar{ω}{\ensuremath{\omega}}
\newunicodechar{Σ}{\ensuremath{\Sigma}}
% majuscules produites par \MakeUppercase dans les titres (α-aminés -> Α)
\newunicodechar{Α}{\ensuremath{\alpha}}
\newunicodechar{Β}{\ensuremath{\beta}}
\newunicodechar{Γ}{\ensuremath{\Gamma}}
\newunicodechar{Δ}{\ensuremath{\Delta}}
\newunicodechar{Ε}{\ensuremath{\varepsilon}}
\newunicodechar{Θ}{\ensuremath{\Theta}}
\newunicodechar{Λ}{\ensuremath{\Lambda}}
\newunicodechar{Μ}{\ensuremath{\mu}}
\newunicodechar{Τ}{\ensuremath{\tau}}
\newunicodechar{Φ}{\ensuremath{\Phi}}
\newunicodechar{Ψ}{\ensuremath{\Psi}}
\newunicodechar{Ω}{\ensuremath{\Omega}}
\newunicodechar{↦}{\ensuremath{\mapsto}}
\newunicodechar{∈}{\ensuremath{\in}}
\newunicodechar{∉}{\ensuremath{\notin}}
\newunicodechar{∧}{\ensuremath{\wedge}}
\newunicodechar{⇔}{\ensuremath{\Leftrightarrow}}
\newunicodechar{⟺}{\ensuremath{\Longleftrightarrow}}
\newunicodechar{⇒}{\ensuremath{\Rightarrow}}
\newunicodechar{⟹}{\ensuremath{\Longrightarrow}}
\newunicodechar{∘}{\ensuremath{\circ}}
\newunicodechar{∪}{\ensuremath{\cup}}
\newunicodechar{∩}{\ensuremath{\cap}}
\newunicodechar{≡}{\ensuremath{\equiv}}
\newunicodechar{⊂}{\ensuremath{\subset}}
\newunicodechar{⊥}{\ensuremath{\perp}}
\newunicodechar{∃}{\ensuremath{\exists}}
\newunicodechar{∀}{\ensuremath{\forall}}
\newunicodechar{∛}{\ensuremath{\sqrt[3]{\,}}}
\newunicodechar{∜}{\ensuremath{\sqrt[4]{\,}}}
\newunicodechar{⃗}{\ensuremath{{}^{\to}}}
\newunicodechar{ⁿ}{\ifmmode^{n}\else\textsuperscript{\ensuremath{n}}\fi}
\newunicodechar{ˣ}{\ifmmode^{x}\else\textsuperscript{\ensuremath{x}}\fi}
\newunicodechar{ᵇ}{\ifmmode^{b}\else\textsuperscript{\ensuremath{b}}\fi}
\newunicodechar{ʳ}{\ifmmode^{r}\else\textsuperscript{\ensuremath{r}}\fi}
\newunicodechar{ᵘ}{\ifmmode^{u}\else\textsuperscript{\ensuremath{u}}\fi}
\newunicodechar{ʸ}{\ifmmode^{y}\else\textsuperscript{\ensuremath{y}}\fi}
\newunicodechar{ᵃ}{\ifmmode^{a}\else\textsuperscript{\ensuremath{a}}\fi}
\newunicodechar{ᵉ}{\ifmmode^{e}\else\textsuperscript{\ensuremath{e}}\fi}
\newunicodechar{ᵏ}{\ifmmode^{k}\else\textsuperscript{\ensuremath{k}}\fi}
\newunicodechar{ᵖ}{\ifmmode^{p}\else\textsuperscript{\ensuremath{p}}\fi}
\newunicodechar{ᴺ}{\ifmmode^{N}\else\textsuperscript{\ensuremath{N}}\fi}
\newunicodechar{ᶜ}{\ifmmode^{c}\else\textsuperscript{\ensuremath{c}}\fi}
\newunicodechar{ʰ}{\ifmmode^{h}\else\textsuperscript{\ensuremath{h}}\fi}
\newunicodechar{ᵠ}{\ifmmode^{q}\else\textsuperscript{\ensuremath{q}}\fi}
\newunicodechar{ₙ}{\ifmmode_{n}\else\textsubscript{\ensuremath{n}}\fi}
\newunicodechar{ₐ}{\ifmmode_{a}\else\textsubscript{\ensuremath{a}}\fi}
\newunicodechar{ᵢ}{\ifmmode_{i}\else\textsubscript{\ensuremath{i}}\fi}
\newunicodechar{ᵣ}{\ifmmode_{r}\else\textsubscript{\ensuremath{r}}\fi}
\newunicodechar{ₖ}{\ifmmode_{k}\else\textsubscript{\ensuremath{k}}\fi}
\newunicodechar{ᵦ}{\ifmmode_{\beta}\else\textsubscript{\ensuremath{\beta}}\fi}
\newunicodechar{ₓ}{\ifmmode_{x}\else\textsubscript{\ensuremath{x}}\fi}
\newfontfamily\popdisplay{ArchivoBlack-Regular.ttf}[Path=../../../fonts/]
\newfontfamily\poplogo{SpaceGrotesk-Bold.ttf}[Path=../../../fonts/]
\newfontfamily\popsemi{PlusJakartaSans-SemiBold.ttf}[Path=../../../fonts/]
\newfontfamily\popxbold{PlusJakartaSans-ExtraBold.ttf}[Path=../../../fonts/]
\newfontfamily\popxboldls{PlusJakartaSans-ExtraBold.ttf}[Path=../../../fonts/,LetterSpace=3.0]

\setlist[enumerate]{leftmargin=1.7em,itemsep=5pt,topsep=4pt}
\setlist[itemize]{leftmargin=1.7em,itemsep=4pt,topsep=4pt,label={\color{poporange}\textbullet}}
\setlength{\parindent}{0pt}
\setlength{\parskip}{5pt}
\renewcommand{\arraystretch}{1.25}

% pgfplots : style Pop par défaut
\pgfplotsset{
  every axis/.append style={axis line style={popink,line width=1pt},tick style={popink},
    grid style={popline,line width=0.5pt},label style={font=\small},tick label style={font=\footnotesize}},
  popcurve/.style={poporange,line width=1.8pt,no markers,smooth},
}
% Même style disponible dans un \draw TikZ simple (hors pgfplots)
\tikzset{popcurve/.style={poporange,line width=1.8pt}}
\tikzset{popgrid/.style={popline,very thin}}
\colorlet{popcurve}{poporange} % le nom du style est parfois utilisé comme couleur

% ============================================================
% Pill / squiggle
% ============================================================
\newcommand{\poppill}[3]{%
  \tikz[baseline=-0.55ex]\node[draw=popink,line width=1.2pt,rounded corners=2.6pt,%
    fill=#2,inner xsep=6.5pt,inner ysep=3pt]{%
    \popxboldls\fontsize{7}{7}\selectfont\color{#3}\MakeUppercase{#1}};%
}
\newcommand{\popsquiggle}[1]{%
  \tikz[baseline=-0.4ex]{
    \draw[#1,line width=2.6pt,line cap=round]
      plot [smooth] coordinates {(0,0.09) (0.34,0.01) (0.68,0.21) (1.02,0.01) (1.36,0.10)};
  }%
}
% Mot-clé surligné (marqueur orange sous le texte)
\newcommand{\cle}[1]{{\popxbold\color{popdark}#1}}

% ============================================================
% En-tête / pied
% ============================================================
\pagestyle{fancy}
\fancyhf{}
\renewcommand{\headrulewidth}{0pt}
\renewcommand{\footrulewidth}{0pt}
\lhead{\raisebox{-0.15ex}{\poplogo\fontsize{13}{13}\selectfont\color{popink}OnBuch\color{poporange}+}}
\rhead{\poppill{\DOCMATIERE\ \textbullet\ \DOCNIVEAU\ \textbullet\ Leçon \DOCLECON}{white}{popink}}
\lfoot{\popxbold\fontsize{7.6}{7.6}\selectfont\color{popmuted}\MakeUppercase{Cours signé OnBuch+}\ \textbullet\ {\popsemi\DOCTITRE}}
\rfoot{\tikz[baseline=-0.6ex]\node[rounded corners=8.5pt,fill=popink,minimum height=17pt,minimum width=17pt,inner xsep=5pt,inner sep=0]{\popxbold\fontsize{8}{8}\selectfont\color{popcream}\thepage\,/\,\pageref*{LastPage}};}

% ============================================================
% Titres
% ============================================================
\newcommand{\chaptitre}[2]{%
  \begin{tcolorbox}[enhanced,colback=poporange,colframe=popink,boxrule=2.6pt,arc=11pt,
      drop shadow=popink,left=15pt,right=15pt,top=12pt,bottom=11pt,before skip=3pt,after skip=18pt]
    \poppill{#2}{popink}{popcream}\par\vspace{9pt}
    {\raggedright\hyphenpenalty=10000\exhyphenpenalty=10000\popdisplay\fontsize{22}{24}\selectfont\color{popcream}\MakeUppercase{#1}\par}
  \end{tcolorbox}
}
% Section numérotée : pastille orange avec le numéro + titre Archivo
\newcounter{popsec}
\newcommand{\coursec}[1]{%
  \stepcounter{popsec}\setcounter{popsub}{0}%
  \par\vspace{16pt}\noindent
  \begin{minipage}{\linewidth}
  \tikz[baseline=-0.7ex]\node[draw=popink,line width=1.6pt,fill=poporange,rounded corners=4pt,minimum size=22pt,inner sep=0]{\popdisplay\fontsize{12}{12}\selectfont\color{popcream}\thepopsec};%
  \hspace{8pt}{\popdisplay\fontsize{15}{17}\selectfont\color{popink}\MakeUppercase{#1}}\par
  \vspace{4pt}\noindent{\color{poporange}\rule{36pt}{3.6pt}}
  \end{minipage}\par\nopagebreak\vspace{8pt}
}
\newcounter{popsub}
\newcommand{\courssub}[1]{%
  \stepcounter{popsub}%
  \par\vspace{9pt}\noindent{\popxbold\fontsize{11.5}{13}\selectfont\color{popink}{\color{poporange}\thepopsec.\thepopsub}\hspace{6pt}#1}\par\nopagebreak\vspace{4pt}
}

% ============================================================
% Boîtes pédagogiques — même recette : fond blanc, bordure épaisse colorée,
% ombre dure encre, badge plein. Titre optionnel [#1].
% ============================================================
\tcbset{popbase/.style={breakable,enhanced,colback=white,boxrule=2.3pt,arc=9pt,
  shadow={3pt}{-3pt}{0pt}{popink},left=14pt,right=14pt,top=11pt,bottom=11pt,before skip=11pt,after skip=15pt,
  fontupper=\popsemi\fontsize{10.6}{14.5}\selectfont\color{popink},
  fontlower=\popsemi\fontsize{10.6}{14.5}\selectfont\color{popink}}}
% Coche dessinée (le glyphe ✓ n'existe pas dans les polices du document)
\newcommand{\popcheck}[1]{\tikz[baseline=-0.5ex]\draw[#1,line width=1.6pt,line cap=round,line join=round](-0.13,0.0)--(-0.04,-0.09)--(0.13,0.1);}
\providecommand{\coche}{\popcheck{popgreen}}
\providecommand{\resultat}[1]{\par\smallskip\noindent\fcolorbox{popgreen}{popgreenL}{\parbox{\dimexpr\linewidth-2\fboxsep-2\fboxrule\relax}{\textbf{Résultat.} #1}}\par\smallskip}
\newcommand{\popboxhead}[3]{\poppill{#1}{#2}{white}\ifx\relax#3\relax\else\hspace{6pt}{\popxbold\fontsize{10.6}{12}\selectfont\color{popink}#3}\fi\par\vspace{7pt}}

\newtcolorbox{aretenir}[1][]{popbase,colframe=popgreen,before upper={\popboxhead{À retenir}{popgreen}{#1}}}
% Texte en chinois (document de lecture, dialogue, extrait) : cadre
% d'étude. Un titre optionnel (source, auteur) s'affiche dans l'en-tête.
\newtcolorbox{textebox}[1][]{popbase,colframe=popblue,before upper={\popboxhead{文本 · Texte}{popblue}{#1}},after upper={}}
% Marques ✓ / ✗ (forme correcte / fautive) souvent utilisées par les modèles
\providecommand{\cross}{\textcolor{poppink}{\ensuremath{\times}}}
\providecommand{\xmark}{\cross}\providecommand{\crossmark}{\cross}\providecommand{\cmark}{\ensuremath{\checkmark}}
% Ligne de réponse / justification à compléter (inventée par les modèles)
\providecommand{\justification}[1]{\par\noindent\textit{Justification~:} \dotfill\par}
% \textipa (package tipa non chargé) : la police ne couvre pas l'API -> simple passage
\providecommand{\textipa}[1]{#1}
% Soulignement (ulem non chargé) : \uline, \ul
\providecommand{\uline}[1]{\underline{#1}}\providecommand{\ul}[1]{\underline{#1}}
% \glossaire{\item ...} inventée par les modèles : liste de glossaire
\providecommand{\glossaire}[1]{\begin{itemize}#1\end{itemize}}
% \alert (beamer) inventée par les modèles : texte mis en évidence
\providecommand{\alert}[1]{\textbf{\textcolor{poppink}{#1}}}
% variantes ASCII d'ordinaux écrites en macro
\providecommand{\iere}{\textsuperscript{re}}\providecommand{\ieme}{\textsuperscript{e}}\providecommand{\ere}{\textsuperscript{re}}\providecommand{\eme}{\textsuperscript{e}}\providecommand{\er}{\textsuperscript{er}}\providecommand{\ie}{\textsuperscript{e}}
% Ordinaux français écrits en macro par les modèles : 1\ière, 3\ième
\newcommand{\ière}{\textsuperscript{re}}\newcommand{\ième}{\textsuperscript{e}}
% \exemple (inventée par les modèles) : étiquette « Exemple : »
\providecommand{\exemple}{\textbf{Exemple~:}\ }
% \square utilisable aussi hors mode math (cases à cocher)
\AtBeginDocument{\let\squareMath\square\renewcommand{\square}{\ensuremath{\squareMath}}}
% Mot ou phrase en chinois à l'intérieur d'un paragraphe français
\newcommand{\zh}[1]{\ifmmode\text{#1}\else{#1}\fi}
\let\eng\zh \let\esp\zh \let\all\zh % alias tolérés
% flèches d'intonation inventées par les modèles
\providecommand{\textsearrow}{}\renewcommand{\textsearrow}{\ensuremath{\searrow}}\providecommand{\textnearrow}{}\renewcommand{\textnearrow}{\ensuremath{\nearrow}}
% \fr{..} : mot français (inventé par les modèles)
\providecommand{\fr}[1]{\textit{#1}}
% zhbox : encadré de texte chinois inventé par les modèles
\newenvironment{zhbox}{\begin{quote}}{\end{quote}}
% \vs : « versus » inventé par les modèles
\providecommand{\vs}{\textit{vs}\ }
% \blank : case à compléter inventée par les modèles
\providecommand{\blank}[1][]{\underline{\hspace{1.6em}}}
\newtcolorbox{exemplebox}[1][]{popbase,colframe=poporange,before upper={\popboxhead{Exemple}{poporange}{#1}}}
\newtcolorbox{definition}[1][]{popbase,colframe=popblue,before upper={\popboxhead{Définition}{popblue}{#1}}}
\newtcolorbox{propriete}[1][]{popbase,colframe=poppurple,before upper={\popboxhead{Propriété}{poppurple}{#1}}}
\newtcolorbox{methode}[1][]{popbase,colframe=popgold,before upper={\popboxhead{Méthode}{popgold}{#1}}}
\newtcolorbox{attention}[1][]{popbase,colframe=poppink,before upper={\popboxhead{Attention}{poppink}{#1}}}
\newtcolorbox{experience}[1][]{popbase,colframe=popblue,colback=popblueL,before upper={\popboxhead{Expérience}{popblue}{#1}}}
% « Le savais-tu ? » : carte violette pleine (contexte, culture, Cameroun)
\newtcolorbox{savaistu}[1][]{popbase,colback=poppurple,colframe=popink,
  fontupper=\popsemi\fontsize{10.4}{14.2}\selectfont\color{white},
  before upper={\poppill{Le savais-tu\,?}{white}{poppurple}\ifx\relax#1\relax\else\hspace{6pt}{\popxbold\color{white}#1}\fi\par\vspace{7pt}}}
% Exercice résolu : énoncé en haut, \tcblower puis la solution sur fond crème
\newcounter{popexo}\newcounter{popexr}
\newtcolorbox{exoresolu}[1][]{popbase,colframe=poporange,colbacklower=popcream,
  segmentation style={popink,line width=1.2pt,dashed},
  before upper={\stepcounter{popexr}\popboxhead{Exercice résolu \thepopexr}{poporange}{#1}},
  before lower={\poppill{Solution}{popink}{popcream}\par\vspace{6pt}}}
% Exercice (non résolu en place) — la correction est dans la partie Corrigés
\newtcolorbox{exercice}[1][]{popbase,colframe=popink,
  before upper={\stepcounter{popexo}\popboxhead{Exercice \thepopexo}{popink}{#1}}}
% Corrigé
\newtcolorbox{corrige}[1][]{popbase,colframe=popgreen,colback=popgreenL,
  before upper={\popboxhead{Corrigé}{popgreen}{#1}}}
% Objectifs / prérequis / bilan
\newtcolorbox{objectifs}{popbase,colframe=popink,colback=poporangeL,
  before upper={\popboxhead{Objectifs de la leçon}{popink}{}}}
\newtcolorbox{prerequis}{popbase,colframe=popink,colback=popblueL,
  before upper={\popboxhead{Prérequis}{popblue}{}}}
\newtcolorbox{bilan}{popbase,colframe=popink,colback=white,boxrule=2.8pt,
  before upper={\popboxhead{Fiche bilan}{poporange}{L'essentiel à connaître pour l'examen}}}
% Figure encadrée avec légende
\newtcolorbox{popfigure}[1][]{enhanced,breakable=false,colback=white,colframe=popline,boxrule=1.4pt,arc=8pt,
  left=10pt,right=10pt,top=10pt,bottom=8pt,before skip=10pt,after skip=12pt,halign=center,
  lower separated=false,halign lower=center,
  fontlower=\popsemi\fontsize{9}{11.5}\selectfont\color{popmuted},
  #1}
\newcommand{\legende}[1]{\par\vspace{6pt}{\popsemi\fontsize{9}{11.5}\selectfont\color{popmuted}#1}}

% ============================================================
% Signature OnBuch+
% ============================================================
\newcommand{\poplogoinline}[1]{{\poplogo\fontsize{#1}{#1}\selectfont\color{popink}OnBuch\color{poporange}+}}
\newcommand{\signatureonbuch}{%
  \begin{tcolorbox}[enhanced,colback=popink,colframe=popink,boxrule=2.2pt,arc=10pt,
      drop shadow=poporange,left=14pt,right=14pt,top=10pt,bottom=10pt,before skip=0pt,after skip=0pt]
    \tikz[baseline=-0.7ex]\node[circle,fill=popgreen,draw=popcream,line width=1.3pt,minimum size=22pt,inner sep=0]{\popcheck{white}};%
    \hspace{9pt}\begin{minipage}[c]{0.62\linewidth}
      {\popxbold\fontsize{12}{13}\selectfont\color{popcream}Signé\ \textbullet\ {\poplogo OnBuch\color{poporange}+}}\par\vspace{2pt}
      {\raggedright\popsemi\fontsize{8}{10.5}\selectfont\color{popline}\MakeUppercase{Équipe pédagogique OnBuch+}\\\MakeUppercase{Conforme au programme officiel MINESEC}\par}
    \end{minipage}\hfill
    {\poplogo\fontsize{22}{22}\selectfont\color{popcream}OnBuch\color{poporange}+}
  \end{tcolorbox}%
}

% ============================================================
% Page de garde
% #1 titre · #2 matière · #3 niveau · #4 module · #5 n° leçon · #6 durée ·
% #7 description / objectifs (texte ou itemize)
% ============================================================
\newcommand{\pagegarde}[7]{%
  \thispagestyle{empty}
  \newgeometry{a4paper,margin=1.7cm,top=1.6cm,bottom=1.4cm}%
  \begin{tikzpicture}[remember picture,overlay]
    \fill[poporange!18] ([xshift=-3.2cm,yshift=-3.6cm]current page.north east) circle (4.6cm);
    \fill[poppurple!14] ([xshift=2.4cm,yshift=7.6cm]current page.south west) circle (3.8cm);
  \end{tikzpicture}%
  {\poplogo\fontsize{30}{30}\selectfont\color{popink}OnBuch\color{poporange}+}\hfill
  \raisebox{0.6ex}{\poppill{Cours signé \textbullet\ Premium}{popink}{popcream}}\par
  \vspace{1pt}\popsquiggle{poppurple}\par
  \vspace{8pt}
  \begin{tcolorbox}[enhanced,colback=poporange,colframe=popink,boxrule=2.8pt,arc=13pt,
      drop shadow=popink,left=18pt,right=18pt,top=12pt,bottom=13pt,after skip=10pt]
    \poppill{#2 \textbullet\ #3}{popink}{popcream}\hspace{5pt}\poppill{#4}{white}{popink}\par\vspace{8pt}
    {\popdisplay\fontsize{14}{14}\selectfont\color{popink}LEÇON #5}\par\vspace{4pt}
    {\raggedright\hyphenpenalty=10000\exhyphenpenalty=10000\popdisplay\fontsize{25}{27}\selectfont\color{popcream}\MakeUppercase{#1}\par}
    \vspace{8pt}
    {\popxbold\fontsize{9.5}{9.5}\selectfont\color{popink}\MakeUppercase{Durée conseillée : #6}}
  \end{tcolorbox}
  \begin{tcolorbox}[enhanced,colback=white,colframe=popink,boxrule=2.2pt,arc=10pt,
      drop shadow=popink,left=16pt,right=16pt,top=10pt,bottom=10pt,after skip=8pt]
    \poppill{Dans cette leçon}{poppurple}{white}\par\vspace{5pt}
    \popsemi\fontsize{9.3}{12.6}\selectfont\color{popink}#7
  \end{tcolorbox}
  \noindent\begin{minipage}[t]{0.487\linewidth}
    \begin{tcolorbox}[enhanced,colback=popgreen,colframe=popink,boxrule=2.2pt,arc=10pt,
        drop shadow=popink,left=11pt,right=11pt,top=10pt,bottom=10pt,height=3.1cm,valign=center]
      \tcbox[colback=white,colframe=popink,boxrule=1.3pt,arc=5pt,boxsep=0pt,nobeforeafter,
        left=3pt,right=3pt,top=3pt,bottom=3pt,valign=center]{\qrcode[height=1.9cm]{https://play.google.com/store/apps/details?id=com.luvvix.onbuch&hl=fr}}%
      \hspace{8pt}\begin{minipage}[c]{0.5\linewidth}
        {\popdisplay\fontsize{11}{12.5}\selectfont\color{white}\MakeUppercase{Télécharge OnBuch}}\par\vspace{4pt}
        {\raggedright\popsemi\fontsize{8.4}{11}\selectfont\color{white}Gratuit \textbullet\ Google Play\\Tuteur IA, annales, concours, résultats\par}
      \end{minipage}
    \end{tcolorbox}
  \end{minipage}\hfill
  \begin{minipage}[t]{0.487\linewidth}
    \begin{tcolorbox}[enhanced,colback=poppurple,colframe=popink,boxrule=2.2pt,arc=10pt,
        drop shadow=popink,left=11pt,right=11pt,top=10pt,bottom=10pt,height=3.1cm,valign=center]
      \tikz[baseline=-0.7ex]\node[circle,fill=white,draw=popink,line width=1.3pt,minimum size=1.6cm,inner sep=0]{\popdisplay\fontsize{24}{24}\selectfont\color{poppurple}+};%
      \hspace{8pt}\begin{minipage}[c]{0.6\linewidth}
        {\popdisplay\fontsize{11}{12.5}\selectfont\color{white}\MakeUppercase{Passe à OnBuch+}}\par\vspace{4pt}
        {\raggedright\popsemi\fontsize{8.4}{11}\selectfont\color{white}Tous les cours signés, Tuteur IA illimité, fiches PDF — onglet OnBuch+ de l'app\par}
      \end{minipage}
    \end{tcolorbox}
  \end{minipage}
  \vspace{8pt}
  \popcontenu
  \vfill
  \signatureonbuch
  \restoregeometry
  \newpage
}

% Rangée « Ce cours contient » (page de garde)
\newcommand{\poptile}[3]{% #1 couleur #2 gros texte #3 légende
  \begin{tcolorbox}[enhanced,colback=#1,colframe=popink,boxrule=2pt,arc=9pt,shadow={2.5pt}{-2.5pt}{0pt}{popink},
      left=6pt,right=6pt,top=9pt,bottom=9pt,halign=center,valign=center,height=1.75cm,width=\linewidth,nobeforeafter]
    {\popdisplay\fontsize{12.5}{13}\selectfont\color{white}\MakeUppercase{#2}}\par\vspace{3pt}
    {\centering\popsemi\fontsize{7.8}{9.5}\selectfont\color{white}#3\par}
  \end{tcolorbox}}
\newcommand{\popcontenu}{%
  \par\noindent{\popxboldls\fontsize{7.5}{7.5}\selectfont\color{popmuted}CE COURS SIGNÉ CONTIENT}\par\vspace{7pt}
  \noindent\begin{minipage}[t]{0.235\linewidth}\poptile{popblue}{Cours}{complet, illustré, pas à pas}\end{minipage}\hfill
  \begin{minipage}[t]{0.235\linewidth}\poptile{poporange}{Méthodes}{et exercices résolus}\end{minipage}\hfill
  \begin{minipage}[t]{0.235\linewidth}\poptile{poppink}{Exercices}{type examen + corrigés}\end{minipage}\hfill
  \begin{minipage}[t]{0.235\linewidth}\poptile{popgreen}{Fiche bilan}{l'essentiel à retenir}\end{minipage}\par}

% Sommaire
\newtcolorbox{sommaire}{enhanced,colback=white,colframe=popink,boxrule=2.2pt,arc=10pt,
  drop shadow=popink,left=18pt,right=18pt,top=13pt,bottom=13pt,before skip=6pt,after skip=20pt,
  before upper={\poppill{Sommaire}{poppurple}{white}\par\vspace{9pt}\popsemi\fontsize{10.6}{14.5}\selectfont\color{popink}}}

% Signature de fin de cours — poussée tout en bas de la dernière page
\newcommand{\findecours}{%
  \par\vfill\nopagebreak
  \begin{center}\popsquiggle{poporange}\end{center}\vspace{4pt}
  \signatureonbuch
}
\AtBeginDocument{\def\llbracket{\mathopen{[\mkern-3.5mu[}}\def\rrbracket{\mathclose{]\mkern-3.5mu]}}} % intervalles d'entiers (glyphe ⟦ absent de la police)
% Glyphes absents de Fira Math : redéfinis à partir de glyphes disponibles
\AtBeginDocument{%
  \def\setminus{\mathbin{\backslash}}\let\smallsetminus\setminus
  \def\vdots{\mathord{\vbox{\baselineskip3.2pt\lineskiplimit0pt\kern4pt\hbox{.}\hbox{.}\hbox{.}}}}%
  \def\ddots{\mathinner{\mkern1mu\raise6.5pt\hbox{.}\mkern2mu\raise3.6pt\hbox{.}\mkern2mu\raise0.7pt\hbox{.}\mkern1mu}}%
  \def\triangle{\mathord{\scalebox{0.9}{$\symup{\Delta}$}}}%
  \def\nabla{\mathord{\raisebox{\depth}{\scalebox{1}[-1]{$\symup{\Delta}$}}}}%
  \def\checkmark{\popcheck{popdark}}%
  \def\star{\mathbin{\ast}}\def\bigstar{\mathord{\ast}}%
  \def\diamond{\mathbin{\scalebox{0.6}{$\Diamond$}}}%
  \def\bigcup{\mathop{\mathchoice{\raisebox{-0.3em}{\scalebox{1.7}{$\cup$}}}{\scalebox{1.25}{$\cup$}}{\cup}{\cup}}\displaylimits}%
  \def\bigcap{\mathop{\mathchoice{\raisebox{-0.3em}{\scalebox{1.7}{$\cap$}}}{\scalebox{1.25}{$\cap$}}{\cap}{\cap}}\displaylimits}%
  \def\lesseqqgtr{\mathrel{\lessgtr}}\def\bigoplus{\mathop{\oplus}\displaylimits}%
}
\providecommand{\ssi}{si et seulement si} % abréviation fréquente
\AtBeginDocument{\providecommand{\wideparen}{\overparen}} % arc de cercle

%% FIGFIX-BEGIN v4 — correctifs globaux des figures (docs/onbuch-generation/tools/figfix)
\usepackage{adjustbox}\usepackage{etoolbox}
% toute figure TikZ/pgfplots plus large que la ligne est réduite à la largeur disponible
\BeforeBeginEnvironment{tikzpicture}{\begin{adjustbox}{max width=\linewidth}}
\AfterEndEnvironment{tikzpicture}{\end{adjustbox}}
% légendes pgfplots lisibles par-dessus les courbes
\pgfplotsset{every axis/.append style={legend style={fill=white,fill opacity=0.92,text opacity=1,draw=black!15,inner sep=3pt}}}
% étiquettes d'arêtes (syntaxe edge["..."]) avec fond blanc
\tikzset{every edge quotes/.append style={fill=white,inner sep=1.2pt}}
%% FIGFIX-END
\begin{document}

\pagegarde{Leçon 8 — Expression écrite : la délinquance juvénile}{Chinois}{1ère A}{Module 1 : Vie familiale et intégration sociale}{ZHA08}{2 h}{Cette leçon d'expression écrite entraîne les élèves de 1ère A à rédiger et traduire un texte simple en chinois sur la délinquance juvénile. Elle combine l'étude d'un modèle de texte commenté, l'écriture correcte des caractères du thème et des exercices de thème et de version.
\vspace{4pt}
\begin{multicols}{2}\small
\begin{itemize}
\item À la fin de cette leçon, je sais lire et comprendre un texte chinois simple sur la délinquance juvénile.
\item Je sais écrire correctement les caractères clés du thème (青少年, 犯罪, 法律, 教育...) en respectant radicaux et ordre des traits.
\item Je sais distinguer les caractères proches faciles à confondre (犯/范, 青/清, 法/去).
\item Je sais utiliser les connecteurs 因为...所以..., 首先/其次/最后, 应该 pour structurer un texte.
\item Je sais rédiger un texte simple de 5 à 8 phrases sur les causes et les solutions de la délinquance juvénile.
\item Je sais traduire des phrases simples du français vers le chinois (thème).
\end{itemize}\end{multicols}}

\begin{sommaire}
\begin{multicols}{2}
\begin{enumerate}
\item Modèle de texte et compréhension
\item Lexique du thème et collocations
\item Écriture correcte des caractères
\item Structure et connecteurs du texte argumentatif
\item Thème : traduction français → chinois
\item Version : traduction chinois → français
\item Production guidée et grille d'évaluation
\item Activité d'intégration
\item Méthodes et exercices résolus
\item Exercices
\item Corrigés des exercices
\item Fiche bilan
\end{enumerate}
\end{multicols}
\end{sommaire}

\begin{objectifs}
\begin{itemize}
\item À la fin de cette leçon, je sais lire et comprendre un texte chinois simple sur la délinquance juvénile.
\item Je sais écrire correctement les caractères clés du thème (青少年, 犯罪, 法律, 教育...) en respectant radicaux et ordre des traits.
\item Je sais distinguer les caractères proches faciles à confondre (犯/范, 青/清, 法/去).
\item Je sais utiliser les connecteurs 因为...所以..., 首先/其次/最后, 应该 pour structurer un texte.
\item Je sais rédiger un texte simple de 5 à 8 phrases sur les causes et les solutions de la délinquance juvénile.
\item Je sais traduire des phrases simples du français vers le chinois (thème).
\item Je sais traduire un court texte du chinois vers le français (version).
\item Je sais m'auto-évaluer à l'aide d'une grille d'évaluation type examen.
\end{itemize}
\end{objectifs}

\begin{prerequis}
\begin{itemize}
\item Structure de base de la phrase chinoise Sujet-Verbe-Objet (Seconde)
\item Connecteurs simples 和, 但是, 因为...所以... (Seconde)
\item Vocabulaire de la famille et de l'école (Module 1, leçons précédentes)
\item Notions de radicaux et d'ordre des traits (Seconde)
\item Le verbe modal 应该 (devoir, il faut)
\end{itemize}
\end{prerequis}

\newpage

\coursec{Modèle de texte et compréhension}

À Douala comme à Yaoundé, les journaux parlent souvent de jeunes impliqués dans les vols de téléphones ou les bagarres de quartier. Au Royaume-Uni, on débat du \emph{youth crime}, et en Chine, la loi sur la protection des mineurs encadre strictement la prévention de la délinquance juvénile. Comment parler de ce phénomène social en chinois ? Comment écrire un court texte argumenté sur les causes et les solutions ? C'est toute la problématique de cette leçon : vous allez apprendre à \cle{lire}, \cle{comprendre} puis \cle{imiter} un texte argumentatif chinois simple, avant de rédiger le vôtre. La première étape est toujours la même chez les bons sinisants : observer un \cle{texte modèle}, en dégager le plan et le vocabulaire.

\courssub{Lecture du texte modèle}

Lisez d'abord le texte en silence, puis à voix haute en suivant le pinyin. Ne cherchez pas encore à tout traduire : repérez simplement les mots que vous reconnaissez (\zh{问题}, \zh{应该}, \zh{因为}…).

\begin{textebox}[青少年犯罪是一个严重的问题 — La délinquance juvénile est un problème grave]
现在，青少年犯罪是一个严重的问题。\\
\emph{Xiànzài, qīngshàonián fànzuì shì yí ge yánzhòng de wèntí.}\\
« Aujourd'hui, la délinquance juvénile est un problème grave. »\\[0.4em]
因为有些青少年缺少家庭教育，所以他们走上了错误的道路。\\
\emph{Yīnwèi yǒuxiē qīngshàonián quēshǎo jiātíng jiàoyù, suǒyǐ tāmen zǒushàngle cuòwù de dàolù.}\\
« Parce que certains adolescents manquent d'éducation familiale, ils s'engagent sur une mauvaise voie. »\\[0.4em]
首先，家长应该多关心孩子；其次，学校应该加强法律教育；最后，社会也应该帮助这些年轻人。\\
\emph{Shǒuxiān, jiāzhǎng yīnggāi duō guānxīn háizi; qícì, xuéxiào yīnggāi jiāqiáng fǎlǜ jiàoyù; zuìhòu, shèhuì yě yīnggāi bāngzhù zhèxiē niánqīngrén.}\\
« D'abord, les parents devraient s'occuper davantage de leurs enfants ; ensuite, l'école devrait renforcer l'éducation juridique ; enfin, la société aussi devrait aider ces jeunes. »\\[0.4em]
我相信，只要大家一起努力，这个问题一定会解决。\\
\emph{Wǒ xiāngxìn, zhǐyào dàjiā yìqǐ nǔlì, zhège wèntí yídìng huì jiějué.}\\
« Je crois que, pourvu que tout le monde fasse des efforts ensemble, ce problème sera certainement résolu. »
\end{textebox}

\begin{definition}[Glossaire du texte modèle]
\begin{center}
\begin{tabularx}{\linewidth}{|X|X|X|X|}
\hline
\rowcolor{popblueL} Caractères & Pinyin & Nature & Traduction \\
\hline
青少年 & qīngshàonián & nom & adolescent, jeune \\
\hline
犯罪 & fànzuì & verbe / nom & commettre un crime ; délinquance \\
\hline
严重 & yánzhòng & adjectif & grave, sérieux \\
\hline
缺少 & quēshǎo & verbe & manquer de \\
\hline
家庭教育 & jiātíng jiàoyù & nom & éducation familiale \\
\hline
走上 & zǒushàng & verbe & s'engager sur (un chemin) \\
\hline
错误 & cuòwù & adjectif / nom & faux, erroné ; erreur \\
\hline
道路 & dàolù & nom & route, voie \\
\hline
家长 & jiāzhǎng & nom & parents (responsables de famille) \\
\hline
关心 & guānxīn & verbe & se soucier de, s'occuper de \\
\hline
加强 & jiāqiáng & verbe & renforcer \\
\hline
法律 & fǎlǜ & nom & loi, droit \\
\hline
社会 & shèhuì & nom & société \\
\hline
年轻人 & niánqīngrén & nom & jeunes, jeunes gens \\
\hline
相信 & xiāngxìn & verbe & croire, être convaincu \\
\hline
努力 & nǔlì & verbe / adjectif & faire des efforts ; appliqué \\
\hline
解决 & jiějué & verbe & résoudre \\
\hline
\end{tabularx}
\end{center}
\end{definition}

\begin{attention}[Prononciation : les pièges du texte]
\begin{itemize}
\item \zh{青少年} \emph{qīngshàonián} : attention à \textbf{qīng} (premier ton, son « ts-ing ») et non \emph{qǐng} ; \zh{少} se lit ici \textbf{shào} (quatrième ton, « jeune ») et non \emph{shǎo} (« peu »). Même caractère, deux lectures !
\item \zh{一个} : le caractère \zh{一} change de ton devant un quatrième ton : on lit \emph{yí ge} et non \emph{yī ge}.
\item \zh{一定} : même phénomène, on lit \emph{yídìng}.
\item \zh{错误} \emph{cuòwù} : deux quatrièmes tons descendants, bien marquer la chute sur chaque syllabe.
\end{itemize}
\end{attention}

\courssub{Comprendre le texte : questions de compréhension}

Avant toute traduction, vérifiez que vous avez compris le sens global. Répondez en français, puis essayez en chinois avec les mots du texte.

\begin{exoresolu}[Questions de compréhension du texte modèle]
Énoncé. Répondez aux questions suivantes : (1) Quel est le problème présenté dans le texte ? (2) Quelle est la cause mentionnée ? (3) Quelles sont les trois solutions proposées, et qui doit agir ? (4) Quelle est l'attitude de l'auteur à la fin du texte ?
\tcblower
Solution détaillée.
\begin{enumerate}
\item Le problème : la délinquance juvénile est un problème grave (\zh{青少年犯罪是一个严重的问题}).
\item La cause : certains adolescents manquent d'éducation familiale (\zh{缺少家庭教育}) ; c'est pourquoi ils s'engagent sur une mauvaise voie.
\item Les trois solutions, annoncées par \zh{首先} / \zh{其次} / \zh{最后} : les \textbf{parents} doivent s'occuper davantage de leurs enfants ; l'\textbf{école} doit renforcer l'éducation juridique ; la \textbf{société} doit aider ces jeunes. Trois acteurs donc : \zh{家长}, \zh{学校}, \zh{社会}.
\item L'auteur est \textbf{optimiste} : il croit (\zh{我相信}) que si tout le monde fait des efforts ensemble, le problème sera résolu.
\end{enumerate}
\end{exoresolu}

\courssub{La structure du texte argumentatif}

Ce petit texte de quatre phrases suit un plan en quatre étapes, typique du texte argumentatif chinois (et très proche du plan « thèse — cause — solutions — conclusion » que vous connaissez en français) :

\begin{enumerate}
\item \cle{Problème} (\zh{问题}) : phrase 1, on pose le sujet avec \zh{是} + jugement (\zh{是一个严重的问题}).
\item \cle{Cause} (\zh{原因}) : phrase 2, introduite par le couple \zh{因为……所以……}.
\item \cle{Solutions} (\zh{办法}) : phrase 3, énumérées par \zh{首先……其次……最后……}, chacune avec \zh{应该} (« devrait »).
\item \cle{Conclusion optimiste} (\zh{结论}) : phrase 4, avec \zh{我相信} et la structure \zh{只要……就/一定……} (« pourvu que… alors… »).
\end{enumerate}

\begin{propriete}[Le plan-type du texte argumentatif simple]
Schéma à retenir : \textbf{Problème → Cause → Solutions → Conclusion}. En chinois, ce plan est \emph{signalé par des connecteurs obligatoires} : \zh{因为……所以……} pour la cause, \zh{首先 / 其次 / 最后} pour l'énumération, \zh{我相信……} pour la conclusion. Contrairement au français où « parce que » et « donc » ne se combinent pas dans la même phrase, le chinois \emph{exige} souvent les deux : \zh{因为……所以……} est la norme, pas une faute de style.
\end{propriete}

\begin{popfigure}
\begin{tikzpicture}[node distance=0.55cm,
  bloc/.style={rectangle, rounded corners, draw=popblue, fill=popblueL, align=center, minimum width=3.1cm, minimum height=1cm, font=\small},
  fleche/.style={-{Stealth[length=3mm]}, thick, popdark}]
\node[bloc, align=center] (p){\textbf{Problème}\\问题 wèntí};
\node[bloc, below=of p, align=center] (c){\textbf{Cause}\\原因 yuányīn\\因为…所以…};
\node[bloc, below=of c, align=center] (s){\textbf{Solutions}\\办法 bànfǎ\\首先 / 其次 / 最后};
\node[bloc, below=of s, xshift=-3.4cm, fill=poporangeL, draw=poporange, align=center] (s1){家长\\jiāzhǎng};
\node[bloc, below=of s, fill=poporangeL, draw=poporange, align=center] (s2){学校\\xuéxiào};
\node[bloc, below=of s, xshift=3.4cm, fill=poporangeL, draw=poporange, align=center] (s3){社会\\shèhuì};
\node[bloc, below=2.1cm of s, fill=popgreenL, draw=popgreen, align=center] (concl){\textbf{Conclusion}\\结论 jiélùn\\我相信…只要…一定…};
\draw[fleche] (p) -- (c);
\draw[fleche] (c) -- (s);
\draw[fleche, poporange] (s.south) -- (s1.north);
\draw[fleche, poporange] (s.south) -- (s2.north);
\draw[fleche, poporange] (s.south) -- (s3.north);
\draw[fleche, popgreen] (s2.south) -- (concl.north);
\end{tikzpicture}
\legende{Organigramme du texte modèle : problème → cause → trois solutions (trois acteurs) → conclusion optimiste.}
\end{popfigure}

\begin{methode}[Lire un texte argumentatif chinois en quatre pas]
\begin{enumerate}
\item \textbf{Première lecture rapide} : ne traduisez rien, identifiez seulement le thème (ici : \zh{青少年犯罪}).
\item \textbf{Repérez les connecteurs} : soulignez \zh{因为}, \zh{所以}, \zh{首先}, \zh{其次}, \zh{最后}, \zh{我相信}. Ils donnent à eux seuls le plan du texte : cause, énumération, conclusion.
\item \textbf{Associez chaque connecteur à une étape du plan} : problème → cause → solutions → conclusion. Rédigez le plan en français en une ligne par étape.
\item \textbf{Lecture fine} : traduisez phrase par phrase en respectant l'ordre chinois (sujet + verbe + complément), puis vérifiez que votre traduction française suit le plan repéré à l'étape 3.
\end{enumerate}
\end{methode}

\courssub{Observer les phrases clés}

Deux phrases du texte méritent une observation attentive, car elles contiennent des structures que vous réutiliserez dans vos propres productions.

\begin{exemplebox}[Deux phrases clés, observées et traduites]
\zh{有些青少年缺少家庭教育。}\\
\emph{Yǒuxiē qīngshàonián quēshǎo jiātíng jiàoyù.}\\
« Certains adolescents manquent d'éducation familiale. »\\
Structure : \zh{有些} (« certains ») + nom + verbe \zh{缺少} (« manquer de ») + nom. Notez que \zh{有些} se place \emph{avant} le nom, comme « certains » en français.\\[0.5em]
\zh{学校应该加强法律教育。}\\
\emph{Xuéxiào yīnggāi jiāqiáng fǎlǜ jiàoyù.}\\
« L'école devrait renforcer l'éducation juridique. »\\
Structure : sujet + \zh{应该} (« devoir, il faudrait ») + verbe + complément. \zh{应该} est le verbe modal du conseil et du devoir moral : il se place toujours \emph{avant} le verbe principal, comme « devrait » en français.
\end{exemplebox}

\begin{attention}[Erreurs fréquentes des francophones]
\begin{itemize}
\item Ne dites pas \zh{犯罪青少年} pour « délinquance juvénile » : cela signifierait « des adolescents criminels ». L'expression correcte est \zh{青少年犯罪} (les jeunes + commettre des crimes), nom + verbe formant un nom composé.
\item \zh{因为} et \zh{所以} s'emploient \emph{ensemble} en chinois : \zh{因为……所以……}. En français, « parce que… donc… » serait incorrect ; en chinois, c'est la construction standard. À la traduction vers le français, n'en gardez qu'un seul.
\item \zh{关心} signifie « se soucier de » et prend directement un complément de personne : \zh{关心孩子}. N'ajoutez pas de préposition comme en français (« s'occuper \emph{de} »).
\end{itemize}
\end{attention}

\begin{popfigure}
\begin{tikzpicture}[mindmap, grow cyclic, every node/.style={concept, align=center, font=\small},
  root concept/.append style={concept color=popblue, text=white, font=\small\bfseries},
  level 1/.append style={level distance=3.1cm, sibling angle=90},
  level 2/.append style={level distance=2.3cm, sibling angle=60}]
\node[root concept, align=center]{青少年犯罪\\qīngshàonián fànzuì}
  child[concept color=poporange!70]{ node {人 Personnes} 
      child { node[align=center]{家长\\jiāzhǎng\\parents} }
      child { node[align=center]{年轻人\\niánqīngrén\\jeunes} } }
  child[concept color=popgreen!70]{ node {原因 Causes}
      child { node[align=center]{缺少\\quēshǎo\\manquer de} }
      child { node[align=center]{家庭教育\\jiātíng jiàoyù} } }
  child[concept color=poppurple!70]{ node {办法 Solutions}
      child { node[align=center]{加强\\jiāqiáng\\renforcer} }
      child { node[align=center]{帮助\\bāngzhù\\aider} } }
  child[concept color=popgold!80]{ node {结果 Résultat}
      child { node[align=center]{解决\\jiějué\\résoudre} }
      child { node[align=center]{努力\\nǔlì\\efforts} } };
\end{tikzpicture}
\legende{Carte mentale du vocabulaire autour de \zh{青少年犯罪} : personnes, causes, solutions, résultat.}
\end{popfigure}

\courssub{Comparer le français et le chinois}

\begin{center}
\begin{tabularx}{\linewidth}{|X|X|X|}
\hline
\rowcolor{popblueL} Notion & Français & Chinois \\
\hline
Cause & « parce que » \emph{ou} « donc » & \zh{因为……所以……} (les deux ensemble) \\
\hline
Énumération & d'abord, ensuite, enfin & \zh{首先、其次、最后} \\
\hline
Conseil & « devrait » + infinitif & \zh{应该} + verbe \\
\hline
Condition & « pourvu que… » & \zh{只要……就……} \\
\hline
Opinion & « je crois que… » & \zh{我相信……} \\
\hline
\end{tabularx}
\end{center}

\begin{savaistu}[La délinquance juvénile en Chine : un cadre légal strict]
En Chine, la prévention de la délinquance juvénile est encadrée par une loi spécifique : \zh{预防未成年人犯罪法} (\emph{Yùfáng wèichéngniánrén fànzuì fǎ}, « loi sur la prévention de la criminalité des mineurs »), révisée en 2020. L'âge de la responsabilité pénale y est en principe fixé à 16 ans, avec des dispositions particulières pour les affaires graves dès 14 ans. La loi insiste sur le rôle conjoint de la famille, de l'école et de la société — exactement les trois acteurs de notre texte modèle ! Au Cameroun aussi, la protection de l'enfance associe familles, écoles et institutions : un beau point de comparaison pour vos productions écrites.
\end{savaistu}

\begin{exercice}[À vous de jouer]
\begin{enumerate}
\item Retrouvez dans le texte modèle les cinq connecteurs et recopiez-les avec leur pinyin et leur traduction.
\item Répondez en chinois : \zh{青少年犯罪的原因是什么？} (Quelle est la cause de la délinquance juvénile ?)
\item Classez les mots suivants dans la carte mentale : \zh{法律}, \zh{社会}, \zh{错误}, \zh{相信}.
\end{enumerate}
\end{exercice}

\begin{corrige}[Exercice « À vous de jouer »]
\begin{enumerate}
\item \zh{因为} \emph{yīnwèi} « parce que » ; \zh{所以} \emph{suǒyǐ} « donc » ; \zh{首先} \emph{shǒuxiān} « d'abord » ; \zh{其次} \emph{qícì} « ensuite » ; \zh{最后} \emph{zuìhòu} « enfin » (on accepte aussi \zh{只要} \emph{zhǐyào} « pourvu que »).
\item Réponse attendue : \zh{因为有些青少年缺少家庭教育。} \emph{Yīnwèi yǒuxiē qīngshàonián quēshǎo jiātíng jiàoyù.} « Parce que certains adolescents manquent d'éducation familiale. »
\item \zh{法律} → Solutions (renforcer l'éducation juridique) ; \zh{社会} → Personnes/acteurs ; \zh{错误} → Causes (la mauvaise voie) ; \zh{相信} → Résultat/conclusion.
\end{enumerate}
\end{corrige}

\begin{aretenir}[L'essentiel de cette section]
Un texte argumentatif chinois simple se lit en repérant ses connecteurs : \zh{因为……所以……} (cause), \zh{首先 / 其次 / 最后} (solutions), \zh{我相信……只要……} (conclusion). Le plan est toujours : problème → cause → solutions → conclusion optimiste. Le vocabulaire clé du thème : \zh{青少年犯罪} (délinquance juvénile), \zh{严重的问题} (problème grave), \zh{缺少} (manquer de), \zh{加强} (renforcer), \zh{解决} (résoudre).
\end{aretenir}

\coursec{Lexique du thème et collocations}

Dans la section précédente, nous avons lu et compris le texte modèle sur la délinquance juvénile. Vous avez remarqué que certains mots revenaient souvent : \zh{青少年}, \zh{问题}, \zh{解决}, \zh{关心}… C'est normal : tout texte argumentatif repose sur un \cle{champ lexical} précis. Pour écrire vous-même un texte sur ce thème, vous devez d'abord posséder ce vocabulaire de manière active : connaître le caractère, le pinyin, la nature grammaticale, le sens exact, et surtout les \cle{collocations}, c'est-à-dire les combinaisons de mots qui vont naturellement ensemble en chinois. C'est l'objet de cette section.

\courssub{Observation : les mots du thème dans leur contexte}

Avant d'apprendre la liste, observons ces mots en action dans un court texte original.

\begin{textebox}[青少年犯罪：一个严重的问题]
现在，青少年犯罪是一个严重的问题。\\
\emph{Xiànzài, qīngshàonián fànzuì shì yí ge yánzhòng de wèntí.}\\
« Aujourd'hui, la délinquance juvénile est un problème grave. »\\[4pt]
为什么有些年轻人会走上错误的道路？很多专家认为，这些孩子缺少家庭教育和关心。\\
\emph{Wèishénme yǒuxiē niánqīngrén huì zǒushàng cuòwù de dàolù? Hěn duō zhuānjiā rènwéi, zhèxiē háizi quēshǎo jiātíng jiàoyù hé guānxīn.}\\
« Pourquoi certains jeunes s'engagent-ils sur une mauvaise voie ? Beaucoup de spécialistes pensent que ces enfants manquent d'éducation familiale et d'attention. »\\[4pt]
家长应该多关心孩子，学校应该加强法律教育，社会也应该帮助他们。\\
\emph{Jiāzhǎng yīnggāi duō guānxīn háizi, xuéxiào yīnggāi jiāqiáng fǎlǜ jiàoyù, shèhuì yě yīnggāi bāngzhù tāmen.}\\
« Les parents doivent se soucier davantage de leurs enfants, l'école doit renforcer l'éducation juridique, et la société doit aussi les aider. »\\[4pt]
我们应该相信：只要大家一起努力，就一定能解决这个问题。\\
\emph{Wǒmen yīnggāi xiāngxìn: zhǐyào dàjiā yìqǐ nǔlì, jiù yídìng néng jiějué zhège wèntí.}\\
« Nous devons croire que tant que tout le monde fait des efforts ensemble, nous pourrons certainement résoudre ce problème. »
\end{textebox}

\textbf{Questions de compréhension}
\begin{enumerate}
\item Selon le texte, pourquoi certains jeunes commettent-ils des crimes ?
\item Que doivent faire les parents, l'école et la société ?
\item Relevez dans le texte les quatre collocations étudiées dans cette leçon.
\end{enumerate}

\courssub{Glossaire du thème}

\begin{definition}[青少年 qīngshàonián : un mot composé]
\zh{青少年} est un mot composé de trois caractères : \zh{青} (qīng, « jeune, vert de jeunesse ») + \zh{少} (shào, « jeune, peu ») + \zh{年} (nián, « année, âge »). Littéralement : « les années vertes de la jeunesse », c'est-à-dire \cle{les adolescents, les jeunes}. C'est un nom collectif : dans un texte argumentatif, on l'utilise généralement pour désigner le groupe (\zh{青少年犯罪}), ou avec un quantificateur (\zh{一些青少年} « certains adolescents », \zh{这些年轻人} « ces jeunes »).
\end{definition}

\begin{center}
\begin{tabularx}{\linewidth}{|X|X|X|X|}
\hline
\rowcolor{popblueL} Caractères & Pinyin & Nature & Traduction \\
\hline
青少年 & qīngshàonián & n. & adolescent, jeune \\
\hline
犯罪 & fànzuì & v. / n. & commettre un crime / crime, délit \\
\hline
严重 & yánzhòng & adj. & grave, sérieux \\
\hline
问题 & wèntí & n. & problème ; question \\
\hline
缺少 & quēshǎo & v. & manquer de \\
\hline
家庭教育 & jiātíng jiàoyù & n. & éducation familiale \\
\hline
法律 & fǎlǜ & n. & loi, droit \\
\hline
社会 & shèhuì & n. & société \\
\hline
关心 & guānxīn & v. & se soucier de, s'intéresser à \\
\hline
加强 & jiāqiáng & v. & renforcer \\
\hline
帮助 & bāngzhù & v. / n. & aider / aide \\
\hline
解决 & jiějué & v. & résoudre \\
\hline
错误 & cuòwù & adj. / n. & erroné / erreur \\
\hline
道路 & dàolù & n. & chemin, voie, route \\
\hline
努力 & nǔlì & v. / adj. & faire des efforts / appliqué, assidu \\
\hline
相信 & xiāngxìn & v. & croire, avoir confiance en \\
\hline
\end{tabularx}
\end{center}

\begin{attention}[问题 : « problème » ou « question » ?]
\zh{问题} (wèntí) a deux sens que le français distingue : « problème » et « question ». Le contexte et le verbe associé décident :
\begin{itemize}
\item \zh{回答问题} (huídá wèntí) : répondre à une \textbf{question} (quelqu'un pose une question, on y répond).
\item \zh{解决问题} (jiějué wèntí) : résoudre un \textbf{problème} (une difficulté qu'on élimine).
\end{itemize}
Ne traduisez donc jamais \zh{问题} mécaniquement : regardez le verbe qui l'accompagne. À l'écrit, c'est une erreur fréquente des francophones de traduire \zh{我有一个问题} par « J'ai un problème » alors qu'il s'agit souvent de « J'ai une question (à poser) ».
\end{attention}

\begin{savaistu}[关心 et 帮助 se construisent sans préposition]
En français, on dit « se soucier \textbf{de} quelqu'un », « s'intéresser \textbf{à} quelqu'un », « aider \textbf{à} faire quelque chose ». En chinois, \zh{关心} (guānxīn) et \zh{帮助} (bāngzhù) se construisent \cle{directement avec leur objet}, sans aucune préposition :
\begin{itemize}
\item \zh{关心孩子} (guānxīn háizi) : se soucier des enfants.
\item \zh{帮助他们} (bāngzhù tāmen) : les aider.
\end{itemize}
N'écrivez jamais \zh{关心对孩子} ou \zh{帮助给他们} : ce sont des calques du français.
\end{savaistu}

\courssub{Les collocations : les mots qui vont ensemble}

Une \cle{collocation} est une combinaison habituelle de mots. En chinois comme en français, on ne combine pas les mots au hasard : on dit « commettre un crime » et non « faire un crime ». Retenez les quatre collocations essentielles du thème.

\begin{center}
\begin{tabularx}{\linewidth}{|X|X|X|}
\hline
\rowcolor{popblueL} Collocation & Pinyin & Traduction \\
\hline
解决问题 & jiějué wèntí & résoudre un problème \\
\hline
走上错误的道路 & zǒushàng cuòwù de dàolù & s'engager sur une mauvaise voie \\
\hline
关心孩子 & guānxīn háizi & se soucier des enfants \\
\hline
加强教育 & jiāqiáng jiàoyù & renforcer l'éducation \\
\hline
\end{tabularx}
\end{center}

\begin{exemplebox}[Chaque collocation dans une phrase]
\zh{我们应该解决这个问题。}\\
\emph{Wǒmen yīnggāi jiějué zhège wèntí.} « Nous devons résoudre ce problème. »\\[4pt]
\zh{缺少关心的孩子容易走上错误的道路。}\\
\emph{Quēshǎo guānxīn de háizi róngyì zǒushàng cuòwù de dàolù.} « Les enfants qui manquent d'attention s'engagent facilement sur une mauvaise voie. »\\[4pt]
\zh{家长应该多关心孩子。}\\
\emph{Jiāzhǎng yīnggāi duō guānxīn háizi.} « Les parents devraient se soucier davantage de leurs enfants. »\\[4pt]
\zh{学校应该加强法律教育。}\\
\emph{Xuéxiào yīnggāi jiāqiáng fǎlǜ jiàoyù.} « L'école devrait renforcer l'éducation juridique. »
\end{exemplebox}

Notez la structure de \zh{走上错误的道路} : le verbe \zh{走上} (zǒushàng, « monter en marchant sur ») est un \cle{verbe à complément de direction} ; \zh{上} (shàng) indique le mouvement vers le haut / l'engagement sur une voie. On retrouve ce \zh{上} directionnel dans \zh{上车} (shàngchē, monter dans un véhicule) ou \zh{上楼} (shànglóu, monter à l'étage).

\courssub{Champ lexical associé : les acteurs sociaux}

Pour compléter votre texte à l'examen, vous aurez besoin des noms des \cle{acteurs} qui interviennent face à la délinquance juvénile.

\begin{center}
\begin{tabularx}{\linewidth}{|X|X|X|X|}
\hline
\rowcolor{popblueL} Caractères & Pinyin & Nature & Traduction \\
\hline
家长 & jiāzhǎng & n. & parents (de famille), responsables d'enfants \\
\hline
学校 & xuéxiào & n. & école \\
\hline
警察 & jǐngchá & n. & police, policier \\
\hline
法院 & fǎyuàn & n. & tribunal \\
\hline
\end{tabularx}
\end{center}

\begin{exemplebox}[Les acteurs en contexte]
\zh{警察帮助这些年轻人。}\\
\emph{Jǐngchá bāngzhù zhèxiē niánqīngrén.} « La police aide ces jeunes. »\\[4pt]
\zh{法院按照法律处理青少年犯罪的问题。}\\
\emph{Fǎyuàn ànzhào fǎlǜ chǔlǐ qīngshàonián fànzuì de wèntí.} « Le tribunal traite les affaires de délinquance juvénile conformément à la loi. »\\[4pt]
\zh{家长和学校应该一起努力。}\\
\emph{Jiāzhǎng hé xuéxiào yīnggāi yìqǐ nǔlì.} « Les parents et l'école doivent faire des efforts ensemble. »
\end{exemplebox}

Remarquez que \zh{家长} (jiāzhǎng) désigne les parents \emph{en tant que responsables de l'enfant} : c'est le mot utilisé dans les contextes scolaires et sociaux (\zh{家长会} : réunion de parents d'élèves), alors que \zh{父母} (fùmǔ) désigne le père et la mère dans la sphère privée familiale.

\courssub{Carte mentale du champ lexical}

\begin{popfigure}
\begin{center}
\begin{tikzpicture}[mindmap, grow cyclic, every node/.style={concept, align=center}, concept color=popblueL, text=popdark,
level 1/.append style={level distance=3.2cm, sibling angle=90},
level 2/.append style={level distance=2.2cm, sibling angle=45}]
\node[concept color=poporangeL, align=center]{青少年犯罪\\ \emph{délinquance juvénile}}
  child[concept color=popgreenL] { node[align=center]{问题\\ problème}
    child { node[align=center]{严重\\ grave} }
    child { node[align=center]{解决\\ résoudre} }
  }
  child[concept color=poppurpleL] { node[align=center]{Causes\\ 原因}
    child { node[align=center]{缺少\\ manquer de} }
    child { node[align=center]{家庭教育\\ éducation familiale} }
  }
  child[concept color=poppinkL] { node[align=center]{Actions\\ 行动}
    child { node[align=center]{关心\\ se soucier de} }
    child { node[align=center]{加强\\ renforcer} }
    child { node[align=center]{帮助\\ aider} }
  }
  child[concept color=popgoldL] { node[align=center]{Acteurs\\ 人员}
    child { node[align=center]{家长\\ parents} }
    child { node[align=center]{警察\\ police} }
    child { node[align=center]{法院\\ tribunal} }
  };
\end{tikzpicture}
\end{center}
\legende{Carte mentale du champ lexical de la délinquance juvénile : le problème, ses causes, les actions, les acteurs.}
\end{popfigure}

\courssub{Comparaison chinois / français : pièges de traduction}

\begin{center}
\begin{tabularx}{\linewidth}{|X|X|X|}
\hline
\rowcolor{popblueL} Structure chinoise & Exemple en chinois + pinyin & Traduction et remarque \\
\hline
Verbe transitif direct & \zh{关心孩子} guānxīn háizi & « se soucier \textbf{des} enfants » : pas de préposition en chinois \\
\hline
Verbe transitif direct & \zh{缺少家庭教育} quēshǎo jiātíng jiàoyù & « manquer \textbf{d'}éducation familiale » : pas de « de » en chinois \\
\hline
Verbe transitif direct & \zh{帮助他们} bāngzhù tāmen & « les aider » / « leur aider » : pas de « à » ni de COI \\
\hline
Adjectif + \zh{的} + nom & \zh{严重的问题} yánzhòng de wèntí & « un problème grave » : \zh{的} relie l'adjectif au nom \\
\hline
Verbe + compl. de direction & \zh{走上错误的道路} zǒushàng cuòwù de dàolù & « s'engager sur une mauvaise voie » : \zh{上} marque la direction \\
\hline
\end{tabularx}
\end{center}

\begin{attention}[Erreurs fréquentes des francophones]
\begin{itemize}
\item \textbf{Oublier} \zh{的} après un adjectif de deux syllabes : on écrit \zh{严重的问题} et non \zh{严重问题} dans un texte soigné (sauf noms composés lexicalisés comme \zh{法律教育}).
\item \textbf{Traduire mot à mot} « commettre un crime » : le chinois utilise le verbe unique \zh{犯罪} (fànzuì), sans « commettre » (\zh{犯} contient déjà l'idée) ni article.
\item \textbf{Confondre} \zh{错误} (cuòwù, adj./n. formel : erroné, erreur) et \zh{错} (cuò, adj./adv. courant : faux, à tort) : \zh{错误的道路} (une voie erronée) mais \zh{你错了} (tu as tort) / \zh{说错了} (mal dit).
\item \textbf{Ajouter une préposition} après \zh{关心}, \zh{帮助}, \zh{缺少} : ces verbes sont transitifs directs.
\end{itemize}
\end{attention}

\begin{exoresolu}[Associer et compléter]
\textbf{Énoncé.} Complétez les phrases suivantes avec la collocation qui convient : \zh{解决问题}, \zh{走上错误的道路}, \zh{关心孩子}, \zh{加强教育}.\\
1. 家长应该多……。\\
2. 没有家庭教育的孩子容易……。\\
3. 学校应该……，教学生认识法律。\\
4. 只要大家一起努力，就能……。
\tcblower
\textbf{Solution détaillée.}\\
1. 家长应该多\textbf{关心孩子}。 (Jiāzhǎng yīnggāi duō guānxīn háizi.) « Les parents doivent se soucier davantage de leurs enfants. » — Le sujet est « les parents », l'action attendue est « se soucier de ».\\
2. 没有家庭教育的孩子容易\textbf{走上错误的道路}。 (Méiyǒu jiātíng jiàoyù de háizi róngyì zǒushàng cuòwù de dàolù.) « Les enfants sans éducation familiale s'engagent facilement sur une mauvaise voie. » — \zh{容易} (facilement) annonce une conséquence négative.\\
3. 学校应该\textbf{加强教育}，教学生认识法律。 (Xuéxiào yīnggāi jiāqiáng jiàoyù, jiāo xuésheng rènshi fǎlǜ.) « L'école doit renforcer l'éducation et enseigner la loi aux élèves. » — Le sujet est « l'école », l'action institutionnelle est « renforcer ».\\
4. 只要大家一起努力，就能\textbf{解决问题}。 (Zhǐyào dàjiā yìqǐ nǔlì, jiù néng jiějué wèntí.) « Tant que tout le monde fait des efforts ensemble, on peut résoudre le problème. » — La structure \zh{只要……就……} (tant que… alors…) appelle un résultat positif : « résoudre le problème ».
\end{exoresolu}

\begin{methode}[Comment mémoriser ce lexique efficacement]
\begin{enumerate}
\item Apprenez chaque mot avec sa \textbf{collocation}, jamais isolément : retenez \zh{解决问题} comme un seul bloc, pas \zh{解决} tout seul.
\item Classez les mots en quatre familles mentales (problème / causes / actions / acteurs), comme sur la carte mentale.
\item Pour chaque famille, fabriquez une phrase personnelle ancrée au Cameroun, par exemple : \zh{雅温得的学校很关心学生。} (Yǎwēndé de xuéxiào hěn guānxīn xuésheng. « Les écoles de Yaoundé se soucient beaucoup des élèves. »)
\item Vérifiez systématiquement la construction syntaxique : verbe + objet direct (sans préposition), adjectif + \zh{的} + nom, structure \zh{只要……就……}.
\end{enumerate}
\end{methode}

\begin{aretenir}[L'essentiel de la section]
\begin{itemize}
\item Le champ lexical du thème s'organise autour de quatre pôles : le problème (\zh{问题}, \zh{严重}, \zh{犯罪}), les causes (\zh{缺少}, \zh{家庭教育}), les actions (\zh{关心}, \zh{加强}, \zh{帮助}, \zh{解决}, \zh{努力}) et les acteurs (\zh{家长}, \zh{学校}, \zh{警察}, \zh{法院}).
\item Les quatre collocations à connaître par cœur : \zh{解决问题}, \zh{走上错误的道路}, \zh{关心孩子}, \zh{加强教育}.
\item \zh{问题} signifie « problème » ou « question » selon le verbe associé ; \zh{关心}, \zh{帮助}, \zh{缺少} prennent leur objet directement, sans préposition.
\end{itemize}
\end{aretenir}

Maintenant que vous maîtrisez le lexique et les collocations du thème, la section suivante vous apprendra à écrire correctement les caractères de ce vocabulaire : traits, radicaux et caractères proches à ne pas confondre.

\coursec{Écriture correcte des caractères}

Dans la section précédente, nous avons découvert le lexique du thème « délinquance juvénile » : \zh{青少年} (qīngshàonián, adolescent), \zh{犯罪} (fànzuì, commettre un crime), \zh{法律} (fǎlǜ, loi), \zh{错误} (cuòwù, erreur), \zh{道德} (dàodé, morale), \zh{解决} (jiějué, résoudre). Mais savoir prononcer et comprendre ces mots ne suffit pas : un élève de Première A doit aussi pouvoir les \cle{écrire correctement}, trait par trait. En chinois, une erreur de trait ou de composant peut changer complètement le sens : écrire \zh{清} au lieu de \zh{青} transforme « adolescent » en « jeune clair », ce qui ne veut rien dire. Cette section vous donne les outils pour écrire juste, du premier coup.

\courssub{Observer avant d'écrire : les caractères sont construits}

\begin{textebox}[Petit texte d'observation]
青少年犯罪是一个社会问题。有的青少年犯了错误，不知道法律，走上了错误的道路。家长和老师应该想办法帮助他们，解决他们的问题。\\
Qīngshàonián fànzuì shì yí ge shèhuì wèntí. Yǒude qīngshàonián fàn le cuòwù, bù zhīdào fǎlǜ, zǒushàng le cuòwù de dàolù. Jiāzhǎng hé lǎoshī yīnggāi xiǎng bànfǎ bāngzhù tāmen, jiějué tāmen de wèntí.\\
« La délinquance juvénile est un problème de société. Certains adolescents commettent des erreurs, ne connaissent pas la loi et s'engagent sur une mauvaise voie. Les parents et les enseignants doivent chercher des moyens de les aider et de résoudre leurs problèmes. »
\end{textebox}

Observez les caractères clés du texte : \zh{法}, \zh{误}, \zh{想}, \zh{道}, \zh{犯}, \zh{罪}, \zh{解}, \zh{青}. Aucun n'est un dessin arbitraire : chacun est \cle{assemblé} à partir de pièces plus petites. Comprendre ces pièces, c'est déjà mémoriser le caractère.

\courssub{Les radicaux clés du thème}

Un caractère chinois se compose le plus souvent de deux parties : un \cle{radical} (ou clé), qui donne une indication de \textbf{sens}, et une partie \cle{phonétique}, qui donne une indication de \textbf{prononciation}. Voici les radicaux essentiels de notre thème.

\begin{definition}[Les quatre radicaux de la leçon]
\begin{itemize}
\item \zh{氵} « eau » (trois gouttes, forme abrégée de \zh{水} shuǐ) : on le trouve à gauche de \zh{法} fǎ (loi). Idée ancienne : la loi doit être plane et juste comme la surface de l'eau.
\item \zh{讠} « parole » (forme abrégée de \zh{言} yán) : on le trouve à gauche de \zh{误} wù (erreur). Une erreur est souvent liée à ce qu'on dit ou à ce qu'on comprend.
\item \zh{心} « cœur » : on le trouve en bas de \zh{想} xiǎng (penser, vouloir). En Chine ancienne, on pensait avec le cœur !
\item \zh{辶} « mouvement, marche » : il enveloppe \zh{道} dào (chemin, voie). Un chemin sert à marcher, à se déplacer.
\end{itemize}
\end{definition}

\begin{center}
\begin{tabularx}{\linewidth}{|X|X|X|X|}
\hline
\rowcolor{popblueL} Radical & Sens du radical & Caractère du thème & Traduction \\
\hline
氵 (eau) & l'eau, la liquidité & 法 fǎ & loi \\
\hline
讠 (parole) & parler, le langage & 误 wù & erreur, se tromper \\
\hline
心 (cœur) & pensée, sentiment & 想 xiǎng & penser, vouloir \\
\hline
辶 (mouvement) & marcher, se déplacer & 道 dào & chemin, voie \\
\hline
犭 (animal) & animal, sauvagerie & 犯 fàn & violer, commettre \\
\hline
罒 (filet) & filet, prise & 罪 zuì & crime, péché \\
\hline
\end{tabularx}
\end{center}

\begin{savaistu}[D'où vient 罪 ?]
Le caractère \zh{罪} zuì (crime) est composé de \zh{罒} (un filet, en haut) et de \zh{非} fēi (faux, mauvais, en bas). Image parlante : celui qui fait le mal (\zh{非}) finit pris dans le filet (\zh{罒}) de la justice. Retenir cette histoire, c'est ne plus jamais oublier le caractère.
\end{savaistu}

\begin{methode}[Mémoriser un caractère en trois étapes]
\begin{enumerate}
\item \textbf{Identifier le radical} : quelle pièce donne le sens ? (ex. \zh{法} → \zh{氵} eau).
\item \textbf{Identifier la partie phonétique} : quelle pièce donne le son ? (ex. \zh{法} → \zh{去} qù, proche de fǎ par la finale).
\item \textbf{Écrire 5 fois en comptant les traits} à voix haute : « un, deux, trois… » en respectant l'ordre officiel. La mémoire de la main aide la mémoire de la tête.
\end{enumerate}
\end{methode}

\courssub{La règle d'or : l'ordre des traits}

\begin{propriete}[Règles générales de l'ordre des traits]
L'écriture chinoise suit des règles strictes, valables pour presque tous les caractères :
\begin{enumerate}
\item De \textbf{haut en bas} : \zh{三} s'écrit trait du haut, trait du milieu, trait du bas.
\item De \textbf{gauche à droite} : \zh{法} s'écrit d'abord \zh{氵} (à gauche), puis \zh{去} (à droite).
\item L'\textbf{extérieur avant l'intérieur}, puis on ferme. Exception classique : les enveloppes « marchantes » \zh{辶} et \zh{廴} se tracent toujours \textbf{en dernier} — pour \zh{道}, on écrit donc d'abord \zh{首} (l'intérieur), puis l'enveloppe \zh{辶}.
\item L'\textbf{horizontal avant le vertical} quand ils se croisent : \zh{十} = 一 puis 丨.
\item Le \textbf{trait central avant les côtés} dans les caractères symétriques : \zh{小} = 亅 puis les deux points.
\end{enumerate}
\end{propriete}

\begin{attention}[Pourquoi l'ordre des traits est-il obligatoire ?]
Les francophones pensent souvent : « si le résultat ressemble, l'ordre importe peu ». Faux ! À l'examen, l'ordre des traits peut être demandé ; surtout, un mauvais ordre produit des caractères déformés et illisibles, et empêche d'utiliser un dictionnaire par nombre de traits. Écrire dans le bon ordre, c'est écrire vite, beau et juste.
\end{attention}

\courssub{Étude détaillée des caractères clés}

\textbf{1. 青 qīng (vert/bleu, jeune) — 8 traits}\par
Ordre : 一 一 丨 一 (la partie supérieure), puis 丨 (trait brisé) 一 一 (la partie inférieure, proche de \zh{月} « lune »). On écrit donc d'abord les trois horizontales et la verticale du haut, puis la partie du bas, de l'extérieur vers l'intérieur. Attention : dans \zh{青}, la partie du bas commence par une \textbf{verticale} 丨, et non par un 撇 tombant à gauche comme dans \zh{月} écrit seul.

\textbf{2. 犯 fàn (violer, commettre) — 5 traits}\par
D'abord le radical \zh{犭} (3 traits : un 撇 tombant à gauche, puis le crochet courbé ㇁ qui traverse, puis un second 撇) ; ensuite la partie droite \zh{㔾} (2 traits : (trait brisé) puis le crochet vertical). Structure : gauche → droite.

\textbf{3. 罪 zuì (crime) — 13 traits}\par
D'abord \zh{罒} le filet en haut (5 traits : 丨 (trait brisé) 丨 丨 一), puis \zh{非} en bas (8 traits : la verticale de gauche 丨, puis les trois horizontales de gauche, puis la verticale de droite 丨, puis les trois horizontales de droite). Structure : haut → bas.

\textbf{4. 法 fǎ (loi) — 8 traits}\par
D'abord \zh{氵} (3 traits : point, point, puis remontée oblique ㇀), ensuite \zh{去} (5 traits : 一 丨 一 (trait brisé) 丶). Structure : gauche → droite.

\textbf{5. 解 jiě (résoudre, expliquer) — 13 traits}\par
D'abord \zh{角} à gauche (7 traits, la « corne »), puis \zh{刀} en haut à droite (2 traits : (trait brisé) puis 丿), puis \zh{牛} en bas à droite (4 traits : 丿 一 一 丨). Structure : gauche → droite, et à droite haut → bas.

\begin{popfigure}
\begin{center}
\begin{tikzpicture}[scale=1, every node/.style={font=\normalsize}]
% 法
\node[draw=popblue, fill=popblueL, rounded corners, minimum width=1.1cm, minimum height=1.1cm] (a1) at (0,0) {\zh{氵}};
\node at (1.1,0) {$+$};
\node[draw=popgreen, fill=popgreenL, rounded corners, minimum width=1.1cm, minimum height=1.1cm] (a2) at (2.2,0) {\zh{去}};
\node at (3.3,0) {$=$};
\node[draw=popdark, fill=popcream, rounded corners, minimum width=1.3cm, minimum height=1.1cm] at (4.5,0) {\zh{法}};
\node[align=center, color=popmuted] at (2.2,-0.95) {eau $+$ aller $=$ loi (fǎ)};
% 罪
\node[draw=poporange, fill=poporangeL, rounded corners, minimum width=1.1cm, minimum height=1.1cm] (b1) at (0,-2.2) {\zh{罒}};
\node at (1.1,-2.2) {$+$};
\node[draw=poppurple, fill=poppurpleL, rounded corners, minimum width=1.1cm, minimum height=1.1cm] (b2) at (2.2,-2.2) {\zh{非}};
\node at (3.3,-2.2) {$=$};
\node[draw=popdark, fill=popcream, rounded corners, minimum width=1.3cm, minimum height=1.1cm] at (4.5,-2.2) {\zh{罪}};
\node[align=center, color=popmuted] at (2.2,-3.15) {filet $+$ mal $=$ crime (zuì)};
% 误
\node[draw=popblue, fill=popblueL, rounded corners, minimum width=1.1cm, minimum height=1.1cm] (c1) at (0,-4.4) {\zh{讠}};
\node at (1.1,-4.4) {$+$};
\node[draw=popgold, fill=popgoldL, rounded corners, minimum width=1.1cm, minimum height=1.1cm] (c2) at (2.2,-4.4) {\zh{吴}};
\node at (3.3,-4.4) {$=$};
\node[draw=popdark, fill=popcream, rounded corners, minimum width=1.3cm, minimum height=1.1cm] at (4.5,-4.4) {\zh{误}};
\node[align=center, color=popmuted] at (2.2,-5.35) {parole $+$ phonétique wú $=$ erreur (wù)};
\end{tikzpicture}
\end{center}
\legende{Décomposition de trois caractères du thème : le radical (sens) $+$ la partie phonétique.}
\end{popfigure}

\courssub{La grille 田字格 : écrire droit et proportionné}

Les écoliers chinois — et vous aussi — s'entraînent dans le \zh{田字格} (tiánzìgé), la « case champ » : un carré divisé en quatre par une ligne verticale et une ligne horizontale en pointillés. Chaque trait a sa place : les radicaux de gauche (\zh{氵}, \zh{讠}, \zh{犭}) occupent le tiers gauche, la phonétique les deux tiers droits.

\begin{popfigure}
\begin{center}
\begin{tikzpicture}[scale=1.6]
% Grille 青
\draw[thick, popdark] (0,0) rectangle (2,2);
\draw[dashed, popmuted] (1,0) -- (1,2);
\draw[dashed, popmuted] (0,1) -- (2,1);
\draw[dashed, popmuted] (0,0) -- (2,2);
\draw[dashed, popmuted] (0,2) -- (2,0);
\node[font=\Huge, color=popblue] at (1,1) {\zh{青}};
\node[below, align=center] at (1,0) {\zh{青} qīng : partie haute\\ au-dessus de la ligne, bas en dessous};
% Grille 犯
\draw[thick, popdark] (3.5,0) rectangle (5.5,2);
\draw[dashed, popmuted] (4.5,0) -- (4.5,2);
\draw[dashed, popmuted] (3.5,1) -- (5.5,1);
\draw[dashed, popmuted] (3.5,0) -- (5.5,2);
\draw[dashed, popmuted] (3.5,2) -- (5.5,0);
\node[font=\Huge, color=poporange] at (4.5,1) {\zh{犯}};
\node[below, align=center] at (4.5,0) {\zh{犯} fàn : \zh{犭} étroit à gauche,\\ \zh{㔾} large à droite};
\end{tikzpicture}
\end{center}
\legende{Le 田字格 : les pointillés aident à placer chaque trait et à équilibrer le caractère.}
\end{popfigure}

\courssub{Caractères proches à ne pas confondre}

\begin{attention}[Pièges classiques des francophones]
\begin{itemize}
\item \zh{犯} fàn (commettre) $\neq$ \zh{范} fàn (modèle) : même pinyin, même partie phonétique, mais radicaux différents (\zh{犭} animal contre \zh{艹} herbe). On écrit \zh{犯罪} (commettre un crime), jamais \zh{范罪} !
\item \zh{青} qīng $\neq$ \zh{清} qīng (clair, pur) : \zh{清} a trois gouttes d'eau \zh{氵} en plus. Erreur fréquente : écrire \zh{清少年} au lieu de \zh{青少年} — le sens devient absurde (« jeunes clairs »).
\item \zh{法} fǎ (loi) $\neq$ \zh{去} qù (aller) : \zh{法} = \zh{氵} + \zh{去}. Sans les gouttes d'eau, ce n'est plus la loi !
\item \zh{道} dào (chemin) $\neq$ \zh{首} shǒu (tête, premier) : \zh{道} contient \zh{首} enveloppé par \zh{辶}. Sans l'enveloppe du mouvement, il n'y a plus de chemin.
\item \zh{错} cuò (faux) $\neq$ \zh{措} cuò (disposer, dans \zh{措施} « mesure ») : \zh{错} prend le radical du métal \zh{钅}, \zh{措} celui de la main \zh{扌}. « Erreur » s'écrit \zh{错误}, jamais \zh{措误}.
\end{itemize}
\end{attention}

\begin{exemplebox}[Exemples traduits avec les bons caractères]
\zh{青少年犯罪是一个严重的问题。} \emph{Qīngshàonián fànzuì shì yí ge yánzhòng de wèntí.} « La délinquance juvénile est un problème grave. »\\
\zh{他知道自己犯了错误。} \emph{Tā zhīdào zìjǐ fàn le cuòwù.} « Il sait qu'il a commis une erreur. » (\zh{误} avec \zh{讠}, le radical de la parole.)\\
\zh{我们应该想办法解决这个问题。} \emph{Wǒmen yīnggāi xiǎng bànfǎ jiějué zhège wèntí.} « Nous devons chercher un moyen de résoudre ce problème. »\\
\zh{每个人都应该遵守法律。} \emph{Měi ge rén dōu yīnggāi zūnshǒu fǎlǜ.} « Chacun doit respecter la loi. »
\end{exemplebox}

\begin{exoresolu}[Corriger les caractères fautifs]
Énoncé : un élève a écrit la phrase suivante avec quatre caractères faux. Trouvez-les et corrigez-les : \zh{清少年范了措误，不知道去律。}
\tcblower
Solution détaillée :
\begin{enumerate}
\item \zh{清} → \zh{青} : on dit \zh{青少年} (adolescent), sans l'eau \zh{氵}.
\item \zh{范} → \zh{犯} : « commettre » s'écrit avec \zh{犭} : \zh{犯了}.
\item \zh{措} → \zh{错} : « erreur » s'écrit \zh{错误} avec le radical du métal \zh{钅} dans \zh{错} (et \zh{讠} dans \zh{误}).
\item \zh{去} → \zh{法} : « loi » s'écrit \zh{法律} avec \zh{氵}.
\end{enumerate}
Phrase corrigée : \zh{青少年犯了错误，不知道法律。} \emph{Qīngshàonián fàn le cuòwù, bù zhīdào fǎlǜ.} « L'adolescent a commis une erreur et ne connaît pas la loi. »
\end{exoresolu}

\courssub{Entraînement : copie guidée et dictée}

\begin{exercice}[Copie guidée]
Sur votre cahier, tracez des cases 田字格 et copiez \textbf{3 lignes} de chaque caractère clé, en comptant les traits à voix haute : \zh{青} (8 traits), \zh{犯} (5 traits), \zh{罪} (13 traits), \zh{法} (8 traits), \zh{解} (13 traits), \zh{误}, \zh{想}, \zh{道}. Vérifiez ensuite que vos radicaux sont bien placés : \zh{氵} et \zh{犭} étroits à gauche, \zh{心} en bas, \zh{辶} tracé en dernier.
\end{exercice}

\begin{exercice}[Dictée de 8 mots du thème]
L'enseignant dicte ; vous écrivez les caractères, le pinyin et la traduction : 1. \zh{青少年} 2. \zh{犯罪} 3. \zh{法律} 4. \zh{错误} 5. \zh{道德} 6. \zh{解决} 7. \zh{想办法} 8. \zh{道路}.
\end{exercice}

\begin{corrige}[Dictée de 8 mots]
1. \zh{青少年} qīngshàonián — adolescent ; 2. \zh{犯罪} fànzuì — commettre un crime ; 3. \zh{法律} fǎlǜ — loi ; 4. \zh{错误} cuòwù — erreur ; 5. \zh{道德} dàodé — morale ; 6. \zh{解决} jiějué — résoudre ; 7. \zh{想办法} xiǎng bànfǎ — chercher un moyen ; 8. \zh{道路} dàolù — chemin, voie. Barème indicatif : 1 point par mot (0,5 caractères + 0,5 pinyin/traduction).
\end{corrige}

\begin{aretenir}[L'essentiel de la section]
Un caractère = un radical (sens) + une phonétique (son). On écrit de haut en bas, de gauche à droite, l'extérieur avant l'intérieur, et l'enveloppe \zh{辶} en dernier. Les paires à surveiller absolument : \zh{犯}/\zh{范}, \zh{青}/\zh{清}, \zh{法}/\zh{去}, \zh{道}/\zh{首}, \zh{错}/\zh{措}. Une écriture correcte est la première marque de sérieux d'une production écrite — et elle vous fera gagner des points dans la section « Production guidée » qui vient.
\end{aretenir}

\coursec{Structure et connecteurs du texte argumentatif}

Dans la section précédente, nous avons appris à écrire correctement les caractères du thème de la délinquance juvénile : \zh{青少年} (qīngshàonián, adolescent), \zh{犯罪} (fànzuì, commettre un crime/délit), \zh{关心} (guānxīn, se soucier de), \zh{教育} (jiàoyù, éducation). Mais un bon texte argumentatif ne se résume pas à de beaux caractères : il faut encore \cle{organiser} ses idées et les \cle{relier} entre elles. C'est exactement le rôle des connecteurs. En chinois comme en français, un texte sans connecteurs ressemble à une liste de courses ; avec eux, il devient un raisonnement. Cette section vous donne le plan type en quatre parties et les cinq outils grammaticaux indispensables pour rédiger votre production écrite sur la délinquance juvénile.

\courssub{Le plan type en quatre parties}

\begin{textebox}[Un petit texte modèle]
现在，青少年犯罪是一个严重的问题。\\
Xiànzài, qīngshàonián fànzuì shì yí ge yánzhòng de wèntí.\\
« Aujourd'hui, la délinquance juvénile est un problème grave. »\\[4pt]
因为有些青少年缺少家庭教育，所以他们走上了错误的道路。\\
Yīnwèi yǒuxiē qīngshàonián quēshǎo jiātíng jiàoyù, suǒyǐ tāmen zǒushàng le cuòwù de dàolù.\\
« Parce que certains jeunes manquent d'éducation familiale, ils s'engagent sur une mauvaise voie. »\\[4pt]
我认为，首先，家长应该多关心孩子；其次，学校应该加强教育；最后，政府应该帮助困难的家庭。\\
Wǒ rènwéi, shǒuxiān, jiāzhǎng yīnggāi duō guānxīn háizi; qícì, xuéxiào yīnggāi jiāqiáng jiàoyù; zuìhòu, zhèngfǔ yīnggāi bāngzhù kùnnán de jiātíng.\\
« Je pense que, d'abord, les parents devraient se soucier davantage de leurs enfants ; ensuite, l'école devrait renforcer l'éducation ; enfin, le gouvernement devrait aider les familles en difficulté. »\\[4pt]
只要大家一起努力，这个问题一定会解决。\\
Zhǐyào dàjiā yìqǐ nǔlì, zhège wèntí yídìng huì jiějué.\\
« Du moment que tout le monde fait des efforts ensemble, ce problème sera sûrement résolu. »
\end{textebox}

\begin{definition}[Les quatre parties du texte argumentatif]
Ce petit texte modèle suit le plan classique de l'argumentation en chinois :
\begin{enumerate}
\item \zh{问题} (wèntí) : le \cle{problème} — on présente le sujet ;
\item \zh{原因} (yuányīn) : la \cle{cause} — on explique pourquoi le problème existe ;
\item \zh{办法} (bànfǎ) : les \cle{solutions} — on propose des mesures ;
\item \zh{结论} (jiélùn) : la \cle{conclusion} — on termine par un espoir ou un appel.
\end{enumerate}
\end{definition}

\begin{aretenir}[Le plan à mémoriser]
Pour tout sujet d'expression écrite argumentative en Première (délinquance juvénile, drogue, travail des enfants…), retenez le schéma : \textbf{问题 → 原因 → 办法 → 结论}. Une ou deux phrases chinoises par partie suffisent à ce niveau. Chaque partie s'ouvre idéalement par un connecteur.
\end{aretenir}

\courssub{Le connecteur de cause : 因为…所以…}

Observez d'abord ces deux phrases :

\begin{exemplebox}[Observation]
\zh{因为有些青少年缺少家庭教育，所以他们走上了错误的道路。}\\
\emph{Yīnwèi yǒuxiē qīngshàonián quēshǎo jiātíng jiàoyù, suǒyǐ tāmen zǒushàng le cuòwù de dàolù.}\\
« Parce que certains jeunes manquent d'éducation familiale, ils s'engagent sur une mauvaise voie. »\\[4pt]
\zh{因为家里很穷，所以这个孩子不上学，在街上卖东西。}\\
\emph{Yīnwèi jiāli hěn qióng, suǒyǐ zhège háizi bù shàngxué, zài jiēshang mài dōngxi.}\\
« Parce que sa famille est très pauvre, cet enfant ne va pas à l'école et vend des choses dans la rue. »
\end{exemplebox}

\begin{propriete}[Règle de 因为…所以…]
Structure : \textbf{因为 + cause，所以 + résultat。}
\begin{itemize}
\item \zh{因为} (yīnwèi, « parce que ») introduit la proposition de \cle{cause} ;
\item \zh{所以} (suǒyǐ, « donc, c'est pourquoi ») introduit la proposition de \cle{résultat} ;
\item en chinois, les deux connecteurs peuvent — et même devraient — apparaître \textbf{ensemble} dans la même phrase. C'est la grande différence avec le français, où « parce que » et « donc » ne se cumulent jamais ;
\item on peut omettre \zh{所以} (phrase plus légère), mais garder les deux est toujours correct et apprécié à l'examen ;
\item l'ordre cause → résultat est l'ordre normal ; la virgule chinoise ，sépare les deux propositions.
\end{itemize}
\end{propriete}

\begin{popfigure}
\begin{tikzpicture}[node distance=0.6cm and 1.2cm]
\node[draw=popblue, fill=popblueL, rounded corners, thick, align=center, minimum width=3.2cm, minimum height=1.1cm] (yin) {因为 yīnwèi\\ \emph{parce que}};
\node[draw=poporange, fill=poporangeL, rounded corners, thick, align=center, minimum width=3.6cm, minimum height=1.1cm, right=of yin] (cause) {proposition 1\\ \textbf{la cause}};
\node[draw=popgreen, fill=popgreenL, rounded corners, thick, align=center, minimum width=3.2cm, minimum height=1.1cm, right=of cause] (suo) {所以 suǒyǐ\\ \emph{donc}};
\node[draw=poppurple, fill=poppurpleL, rounded corners, thick, align=center, minimum width=3.6cm, minimum height=1.1cm, right=of suo] (res) {proposition 2\\ \textbf{le résultat}};
\draw[-{Stealth}, very thick, popblue] (yin) -- (cause);
\draw[-{Stealth}, very thick, popgreen] (cause) -- (suo);
\draw[-{Stealth}, very thick, poppurple] (suo) -- (res);
\node[below=0.35cm of cause, align=center, text=popmuted] (ex) {\zh{有些青少年缺少家庭教育}};
\node[below=0.35cm of res, align=center, text=popmuted] (ex2) {\zh{他们走上了错误的道路}};
\draw[dashed, popmuted] (ex) -- (cause);
\draw[dashed, popmuted] (ex2) -- (res);
\end{tikzpicture}
\legende{Schéma de la structure 因为 + cause，所以 + résultat}
\end{popfigure}

\begin{attention}[Piège des francophones]
L'erreur numéro un : traduire « parce que… donc… » mot à mot en français mental, puis \textbf{omettre} \zh{所以} parce qu'en français « donc » serait de trop. Faux ! En chinois, la double marque \zh{因为…所以…} est \cle{naturelle}, élégante et valorisée à l'examen. À l'inverse, n'écrivez jamais \zh{因为…但是…} : « parce que » et « mais » ne s'associent pas. Autre erreur fréquente : inverser l'ordre (\zh{所以…因为…}) — gardez toujours la cause d'abord.
\end{attention}

\courssub{Les connecteurs d'énumération : 首先…其次…最后…}

Pour présenter plusieurs solutions (la partie \zh{办法} du plan), le chinois utilise une énumération en trois temps, exactement comme « d'abord… ensuite… enfin… ».

\begin{propriete}[Énumérer les solutions]
Structure : \textbf{首先 + solution 1 ；其次 + solution 2 ；最后 + solution 3。}
\begin{itemize}
\item \zh{首先} (shǒuxiān) : « d'abord, en premier lieu » ;
\item \zh{其次} (qícì) : « ensuite, en second lieu » ;
\item \zh{最后} (zuìhòu) : « enfin, pour finir ».
\end{itemize}
Chaque connecteur se place en tête de proposition, suivi d'une virgule. Entre les propositions, on utilise le point-virgule chinois ；. S'il n'y a que deux solutions, on dit \zh{首先…然后…} (d'abord… puis…).
\end{propriete}

\begin{popfigure}
\begin{tikzpicture}
\draw[-{Stealth}, very thick, popline] (0,0) -- (12.5,0);
\foreach \x/\mot/\py/\fr/\col in {0.8/首先/shǒuxiān/d'abord/popblue, 6.2/其次/qícì/ensuite/poporange, 11.6/最后/zuìhòu/enfin/popgreen}{
  \fill[\col] (\x,0) circle (0.12);
  \node[above=0.25cm, align=center] at (\x,0.25){\textbf{\mot}\\ \emph{\py}\\ \fr};
}
\node[below=0.3cm, align=center, text=popmuted, font=\small] at (0.8,-0.3) {solution 1\\ \zh{家长多关心孩子}};
\node[below=0.3cm, align=center, text=popmuted, font=\small] at (6.2,-0.3) {solution 2\\ \zh{学校加强教育}};
\node[below=0.3cm, align=center, text=popmuted, font=\small] at (11.6,-0.3) {solution 3\\ \zh{政府帮助家庭}};
\end{tikzpicture}
\legende{Frise des connecteurs d'énumération : 首先 → 其次 → 最后}
\end{popfigure}

\begin{exemplebox}[Exemples traduits]
\zh{首先，家长应该多关心孩子。}\\
\emph{Shǒuxiān, jiāzhǎng yīnggāi duō guānxīn háizi.}\\
« D'abord, les parents devraient se soucier davantage de leurs enfants. »\\[4pt]
\zh{其次，学校应该教学生什么是对的，什么是错的。}\\
\emph{Qícì, xuéxiào yīnggāi jiāo xuésheng shénme shì duì de, shénme shì cuò de.}\\
« Ensuite, l'école devrait enseigner aux élèves ce qui est juste et ce qui est faux. »\\[4pt]
\zh{最后，政府应该给年轻人更多的工作机会。}\\
\emph{Zuìhòu, zhèngfǔ yīnggāi gěi niánqīngrén gèng duō de gōngzuò jīhuì.}\\
« Enfin, le gouvernement devrait offrir plus d'opportunités d'emploi aux jeunes. »
\end{exemplebox}

\courssub{Exprimer son opinion : 我认为 / 我相信}

Un texte argumentatif n'est pas seulement une liste de faits : l'élève doit dire ce qu'il \cle{pense}. Deux verbes sont parfaits à ce niveau.

\begin{definition}[Les verbes d'opinion]
\begin{itemize}
\item \zh{我认为} (wǒ rènwéi) : « je pense que, j'estime que » — le plus courant à l'écrit, légèrement formel ;
\item \zh{我相信} (wǒ xiāngxìn) : « je crois que » — exprime la conviction et l'espoir, idéal dans la conclusion.
\end{itemize}
Structure : \textbf{我认为 / 我相信 + proposition complète (sujet + verbe).} Le connecteur \zh{ que} du français ne se traduit pas : on enchaîne directement la proposition.
\end{definition}

\begin{exemplebox}[Exemples traduits]
\zh{我认为，家庭教育是最重要的。}\\
\emph{Wǒ rènwéi, jiātíng jiàoyù shì zuì zhòngyào de.}\\
« Je pense que l'éducation familiale est la plus importante. »\\[4pt]
\zh{我相信，雅温得的年轻人可以建设更美好的社会。}\\
\emph{Wǒ xiāngxìn, Yǎwēndé de niánqīngrén kěyǐ jiànshè gèng měihǎo de shèhuì.}\\
« Je crois que les jeunes de Yaoundé peuvent construire une société meilleure. »
\end{exemplebox}

\begin{attention}[Ne traduisez pas « que »]
Les francophones cherchent souvent un mot pour « que » après \zh{我认为}. Il n'y en a pas : \zh{我认为这个问题很严重} (et jamais un mot-outil entre les deux). De même, ne confondez pas \zh{认为} (penser, estimer — opinion réfléchie) et \zh{想} (penser, avoir envie — plus spontané). À l'écrit argumentatif, préférez \zh{认为}.
\end{attention}

\courssub{Le connecteur de condition : 只要…就…}

Ce connecteur est un excellent « plus » pour la conclusion : il exprime qu'\textbf{une seule condition suffit} pour obtenir le résultat.

\begin{propriete}[Règle de 只要…就…]
Structure : \textbf{只要 + condition，就 + résultat。}
\begin{itemize}
\item \zh{只要} (zhǐyào) : « du moment que, pourvu que » ;
\item \zh{就} (jiù) : « alors » — placé juste avant le verbe (ou l'auxiliaire) de la seconde proposition ;
\item sens : la condition énoncée est \cle{suffisante} pour que le résultat se réalise ;
\item équivalent français : « du moment que… alors… », « pourvu que… ».
\end{itemize}
\end{propriete}

\begin{exemplebox}[Exemples traduits]
\zh{只要大家一起努力，这个问题一定会解决。}\\
\emph{Zhǐyào dàjiā yìqǐ nǔlì, zhège wèntí yídìng huì jiějué.}\\
« Du moment que tout le monde fait des efforts ensemble, ce problème sera sûrement résolu. »\\[4pt]
\zh{只要家长多关心孩子，孩子就不会走上错误的道路。}\\
\emph{Zhǐyào jiāzhǎng duō guānxīn háizi, háizi jiù bú huì zǒushàng cuòwù de dàolù.}\\
« Pourvu que les parents se soucient davantage de leurs enfants, ceux-ci ne s'engageront pas sur une mauvaise voie. »
\end{exemplebox}

\courssub{Exprimer la solution : le modal 应该 + verbe}

Dans la partie \zh{办法} (solutions), chaque proposition a besoin d'un verbe modal de \cle{devoir} ou de \cle{conseil} : c'est \zh{应该} (yīnggāi, « devrait, il faudrait »).

\begin{propriete}[Le modal 应该]
Structure : \textbf{Sujet + 应该 + verbe (+ complément).}
\begin{itemize}
\item \zh{应该} exprime le devoir moral, la recommandation : « devrait » ;
\item la négation est \zh{不应该} (bù yīnggāi, « ne devrait pas ») ;
\item il se place toujours \textbf{avant} le verbe principal, jamais après ;
\item sujets typiques dans notre thème : \zh{家长} (les parents), \zh{学校} (l'école), \zh{政府} (le gouvernement), \zh{我们} (nous).
\end{itemize}
\end{propriete}

\begin{exemplebox}[Exemples traduits]
\zh{家长应该多关心孩子。}\\
\emph{Jiāzhǎng yīnggāi duō guānxīn háizi.}\\
« Les parents devraient se soucier davantage de leurs enfants. »\\[4pt]
\zh{青少年不应该喝酒，也不应该吸烟。}\\
\emph{Qīngshàonián bù yīnggāi hē jiǔ, yě bù yīnggāi xī yān.}\\
« Les adolescents ne devraient ni boire d'alcool ni fumer. »\\[4pt]
\zh{我们应该帮助有困难的同学。}\\
\emph{Wǒmen yīnggāi bāngzhù yǒu kùnnán de tóngxué.}\\
« Nous devrions aider les camarades en difficulté. »
\end{exemplebox}

\courssub{Tableau récapitulatif et comparaison chinois / français}

\begin{center}
\begin{tabularx}{\linewidth}{|X|X|X|X|}
\hline
\rowcolor{popblueL} Caractères & Pinyin & Nature & Traduction \\
\hline
因为…所以… & yīnwèi… suǒyǐ… & Structure corrélative (cause-résultat) & parce que… (donc)… \\
\hline
首先 & shǒuxiān & Adverbe (énumération) & d'abord \\
\hline
其次 & qícì & Adverbe (énumération) & ensuite \\
\hline
最后 & zuìhòu & Adverbe (énumération) & enfin \\
\hline
我认为 & wǒ rènwéi & Verbe d'opinion & je pense que \\
\hline
我相信 & wǒ xiāngxìn & Verbe de conviction & je crois que \\
\hline
只要…就… & zhǐyào… jiù… & Structure corrélative (condition-résultat) & du moment que… alors… \\
\hline
应该 & yīnggāi & Verbe modal & devrait, il faudrait \\
\hline
\end{tabularx}
\end{center}

\begin{center}
\begin{tabularx}{\linewidth}{|X|X|X|}
\hline
\rowcolor{popgreenL} Point de comparaison & Chinois & Français \\
\hline
Double marque de cause & \zh{因为} et \zh{所以} ensemble : naturel & « parce que » et « donc » : jamais ensemble \\
\hline
« que » après « je pense » & ne se traduit pas & obligatoire \\
\hline
Place du modal & avant le verbe : \zh{应该关心} & « devrait se soucier » (deux verbes) \\
\hline
Énumération & \zh{首先、其次、最后} fixes & d'abord, ensuite, enfin (plus souples) \\
\hline
\end{tabularx}
\end{center}

\courssub{Méthode de rédaction et entraînement}

\begin{methode}[Rédiger en quatre étapes]
\begin{enumerate}
\item \textbf{Plan en français} : sur le brouillon, écrivez quatre mots : problème – cause – solutions – conclusion. Notez une idée simple pour chaque partie.
\item \textbf{Une phrase chinoise par partie} : rédigez une phrase courte et correcte pour chaque partie, en puisant dans le lexique de la leçon (\zh{青少年}, \zh{家庭}, \zh{教育}, \zh{关心}, \zh{帮助}…).
\item \textbf{Un connecteur par phrase} : ouvrez la cause avec \zh{因为…所以…}, les solutions avec \zh{首先…其次…最后…} + \zh{应该}, la conclusion avec \zh{我相信} ou \zh{只要…就…}.
\item \textbf{Relecture} : vérifiez les tons du pinyin au brouillon, l'ordre des mots (sujet + modal + verbe), la présence de \zh{所以} après \zh{因为}, et l'écriture des caractères.
\end{enumerate}
\end{methode}

\begin{exoresolu}[Relier avec le bon connecteur]
Énoncé : complétez la phrase avec le connecteur qui convient, puis traduisez : \zh{＿＿＿学校加强教育，＿＿＿学生能分清对错。}
\tcblower
Solution détaillée : la phrase oppose une condition (« l'école renforce l'éducation ») à un résultat attendu (« les élèves distingueront le bien du mal »). Le connecteur de condition suffisante convient parfaitement : \zh{只要…就…}.\\
Phrase complète : \zh{只要学校加强教育，学生就能分清对错。}\\
\emph{Zhǐyào xuéxiào jiāqiáng jiàoyù, xuésheng jiù néng fēnqīng duìcuò.}\\
Traduction : « Du moment que l'école renforce l'éducation, les élèves pourront distinguer le juste du faux. »\\
Remarque : \zh{因为…所以…} serait grammaticalement possible, mais il présenterait le lien comme un fait déjà acquis ; \zh{只要…就…} exprime mieux l'espoir d'une solution, ce qui convient à une conclusion.
\end{exoresolu}

\begin{exercice}[À vous de jouer]
Remettez les quatre phrases suivantes dans l'ordre du plan 问题 → 原因 → 办法 → 结论, puis traduisez-les en français :\\
1. \zh{只要我们一起努力，社会会更美好。}\\
2. \zh{青少年犯罪是一个严重的问题。}\\
3. \zh{首先，家长应该多关心孩子。}\\
4. \zh{因为有些孩子缺少父母的爱，所以他们犯了错误。}
\end{exercice}

\begin{corrige}[Exercice « À vous de jouer »]
Ordre correct : 2 → 4 → 3 → 1.\\
2. \zh{青少年犯罪是一个严重的问题。} « La délinquance juvénile est un problème grave. » (问题)\\
4. \zh{因为有些孩子缺少父母的爱，所以他们犯了错误。} « Parce que certains enfants manquent d'amour parental, ils commettent des erreurs. » (原因, avec 因为…所以…)\\
3. \zh{首先，家长应该多关心孩子。} « D'abord, les parents devraient se soucier davantage de leurs enfants. » (办法, avec 首先 + 应该)\\
1. \zh{只要我们一起努力，社会会更美好。} « Du moment que nous faisons des efforts ensemble, la société sera meilleure. » (结论, avec 只要…就…)
\end{corrige}

Vous disposez maintenant de tous les outils pour structurer un raisonnement en chinois : le plan en quatre parties, la double marque de cause \zh{因为…所以…}, l'énumération \zh{首先…其次…最后…}, les verbes d'opinion \zh{我认为} et \zh{我相信}, la condition \zh{只要…就…} et le modal des solutions \zh{应该}. La section suivante vous entraînera à mobiliser ces connecteurs dans la traduction du français vers le chinois (thème).

\coursec{Thème : traduction français → chinois}

Dans la section précédente, nous avons appris à structurer un texte argumentatif et à utiliser les connecteurs \zh{因为……所以……}, \zh{应该}, \zh{但是}. Nous allons maintenant mettre ces outils en pratique dans l'exercice roi de l'expression écrite au lycée : le \cle{thème}, c'est-à-dire la traduction du français vers le chinois. Le thème exige deux compétences : trouver le bon mot chinois (le lexique de la délinquance juvénile que nous avons étudié) et le placer au bon endroit dans la phrase (la syntaxe). C'est cette deuxième compétence qui fait la différence à l'examen.

\begin{definition}[Lexique clé du thème : la délinquance juvénile]
\begin{center}
\begin{tabularx}{\linewidth}{|X|X|X|X|}
\hline
\rowcolor{popblueL} \textbf{Caractères} & \textbf{Pinyin} & \textbf{Nature} & \textbf{Traduction} \\
\hline
青少年犯罪 & qīngshàonián fànzuì & nom & délinquance juvénile / criminalité des mineurs \\
\hline
问题 & wèntí & nom & problème \\
\hline
严重 & yánzhòng & adj & grave, sérieux \\
\hline
缺少 & quēshǎo & verbe & manquer de, faire défaut (transitif direct) \\
\hline
家庭教育 & jiātíng jiàoyù & nom & éducation familiale \\
\hline
加强 & jiāqiáng & verbe & renforcer, intensifier \\
\hline
法律教育 & fǎlǜ jiàoyù & nom & éducation juridique \\
\hline
走上 & zǒushàng & verbe & emprunter (une voie), s'engager dans \\
\hline
错误 & cuòwù & adj/nom & erroné, faux / erreur \\
\hline
道路 & dàolù & nom & route, voie, chemin \\
\hline
帮助 & bāngzhù & verbe & aider, assistance (transitif direct) \\
\hline
解决 & jiějué & verbe & résoudre, régler \\
\hline
相信 & xiāngxìn & verbe & croire, avoir confiance en \\
\hline
一定 & yídìng & adv & certainement, sans doute \\
\hline
\end{tabularx}
\end{center}
\end{definition}

\begin{propriete}[Écriture correcte : traits, radicaux et caractères proches]
\begin{itemize}
\item \textbf{犯} : radical \zh{犭} (chien/quadrupède) + \zh{然}. Ordre : traits du radical (3 traits : point, courbe, point) puis \zh{然} (haut → bas, gauche → droite). \emph{Attention :} ne pas confondre avec \zh{范} (modèle, radical \zh{艹}).
\item \textbf{罪} : radical \zh{网} (filet) + \zh{非} + \zh{辛}. Le haut \zh{网} s'écrit : horizontales, verticale, horizontale de fermeture, puis contenu. \emph{Caractère proche :} \zh{辠} (forme traditionnelle de 罪).
\item \textbf{法} : radical \zh{氵} (eau) + \zh{去}. 3 traits pour l'eau (point, point, trait montant), puis \zh{去} (haut → bas).
\item \textbf{律} : radical \zh{彳} (pas) + \zh{聿} (pinceau). \zh{彳} : 3 traits (horizontal, courbe, horizontal). \emph{Ne pas confondre avec} \zh{率} (taux, radical \zh{玄}/\zh{亠}).
\item \textbf{缺} : radical \zh{缶} (vase) + \zh{夬} (décider). Haut → bas. \emph{Caractère proche :} \zh{阙} (porte de palais, radical \zh{门}).
\item \textbf{错} : radical \zh{钅} (métal) + \zh{昔}. \zh{钅} : 5 traits (point, horizontale, verticale, point, horizontale). \emph{Ne pas confondre avec} \zh{措} (mesure, radical \zh{扌} main).
\item \textbf{解} : radical \zh{角} (corne) + \zh{刀} (sabre) + \zh{牛} (bœuf). Structure complexe : haut → bas, gauche → droite.
\end{itemize}
\end{propriete}

\begin{propriete}[Connecteurs utiles pour l'argumentation (niveau HSK 3-4)]
\begin{center}
\begin{tabularx}{\linewidth}{|X|X|X|}
\hline
\rowcolor{popblueL} \textbf{Connecteur (Chinois + Pinyin)} & \textbf{Sens / Fonction} & \textbf{Exemple court} \\
\hline
\zh{而且} \emph{érqiě} & De plus, en outre (addition) & \zh{他缺少教育，而且不守法律。} \\
\hline
\zh{但是} \emph{dànshì} / \zh{可是} \emph{kěshì} & Mais, cependant (opposition) & \zh{问题严重，但是我们有办法。} \\
\hline
\zh{因此} \emph{yīncǐ} / \zh{所以} \emph{suǒyǐ} & Par conséquent, donc (conséquence) & \zh{他犯罪了，因此必须受罚。} \\
\hline
\zh{首先} \emph{shǒuxiān} & Premièrement (énumération) & \zh{首先，家庭应该负责。} \\
\hline
\zh{其次} \emph{qícì} & Deuxièmement (énumération) & \zh{其次，学校要加强教育。} \\
\hline
\zh{总之} \emph{zǒngzhī} & En bref, en conclusion & \zh{总之，社会要帮助他们。} \\
\hline
\end{tabularx}
\end{center}
\end{propriete}

\courssub{Observer avant de traduire : l'ordre des mots}

Avant toute méthode, observons deux phrases déjà traduites :

\begin{exemplebox}[Observation : où sont les mots ?]
\zh{学校应该加强法律教育。}\\
\emph{Xuéxiào yīnggāi jiāqiáng fǎlǜ jiàoyù.}\\
« L'école devrait renforcer l'éducation juridique. »

\medskip
\zh{因为这些年轻人走上了错误的道路，所以社会应该帮助他们。}\\
\emph{Yīnwèi zhèxiē niánqīngrén zǒushàng le cuòwù de dàolù, suǒyǐ shèhuì yīnggāi bāngzhù tāmen.}\\
« Parce que ces jeunes ont pris une mauvaise voie, la société doit les aider. »
\end{exemplebox}

Que remarque-t-on ? En français, « devrait » est un verbe conjugué ; en chinois, \zh{应该} est un \cle{verbe modal} invariable placé \textbf{avant le verbe principal} et \textbf{après le sujet}. De même, le complément de cause introduit par \zh{因为} se place \textbf{au début} de la phrase, jamais à la fin comme en français (« la société doit les aider parce que… » est impossible à calquer mot à mot).

\begin{propriete}[L'ordre fondamental de la phrase chinoise]
La phrase chinoise suit l'ordre rigide :

\medskip
\textbf{Sujet + Temps + Lieu + Modal (\zh{应该}, \zh{会}, \zh{能}, \zh{一定}) + Verbe + Objet}

\medskip
Contrairement au français, le chinois ne conjugue pas : c'est la \textbf{position} des mots qui porte le sens. Tout déplacement change ou détruit le sens de la phrase.
\end{propriete}

\begin{popfigure}
\begin{tikzpicture}[
  bloc/.style={rectangle, rounded corners=3pt, draw=popblue, fill=popblueL, align=center, minimum height=0.9cm, font=\small, text width=2.2cm},
  fleche/.style={-{Stealth[length=2.5mm]}, thick, popdark}
]
\node[bloc, align=center] (s) at (0,0){\textbf{Sujet}\\ 学校};
\node[bloc, fill=popgoldL, draw=popgold, align=center] (t) at (2.8,0){\textbf{Temps}\\ 现在};
\node[bloc, fill=poporangeL, draw=poporange, align=center] (l) at (5.6,0){\textbf{Lieu}\\ 在学校};
\node[bloc, fill=popgreenL, draw=popgreen, align=center] (m) at (8.4,0){\textbf{Modal}\\ 应该};
\node[bloc, fill=poppinkL, draw=poppink, align=center] (v) at (11.2,0){\textbf{Verbe}\\ 加强};
\node[bloc, fill=popblueL, draw=popblue, align=center] (o) at (14,0){\textbf{Objet}\\ 法律教育};
\draw[fleche] (s) -- (t);
\draw[fleche] (t) -- (l);
\draw[fleche] (l) -- (m);
\draw[fleche] (m) -- (v);
\draw[fleche] (v) -- (o);
\node[align=center, font=\small\itshape, popmuted] at (7,-1.3) {Xuéxiào xiànzài zài xuéxiào yīnggāi jiāqiáng fǎlǜ jiàoyù. — « L'école maintenant à l'école devrait renforcer l'éducation juridique. »};
\end{tikzpicture}
\legende{Le train de la phrase chinoise : chaque wagon (Temps, Lieu, Modal) a sa place fixe avant le Verbe.}
\end{popfigure}

\courssub{La méthode du thème en quatre étapes}

\begin{methode}[Traduire du français vers le chinois : les 4 étapes]
\begin{enumerate}
\item \textbf{Souligner les mots-clés} de la phrase française : qui ? (sujet) fait quoi ? (verbe) quoi / qui ? (objet) quand ? où ? (compléments circonstanciels).
\item \textbf{Trouver l'équivalent chinois} de chaque mot-clé dans le lexique du thème (délinquance juvénile) : \zh{青少年犯罪}, \zh{问题}, \zh{缺少}, \zh{家庭教育}, \zh{加强}, \zh{帮助}, \zh{走上}, \zh{错误}, \zh{道路}, \zh{解决}…
\item \textbf{Ordonner selon S + Temps + Lieu + Modal + V + O} : placer les compléments de temps et de lieu \emph{avant} le modal/verbe, le modal juste avant le verbe principal. Pour les phrases complexes : \zh{因为} + Cause (sujet + verbe), \zh{所以} + Conséquence (sujet + modal + verbe).
\item \textbf{Vérifier les caractères} : traits complets, bons radicaux, ne pas confondre les caractères proches (\zh{未} / \zh{末}, \zh{己} / \zh{已}, \zh{钅} vs \zh{扌}) ; relire en prononçant le pinyin pour vérifier que la phrase « sonne » chinois.
\end{enumerate}
\end{methode}

\courssub{Les cinq phrases du thème, pas à pas}

\textbf{Phrase 1 : « La délinquance juvénile est un problème grave. »}

Étape 1 : sujet = « la délinquance juvénile » ; verbe = « est » (verbe d'état/identité) ; attribut = « un problème grave ». Étape 2 : \zh{青少年犯罪} + \zh{是} + \zh{一个严重的问题}. Étape 3 : phrase nominale Sujet + \zh{是} + Nom (classificateur \zh{个}).

\zh{青少年犯罪是一个严重的问题。}\\
\emph{Qīngshàonián fànzuì shì yí ge yánzhòng de wèntí.}

\begin{attention}[Le piège n° 1 : oublier \zh{是} devant un nom]
En français, « est » semble léger, on l'oublie. En chinois, la phrase nominale (Sujet = Nom) exige \textbf{obligatoirement} \zh{是} entre le sujet et le nom attribut : « est un problème grave » = \zh{是一个严重的问题}. La phrase \zh{青少年犯罪一个严重的问题} est \textbf{fausse} : sans \zh{是}, elle ne veut rien dire. \\
\textbf{Exception :} avec un adjectif seul comme attribut, on n'utilise pas \zh{是} mais souvent l'adverbe de degré \zh{很} : \zh{这个问题很严重。} \emph{Zhège wèntí hěn yánzhòng.} « Ce problème est grave. »
\end{attention}

\textbf{Phrase 2 : « Certains adolescents manquent d'éducation familiale. »}

« Certains » = \zh{有些} (suivi directement du nom, sans article, pas de \zh{的} nécessaire ici). « Manquer de » = \zh{缺少} + objet direct, sans préposition (verbe transitif direct).

\zh{有些青少年缺少家庭教育。}\\
\emph{Yǒuxiē qīngshàonián quēshǎo jiātíng jiàoyù.}

\textbf{Phrase 3 : « L'école devrait renforcer l'éducation juridique. »}

« Devrait » = le modal \zh{应该}, placé entre le sujet et le verbe \zh{加强}. Complément de lieu absent, temps absent.

\zh{学校应该加强法律教育。}\\
\emph{Xuéxiào yīnggāi jiāqiáng fǎlǜ jiàoyù.}

\textbf{Phrase 4 : « Parce que ces jeunes ont pris une mauvaise voie, la société doit les aider. »}

Le couple de connecteurs \zh{因为……所以……} encadre les deux propositions ; la cause en tête. « Prendre une mauvaise voie » = \zh{走上错误的道路} (verbe \zh{走上} + aspect accompli \zh{了} + objet). « Les » = \zh{他们} placé après le verbe \zh{帮助}.

\zh{因为这些年轻人走上了错误的道路，所以社会应该帮助他们。}\\
\emph{Yīnwèi zhèxiē niánqīngrén zǒushàng le cuòwù de dàolù, suǒyǐ shèhuì yīnggāi bāngzhù tāmen.}

\textbf{Phrase 5 (niveau examen) : « Je crois que ce problème sera résolu. »}

« Croire que » = \zh{相信} + proposition complète (sujet + verbe) ; « sera résolu » (passif futur / certitude) se rend par \zh{一定会解决} : l'adverbe \zh{一定} « certainement » et le modal \zh{会} « futur / certitude » précèdent le verbe \zh{解决}.

\zh{我相信这个问题一定会解决。}\\
\emph{Wǒ xiāngxìn zhège wèntí yídìng huì jiějué.}

\courssub{Tableau récapitulatif des correspondances}

\begin{center}
\begin{tabularx}{\linewidth}{|X|X|X|}
\hline
\rowcolor{popblueL} \textbf{Français} & \textbf{Chinois + pinyin} & \textbf{Point de grammaire} \\
\hline
La délinquance juvénile est un problème grave. & 青少年犯罪是一个严重的问题。 \emph{Qīngshàonián fànzuì shì yí ge yánzhòng de wèntí.} & \zh{是} obligatoire devant un nom attribut \\
\hline
Certains adolescents manquent d'éducation familiale. & 有些青少年缺少家庭教育。 \emph{Yǒuxiē qīngshàonián quēshǎo jiātíng jiàoyù.} & \zh{缺少} + objet direct (sans préposition) \\
\hline
L'école devrait renforcer l'éducation juridique. & 学校应该加强法律教育。 \emph{Xuéxiào yīnggāi jiāqiáng fǎlǜ jiàoyù.} & Modal \zh{应该} avant le verbe principal \\
\hline
Parce que ces jeunes ont pris une mauvaise voie, la société doit les aider. & 因为这些年轻人走上了错误的道路，所以社会应该帮助他们。 \emph{Yīnwèi zhèxiē niánqīngrén zǒushàng le cuòwù de dàolù, suǒyǐ shèhuì yīnggāi bāngzhù tāmen.} & \zh{因为……所以……} en tête ; \zh{了} aspect accompli \\
\hline
Je crois que ce problème sera résolu. & 我相信这个问题一定会解决。 \emph{Wǒ xiāngxìn zhège wèntí yídìng huì jiějué.} & \zh{相信} + proposition ; \zh{一定会} = certitude future \\
\hline
\end{tabularx}
\end{center}

\begin{exemplebox}[Le modal \zh{应该} avant le verbe principal]
\zh{我们应该帮助这些年轻人。}\\
\emph{Wǒmen yīnggāi bāngzhù zhèxiē niánqīngrén.}\\
« Nous devons aider ces jeunes. »

\medskip
Autres exemples : \zh{青少年应该遵守法律。} \emph{Qīngshàonián yīnggāi zūnshǒu fǎlǜ.} « Les adolescents doivent respecter la loi. » — \zh{家长应该多关心孩子。} \emph{Jiāzhǎng yīnggāi duō guānxīn háizi.} « Les parents devraient s'occuper davantage de leurs enfants. »
\end{exemplebox}

\begin{attention}[Deux pièges de francophones]
\begin{itemize}
\item \textbf{Pas de préposition calquée} : \zh{帮助} prend directement un objet : \zh{帮助他们} « les aider ». Ne jamais écrire une préposition imitant le français « aider à ». De même : \zh{缺少家庭教育} « manquer d'éducation familiale », sans « de ». \zh{加强法律教育} « renforcer l'éducation juridique », sans « de ».
\item \textbf{Place de \zh{应该}} : il se place \emph{avant} le verbe principal, jamais après : on dit \zh{社会应该帮助他们}, jamais \zh{社会帮助应该他们}. Avec la négation : \zh{不应该} + verbe : \zh{我们不应该放弃这些年轻人。} \emph{Wǒmen bù yīnggāi fàngqì zhèxiē niánqīngrén.} « Nous ne devrions pas abandonner ces jeunes. »
\end{itemize}
\end{attention}

\begin{exoresolu}[Thème guidé : deux phrases]
Énoncé. Traduisez en chinois : 1) « Les parents doivent aider leurs enfants. » 2) « Parce que certains jeunes manquent d'éducation familiale, ils prennent une mauvaise voie. »
\tcblower
Solution détaillée.

\textbf{Phrase 1.} Étape 1 : sujet « les parents » (\zh{家长} ou \zh{父母}), modal « doivent » (\zh{应该}), verbe « aider » (\zh{帮助}), objet « leurs enfants » (\zh{他们的孩子}). Étape 3 : S + Modal + V + O.

\zh{家长应该帮助他们的孩子。} \emph{Jiāzhǎng yīnggāi bāngzhù tāmen de háizi.}

\textbf{Phrase 2.} Étape 1 : la cause vient en tête avec \zh{因为}. Sujet cause : « certains jeunes » (\zh{些青少年} / \zh{些年轻人}). Verbe cause : « manquent d'éducation familiale » (\zh{缺少家庭教育}). Conséquence : sujet « ils » (\zh{他们}), verbe « prennent une mauvaise voie » (\zh{走上了错误的道路}).

\zh{因为有些青少年缺少家庭教育，所以他们走上了错误的道路。}\\
\emph{Yīnwèi yǒuxiē qīngshàonián quēshǎo jiātíng jiàoyù, suǒyǐ tāmen zǒushàng le cuòwù de dàolù.}

Vérification (étape 4) : \zh{应该} bien avant \zh{帮助} ; pas de préposition après \zh{帮助} ni après \zh{缺少} ; caractères \zh{错误} (radical \zh{钅}) et \zh{道路} correctement tracés ; \zh{了} après \zh{走上} pour l'accompli.
\end{exoresolu}

\begin{exercice}[À vous de traduire]
Traduisez en chinois, en appliquant les quatre étapes de la méthode :
\begin{enumerate}
\item La délinquance juvénile est un problème grave au Cameroun.
\item L'école et la famille doivent aider ces jeunes.
\item Je crois que la société devrait renforcer l'éducation juridique.
\end{enumerate}
\end{exercice}

\begin{corrige}[Exercice : thème]
1) \zh{青少年犯罪在喀麦隆是一个严重的问题。} \emph{Qīngshàonián fànzuì zài Kāmàilóng shì yí ge yánzhòng de wèntí.} — Le complément de lieu \zh{在喀麦隆} se place \emph{avant} le verbe \zh{是} (ordre S + Lieu + V).

2) \zh{学校和家庭应该帮助这些年轻人。} \emph{Xuéxiào hé jiātíng yīnggāi bāngzhù zhèxiē niánqīngrén.} — Double sujet relié par \zh{和} ; \zh{应该} avant \zh{帮助} ; \zh{这些} démonstratif pluriel.

3) \zh{我相信社会应该加强法律教育。} \emph{Wǒ xiāngxìn shèhuì yīnggāi jiāqiáng fǎlǜ jiàoyù.} — \zh{相信} introduit une proposition complète (S + Modal + V + O) ; pas de \zh{的} après \zh{社会} car sujet de la subordonnée.
\end{corrige}

\begin{aretenir}[L'essentiel du thème]
\begin{itemize}
\item Ordre rigide : \textbf{Sujet + Temps + Lieu + Modal + Verbe + Objet}.
\item Phrase nominale (Sujet = Nom) : \zh{是} obligatoire devant un nom attribut ; \zh{很} avec un adjectif seul.
\item Modaux \zh{应该}, \zh{会}, \zh{能}, \zh{一定} se placent \emph{avant} le verbe principal. Négation : \zh{不} + Modal + V.
\item \zh{因为……所以……} : la cause en tête de phrase principale.
\item Verbes transitifs directs : \zh{帮助}, \zh{缺少}, \zh{加强}, \zh{解决}, \zh{遵守} prennent l'objet sans préposition.
\item Connecteurs d'argumentation : \zh{而且}, \zh{但是}, \zh{因此}, \zh{首先}, \zh{其次}, \zh{总之}.
\end{itemize}
\end{aretenir}

\coursec{Version : traduction chinois → français}

Dans la section précédente, nous avons appris à traduire du français vers le chinois (\emph{thème}) : il fallait penser en chinois, choisir les bons mots-outils et respecter l'ordre des mots chinois. Nous faisons maintenant le chemin inverse : la \cle{version}, c'est-à-dire la traduction du chinois vers le français. L'exercice semble plus facile — on comprend « à peu près » le chinois — mais il cache un piège redoutable : le \emph{calque}, c'est-à-dire la phrase française qui copie mot à mot la structure chinoise et qui sonne faux. Une bonne version exige deux qualités : la \textbf{fidélité au sens} et la \textbf{qualité du français}.

\courssub{Observer d'abord : le texte de version}

Lisons d'abord le texte à traduire, sans chercher encore à traduire. Lisons-le entièrement, à voix haute si possible.

\begin{textebox}[Texte de version : 青少年问题]
现在，很多青少年缺少家长的关心。因为他们常常犯错误，所以学校应该加强教育。我相信，只要家长、学校和社会一起努力，这个问题一定会解决。\\
\emph{Xiànzài, hěn duō qīngshàonián quēshǎo jiāzhǎng de guānxīn. Yīnwèi tāmen chángcháng fàn cuòwù, suǒyǐ xuéxiào yīnggāi jiāqiáng jiàoyù. Wǒ xiāngxìn, zhǐyào jiāzhǎng, xuéxiào hé shèhuì yìqǐ nǔlì, zhège wèntí yídìng huì jiějué.}
\end{textebox}

\begin{center}
\begin{tabularx}{\linewidth}{|X|X|X|X|}
\hline
\rowcolor{popblueL} Caractères & Pinyin & Nature & Traduction \\
\hline
青少年 & qīngshàonián & nom & adolescent, jeune \\
\hline
缺少 & quēshǎo & verbe & manquer de \\
\hline
家长 & jiāzhǎng & nom & parents (de famille) \\
\hline
关心 & guānxīn & nom / verbe & attention, sollicitude ; se soucier de \\
\hline
犯错误 & fàn cuòwù & locution verbale & commettre des erreurs \\
\hline
加强 & jiāqiáng & verbe & renforcer \\
\hline
教育 & jiàoyù & nom / verbe & éducation ; éduquer \\
\hline
相信 & xiāngxìn & verbe & croire, être convaincu \\
\hline
只要 & zhǐyào & conjonction & pourvu que, si seulement \\
\hline
社会 & shèhuì & nom & société \\
\hline
一起 & yìqǐ & adverbe & ensemble \\
\hline
努力 & nǔlì & verbe / adjectif & faire des efforts ; appliqué \\
\hline
解决 & jiějué & verbe & résoudre \\
\hline
一定 & yídìng & adverbe & certainement, sûrement \\
\hline
\end{tabularx}
\end{center}

\courssub{La méthode de la version en quatre étapes}

\begin{methode}[Réussir une version en quatre étapes]
\begin{enumerate}
\item \textbf{Lire le texte entier} avant de traduire la moindre phrase. On doit pouvoir dire en une phrase française de quoi parle le texte. Ici : \emph{la délinquance juvénile et les moyens d'y remédier}.
\item \textbf{Repérer les connecteurs} : \zh{因为……所以……} (comme / parce que … donc …), \zh{只要} (pourvu que), \zh{我相信} (je crois que). Ils donnent l'architecture logique du texte : cause → conséquence → conviction.
\item \textbf{Traduire par blocs de sens}, jamais mot à mot. Un bloc = un groupe de mots qui forme une unité de sens : \zh{缺少家长的关心} se traduit d'un seul geste : « manquer de l'attention de leurs parents ».
\item \textbf{Relire la traduction française à voix haute} : est-ce du vrai français ? Est-ce qu'un Français dirait cela naturellement ? Si une phrase sonne « chinoise », la reformuler.
\end{enumerate}
\end{methode}

\begin{popfigure}
\begin{tikzpicture}[
  node distance=0.55cm,
  etape/.style={rectangle, rounded corners=3pt, draw=popblue, fill=popblueL, text width=4.6cm, align=center, font=\small, inner sep=5pt},
  fleche/.style={-{Stealth[length=2.5mm]}, thick, poporange}
]
\node[etape, align=center] (e1){\textbf{1. Lire} le texte entier\\ (saisir le sens global)};
\node[etape, below=of e1, align=center] (e2){\textbf{2. Repérer} les connecteurs\\ \zh{因为……所以……、只要、我相信}};
\node[etape, below=of e2, align=center] (e3){\textbf{3. Traduire} par blocs de sens\\ (pas mot à mot)};
\node[etape, below=of e3, align=center] (e4){\textbf{4. Relire} en français\\ (corriger les calques)};
\draw[fleche] (e1) -- (e2);
\draw[fleche] (e2) -- (e3);
\draw[fleche] (e3) -- (e4);
\end{tikzpicture}
\legende{Organigramme des quatre étapes de la version chinois → français}
\end{popfigure}

\courssub{Traduction de référence commentée}

\begin{aretenir}[Traduction de référence]
« Aujourd'hui, beaucoup d'adolescents manquent de l'attention de leurs parents. Comme ils commettent souvent des erreurs, l'école devrait renforcer l'éducation. Je crois que si les parents, l'école et la société travaillent ensemble, ce problème sera sûrement résolu. »
\end{aretenir}

Analysons phrase par phrase comment on passe du chinois au français.

\textbf{Phrase 1.} \zh{现在，很多青少年缺少家长的关心。} (\emph{Xiànzài, hěn duō qīngshàonián quēshǎo jiāzhǎng de guānxīn.})

Le mot de temps \zh{现在} se place en tête en chinois ; en français, « aujourd'hui » peut aussi ouvrir la phrase : aucune difficulté. Le bloc \zh{家长的关心} (« l'attention des parents ») demande un petit travail de naturel : « manquent de l'attention de leurs parents » est plus élégant que « manquent l'attention des parents ».

\textbf{Phrase 2.} \zh{因为他们常常犯错误，所以学校应该加强教育。} (\emph{Yīnwèi tāmen chángcháng fàn cuòwù, suǒyǐ xuéxiào yīnggāi jiāqiáng jiàoyù.})

Piège classique : le chinois utilise \textbf{les deux} connecteurs \zh{因为……所以……} (« parce que … donc … »), alors que le français n'en tolère qu'un seul. Il faut choisir : « \emph{Comme} ils commettent souvent des erreurs, l'école devrait renforcer l'éducation » ou « Ils commettent souvent des erreurs, \emph{donc} l'école devrait renforcer l'éducation ». Dire « parce que … donc … » en français est une faute de calque.

\textbf{Phrase 3.} \zh{我相信，只要家长、学校和社会一起努力，这个问题一定会解决。} (\emph{Wǒ xiāngxìn, zhǐyào jiāzhǎng, xuéxiào hé shèhuì yìqǐ nǔlì, zhège wèntí yídìng huì jiějué.})

\zh{我相信} introduit une opinion : « je crois que ». \zh{只要} correspond à « si (seulement) / pourvu que ». \zh{一起努力} se rend naturellement par « travailler ensemble » ou « unir leurs efforts ». Enfin, \zh{这个问题一定会解决} est un passif implicite : littéralement « ce problème certainement va se résoudre », donc en bon français : « ce problème sera sûrement résolu ».

\courssub{Les trois points délicats du texte}

\begin{exemplebox}[Collocations et faux amis de traduction]
\begin{itemize}
\item \zh{犯错误} (\emph{fàn cuòwù}) = « commettre des erreurs ». Le verbe \zh{犯} signifie « commettre » (une faute, un crime : \zh{犯罪} « commettre un crime »). On ne peut pas traduire littéralement : « marcher erreur » ou « faire erreur » seraient fautifs. En chinois, on ne dit pas non plus \zh{做错误}.
\item \zh{一起努力} (\emph{yìqǐ nǔlì}) = « faire des efforts ensemble », « unir leurs efforts », « travailler ensemble ». Le français préfère souvent un nom (« efforts ») là où le chinois utilise un verbe.
\item \zh{一定会解决} (\emph{yídìng huì jiějué}) = « sera sûrement résolu ». Le chinois n'a pas de forme passive obligatoire ici ; le français, lui, préfère le passif (« sera résolu ») ou le pronominal (« se résoudra »).
\end{itemize}
\end{exemplebox}

\begin{attention}[会 : certitude future, pas capacité]
Le mot \zh{会} a deux emplois principaux : la \textbf{capacité acquise} (\zh{我会说汉语。} « Je sais parler chinois. ») et la \textbf{certitude future} (\zh{这个问题一定会解决。} « Ce problème sera sûrement résolu. »). Dans notre texte, c'est le second emploi : traduire « ce problème saura se résoudre » serait un contresens. Indice : quand le sujet de \zh{会} est une chose inanimée (\zh{问题}, « problème »), il s'agit presque toujours du futur ou de la probabilité, jamais de la capacité.
\end{attention}

\begin{center}
\begin{tabularx}{\linewidth}{|X|X|X|}
\hline
\rowcolor{popblueL} Bloc chinois & Traduction littérale (à éviter) & Bonne traduction française \\
\hline
\zh{缺少家长的关心} & manquer l'attention des parents & manquer de l'attention de leurs parents \\
\hline
\zh{因为……所以……} & parce que … donc … & comme …, … / …, donc … \\
\hline
\zh{犯错误} & faire erreur & commettre des erreurs \\
\hline
\zh{加强教育} & ajouter-fort l'éducation & renforcer l'éducation \\
\hline
\zh{一起努力} & ensemble effort & unir leurs efforts / travailler ensemble \\
\hline
\zh{一定会解决} & certainement va résoudre & sera sûrement résolu \\
\hline
\end{tabularx}
\end{center}

\courssub{Exercice résolu}

\begin{exoresolu}[Version guidée]
Traduisez en français correct : \zh{因为很多青少年缺少家长的关心，所以他们常常犯错误。} (\emph{Yīnwèi hěn duō qīngshàonián quēshǎo jiāzhǎng de guānxīn, suǒyǐ tāmen chángcháng fàn cuòwù.})
\tcblower
\textbf{Étape 1 — Sens global} : lien de cause à effet entre le manque d'attention parentale et les erreurs des jeunes.\\
\textbf{Étape 2 — Connecteurs} : \zh{因为……所以……} = un seul connecteur en français.\\
\textbf{Étape 3 — Blocs} : \zh{很多青少年} « beaucoup d'adolescents » ; \zh{缺少家长的关心} « manquent de l'attention de leurs parents » ; \zh{常常犯错误} « commettent souvent des erreurs ».\\
\textbf{Étape 4 — Rédaction} : « Comme beaucoup d'adolescents manquent de l'attention de leurs parents, ils commettent souvent des erreurs. »\\
\textbf{À ne pas écrire} : « Parce que beaucoup d'adolescents manquent l'attention des parents, donc ils font souvent des erreurs » (double connecteur + calque de \zh{犯}).
\end{exoresolu}

\courssub{Entraînement : version guidée à trous}

\begin{exercice}[Complétez la traduction]
Traduisez les phrases suivantes en complétant les trous.
\begin{enumerate}
\item \zh{我相信这个问题一定会解决。} (\emph{Wǒ xiāngxìn zhège wèntí yídìng huì jiějué.})\\
« Je \ldots\ldots\ldots\ldots que ce problème sera \ldots\ldots\ldots\ldots résolu. »
\item \zh{只要家长和学校一起努力，青少年就会进步。} (\emph{Zhǐyào jiāzhǎng hé xuéxiào yìqǐ nǔlì, qīngshàonián jiù huì jìnbù.})\\
« \ldots\ldots\ldots\ldots les parents et l'école \ldots\ldots\ldots\ldots leurs efforts, les adolescents progresseront. »
\item \zh{学校应该加强教育，帮助学生不犯错误。} (\emph{Xuéxiào yīnggāi jiāqiáng jiàoyù, bāngzhù xuéshēng bù fàn cuòwù.})\\
« L'école devrait \ldots\ldots\ldots\ldots l'éducation et aider les élèves à ne pas \ldots\ldots\ldots\ldots d'erreurs. »
\end{enumerate}
\end{exercice}

\begin{corrige}[Exercice : version guidée à trous]
\begin{enumerate}
\item « Je \textbf{crois} que ce problème sera \textbf{sûrement} résolu. » (\zh{相信} = croire ; \zh{一定} = sûrement ; \zh{会解决} = futur passif implicite.)
\item « \textbf{Si} les parents et l'école \textbf{unissent} leurs efforts, les adolescents progresseront. » (\zh{只要} = si / pourvu que ; \zh{一起努力} = unir leurs efforts ; \zh{就会} marque la conséquence future.)
\item « L'école devrait \textbf{renforcer} l'éducation et aider les élèves à ne pas \textbf{commettre} d'erreurs. » (\zh{加强} = renforcer ; \zh{犯错误} = commettre des erreurs.)
\end{enumerate}
\end{corrige}

\begin{savaistu}[La version à l'examen : deux critères de notation]
Au baccalauréat et dans les évaluations MINESEC, la version est notée sur \textbf{deux} critères : la fidélité au sens du texte chinois \emph{et} la qualité du français écrit. Une traduction exacte mais rédigée en français maladroit perd des points, tout comme une belle phrase française qui trahit le sens. Astuce de correcteur : relisez votre phrase à voix haute (ou mentalement, en articulant) ; si elle « sonne chinois », c'est qu'un calque s'y cache. Les grands écrivains chinois modernes, comme Lu Xun, ont d'ailleurs longuement débattu de la manière de traduire les langues étrangères en chinois sans déformer ni la langue source ni la langue d'arrivée : le même souci de naturel s'applique à vous, dans l'autre sens !
\end{savaistu}

\begin{aretenir}[L'essentiel de la version]
\begin{itemize}
\item Lire tout le texte, repérer les connecteurs, traduire par blocs de sens, relire en français.
\item \zh{因为……所以……} = un seul connecteur en français (« comme » ou « donc », jamais les deux).
\item \zh{犯错误} = « commettre des erreurs » ; \zh{一起努力} = « unir leurs efforts » ; \zh{一定会解决} = « sera sûrement résolu ».
\item \zh{会} devant un verbe dont le sujet est une chose exprime la certitude future (« sera »), non la capacité.
\item La note dépend du sens \emph{et} de la qualité du français : la relecture à voix haute est votre meilleure alliée.
\end{itemize}
\end{aretenir}

\coursec{Leçon 8 — Expression écrite : la délinquance juvénile}

\courssub{Modèle de texte et compréhension}

\begin{textebox}[Reportage : La délinquance juvénile au Cameroun et en Chine]
\zh{喀麦隆和中国都关注青少年犯罪问题。}\\
\emph{Kāmàilóng hé Zhōngguó dōu guānzhù qīngshàonián fànzuì wèntí.}\\
« Le Cameroun et la Chine s'inquiètent tous deux du problème de la délinquance juvénile. »\\[4pt]
\zh{在喀麦隆的大城市，如雅温得和杜阿拉，一些青少年因为缺少家庭教育，走上了犯罪的道路。}\\
\emph{Zài Kāmàilóng de dà chéngshì, rú Yǎwēndé hé Dùālá, yìxiē qīngshàonián yīnwèi quēshǎo jiātíng jiàoyù, zǒushàng le fànzuì de dàolù.}\\
« Dans les grandes villes du Cameroun, comme Yaoundé et Douala, certains jeunes, par manque d'éducation familiale, s'engagent sur la voie de la délinquance. »\\[4pt]
\zh{中国政府采取了严厉的法律措施，同时加强学校和家庭的教育，青少年犯罪率已经下降。}\\
\emph{Zhōngguó zhèngfǔ cǎiqǔ le yánlì de fǎlǜ cuòshī, tóngshí jiāqiáng xuéxiào hé jiātíng de jiàoyù, qīngshàonián fànzuì lǜ yǐjīng xiàjiàng.}\\
« Le gouvernement chinois a pris des mesures légales strictes, tout en renforçant l'éducation à l'école et dans la famille ; le taux de délinquance juvénile a déjà baissé. »\\[4pt]
\zh{专家认为，解决这个问题，首先要靠父母多关心孩子，其次要靠学校教学生遵守法律，最后要靠社会提供更多机会。}\\
\emph{Zhuānjiā rènwéi, jiějué zhège wèntí, shǒuxiān yào kào fùmǔ duō guānxīn háizi, qícì yào kào xuéxiào jiāo xuésheng zūnshǒu fǎlǜ, zuìhòu yào kào shèhuì tígōng gèng duō jīhuì.}\\
« Les experts estiment que pour résoudre ce problème, il faut d'abord que les parents s'occupent plus de leurs enfants, ensuite que l'école enseigne aux élèves à respecter la loi, enfin que la société offre plus d'opportunités. »
\end{textebox}

\begin{exercice}[Compréhension du texte]
\begin{enumerate}
\item Selon le texte, quelles sont les deux grandes villes camerounaises citées ? Écrivez-les en chinois.
\item Quelle est la cause principale mentionnée pour le Cameroun ? (Répondez en français).
\item Quelles sont les trois mesures prises par la Chine ? Utilisez les connecteurs \zh{首先}、\zh{其次}、\zh{最后} pour répondre en chinois.
\item Traduisez en français : \zh{专家认为，解决这个问题，首先要靠父母多关心孩子。}
\end{enumerate}
\end{exercice}

\begin{corrige}[Compréhension du texte]
\begin{enumerate}
\item \zh{雅温得} (Yaoundé) et \zh{杜阿拉} (Douala).
\item Le manque d'éducation familiale (\zh{缺少家庭教育}).
\item \zh{首先，中国政府采取了严厉的法律措施；其次，加强学校和家庭的教育；最后，青少年犯罪率已经下降。} (Note : le texte cite les mesures gouvernementales comme cause de la baisse).
\item « Les experts estiment que pour résoudre ce problème, il faut avant tout que les parents portent plus d'attention à leurs enfants. »
\end{enumerate}
\end{corrige}

\courssub{Lexique du thème et collocations}

\begin{definition}[Vocabulaire de base — Thème : La délinquance juvénile]
\begin{center}
\begin{tabularx}{\linewidth}{|X|X|X|X|}\hline
\rowcolor{popblueL} \textbf{Caractères} & \textbf{Pinyin} & \textbf{Nature} & \textbf{Traduction} \\ \hline
青少年 & qīngshàonián & n. & jeune, adolescent \\ \hline
犯罪 & fànzuì & v./n. & commettre un crime / délinquance, crime \\ \hline
严重 & yánzhòng & adj. & grave, sérieux \\ \hline
缺少 & quēshǎo & v. & manquer de, faire défaut \\ \hline
应该 & yīnggāi & modal v. & devoir, devrait \\ \hline
解决 & jiějué & v. & résoudre (un problème) \\ \hline
原因 & yuányīn & n. & cause, raison \\ \hline
解决办法 & jiějué bànfǎ & n. & solution, moyen de résoudre \\ \hline
家庭教育 & jiātíng jiàoyù & n. & éducation familiale \\ \hline
法律 & fǎlǜ & n. & loi \\ \hline
遵守 & zūnshǒu & v. & respecter (loi, règle), observer \\ \hline
社会 & shèhuì & n. & société \\ \hline
机会 & jīhuì & n. & occasion, opportunité \\ \hline
关心 & guānxīn & v. & se soucier de, s'intéresser à \\ \hline
陪伴 & péibàn & v. & accompagner, tenir compagnie \\ \hline
错误 & cuòwù & adj./n. & erroné / erreur, faute \\ \hline
道路 & dàolù & n. & route, chemin, voie \\ \hline
下降 & xiàjiàng & v. & baisser, diminuer \\ \hline
采取 & cǎiqǔ & v. & prendre (mesures), adopter \\ \hline
措施 & cuòshī & n. & mesure \\ \hline
专家 & zhuānjiā & n. & expert \\ \hline
\end{tabularx}
\end{center}
\end{definition}

\begin{propriete}[Collocations fréquentes (搭配)]
\begin{center}
\begin{tabularx}{\linewidth}{|X|X|X|}\hline
\rowcolor{popblueL} \textbf{Structure} & \textbf{Exemple (Caractères + Pinyin)} & \textbf{Traduction} \\ \hline
\zh{缺少 + N} & \zh{缺少关心} (quēshǎo guānxīn) & manquer d'attention \\ \hline
\zh{犯罪 + 的 + N} & \zh{犯罪的道路} (fànzuì de dàolù) & la voie de la délinquance \\ \hline
\zh{严重 + 的 + N} & \zh{严重的问题} (yánzhòng de wèntí) & un problème grave \\ \hline
\zh{应该 + V} & \zh{应该解决} (yīnggāi jiějué) & devrait résoudre \\ \hline
\zh{解决 + N} & \zh{解决问题} (jiějué wèntí) & résoudre le problème \\ \hline
\zh{遵守 + 法律/规则} & \zh{遵守法律} (zūnshǒu fǎlǜ) & respecter la loi \\ \hline
\zh{提供 + 机会} & \zh{提供机会} (tígōng jīhuì) & offrir des opportunités \\ \hline
\zh{加强 + 教育} & \zh{加强教育} (jiāqiáng jiàoyù) & renforcer l'éducation \\ \hline
\end{tabularx}
\end{center}
\end{propriete}

\begin{attention}[Faux amis et pièges lexicaux]
\begin{itemize}
\item \zh{犯} (fàn, commettre, clé 犭) \textbf{\neq} \zh{范} (fàn, modèle, norme, clé 艹). On écrit \zh{犯罪}, \zh{犯错误} (faire une erreur), jamais \zh{范罪} ni \zh{范错误}.
\item \zh{青少年} = \zh{青年} (jeunes adultes, 18-30 ans) + \zh{少年} (adolescents, 10-18 ans). Ne pas confondre avec \zh{儿童} (enfants, 0-12 ans).
\item \zh{解决} (résoudre un problème, un conflit) vs \zh{决定} (jué dìng, décider). Sons proches, sens radicalement différents.
\item \zh{严重} (grave, pour situation, problème, conséquence) vs \zh{严厉} (yánlì, sévère, pour personne, règle, ton, mesure).
\end{itemize}
\end{attention}

\courssub{Écriture correcte des caractères}

\begin{methode}[Règles d'ordre des traits — Rappel fondamental]
\begin{enumerate}
\item \textbf{De haut en haut} (从上到下) : \zh{青} (qīng), \zh{严} (yán).
\item \textbf{De gauche à droite} (从左到右) : \zh{缺} (quē), \zh{解} (jiě).
\item \textbf{Horizontal avant vertical} (先横后竖) : \zh{十} (shí), trait horizontal du \zh{土} dans \zh{少}.
\item \textbf{Extérieur avant intérieur, puis fermeture} (先外后内，再封口) : \zh{严} (yán), \zh{囚} (qiú, dans \zh{圉}... non, ici \zh{严} structure haut-bas). Pour \zh{囗} (wei) : \zh{国} (guó).
\item \textbf{Le trait central avant les ailes} (先中间后两边) : \zh{小} (xiǎo), \zh{水} (shuǐ).
\end{enumerate}
\end{methode}

\begin{exemplebox}[Analyse des caractères clés du thème]
\begin{itemize}
\item \zh{\textbf{犯} (fàn)} : Clé \zh{犭} (chien/quadrupède, 3 traits) + phonétique \zh{范} (fàn). Ordre : 犭 (1. courbe, 2. crochet, 3. point) → \zh{范} (herbe 艹 → 車). \textit{Sens : l'animal qui outrepasse les bornes → commettre une faute.}
\item \zh{\textbf{罪} (zuì)} : Clé \zh{网} (filet, 5 traits : horizontales, verticales, crochets) + \zh{非} (fēi, non/erreur). Ordre : \zh{网} (cadre haut, côtés, bas, intérieur) → \zh{非} (haut → bas). \textit{Sens : capturé dans le filet pour une faute.}
\item \zh{\textbf{严} (yán)} : Clé \zh{宀} (toit, 2 points + trait horizontal) + \zh{工} (gōng, travail). Ordre : \zh{宀} → \zh{工} (horizontal, vertical, horizontal). \textit{Structure haut-bas.}
\item \zh{\textbf{重} (zhòng)} : \zh{里} (lǐ) + \zh{里} (lǐ) ? Non, structure : \zh{千} (qiān) + \zh{里} (lǐ) ? En fait : haut \zh{厂} (hàn, falaise) + \zh{东} (dōng) sans le trait du bas ? Standard : \zh{重} = \zh{里} + \zh{里} (simplifié) ou \zh{千} + \zh{里} (traditionnel \zh{重}). Simplifié : \zh{里} (7 traits) + \zh{里} (7 traits) superposés ? Non. Ordre officiel simplifié : 1. verticale gauche \zh{里}, 2. horizontales \zh{里}, 3. verticale droite \zh{里}, 4. horizontales bas \zh{里}. \textit{Mémoriser : "double里".}
\item \zh{\textbf{缺} (quē)} : Clé \zh{缶} (fǒu, pot, 6 traits) + \zh{欠} (qiàn, devoir, 4 traits). Structure gauche-droite. \textit{Ordre : clé \zh{缶} complète → \zh{欠} complète.}
\item \zh{\textbf{解} (jiě)} : Clé \zh{角} (jiǎo, corne, 7 traits) + \zh{刀} (dāo, sabre) + \zh{牛} (niú, bœuf). Structure haut-bas. \textit{Ordre : \zh{角} → \zh{刀} → \zh{牛}. Sens : dénouer les cornes du bœuf avec un sabre → résoudre.}
\item \zh{\textbf{决} (jué)} : Clé \zh{氵} (shuǐ, eau, 3 traits) + \zh{夂} (zhǐ, aller, 3 traits). Gauche-droite. \textit{Ordre : \zh{氵} (points, trait montant) → \zh{夂} (horizontal, diagonal gauche, diagonal droite). Sens : l'eau qui perce le barrage → décider, trancher.}
\end{itemize}
\end{exemplebox}

\begin{exercice}[Entraînement à l'écriture]
\begin{enumerate}
\item Écrivez 5 fois chaque caractère ci-dessous en respectant l'ordre des traits : \zh{犯}, \zh{罪}, \zh{严}, \zh{重}, \zh{缺}, \zh{解}, \zh{决}.
\item Complétez les caractères en ajoutant la clé ou le composant manquant :
      \begin{itemize}
      \item \zh{犭 + \_\_\_\_ = 犯} (réponse : \zh{范})
      \item \zh{网 + \_\_\_\_ = 罪} (réponse : \zh{非})
      \item \zh{宀 + \_\_\_\_ = 严} (réponse : \zh{工})
      \item \zh{氵 + \_\_\_\_ = 决} (réponse : \zh{夂})
      \end{itemize}
\item Identifiez l'erreur de caractère dans chaque phrase et corrigez-la :
      \begin{enumerate}
      \item \zh{青少年范罪是一个严重的问题。}
      \item \zh{我们应该加强家庭教肓。} (\zh{肓} huāng vs \zh{育} yù)
      \item \zh{社会应该提工更多机会。} (\zh{工} gōng vs \zh{供} gōng)
      \end{enumerate}
\end{enumerate}
\end{exercice}

\begin{corrige}[Entraînement à l'écriture]
\begin{enumerate}
\item (Exercice graphomoteur, vérification par l'enseignant).
\item \zh{范}, \zh{非}, \zh{工}, \zh{夂}.
\item 1. \zh{范罪} $\rightarrow$ \zh{犯罪}. 2. \zh{教肓} $\rightarrow$ \zh{教育} (\zh{育} : clé \zh{月} viande/corps). 3. \zh{提工} $\rightarrow$ \zh{提供} (\zh{供} : clé \zh{亻} personne + \zh{共}).
\end{enumerate}
\end{corrige}

\courssub{Structure et connecteurs du texte argumentatif}

\begin{propriete}[Structure \zh{因为……所以……} (yīnwèi... suǒyǐ...)]
\textbf{Schéma :} \zh{因为} + Cause (phrase), \zh{所以} + Conséquence (phrase).\\
\textbf{Emploi :} Lien logique fort cause-effet. \zh{所以} ne s'omet pas à l'écrit formel (contrairement à l'oral où il peut disparaître).
\begin{center}
\begin{tabularx}{\linewidth}{|X|X|X|}\hline
\rowcolor{popblueL} \textbf{Structure} & \textbf{Exemple} & \textbf{Traduction} \\ \hline
\zh{因为 S + V, 所以 S + V} & \zh{因为他缺少爱，所以他犯罪。} & Parce qu'il manque d'amour, il commet des crimes. \\ \hline
\zh{因为原因, 所以结果} & \zh{因为家庭教育缺失，所以青少年走上错误道路。} & À cause du manque d'éducation familiale, les jeunes prennent le mauvais chemin. \\ \hline
\end{tabularx}
\end{center}
\textbf{Variante :} \zh{由于} (yóuyú, "du à", plus formel) + \zh{因此} (yīncǐ, "par conséquent") ou \zh{所以}.
\end{propriete}

\begin{propriete}[Connecteurs d'énumération : \zh{首先 / 其次 / 最后}]
\textbf{Emploi :} Structurer une argumentation ou une liste de solutions. Toujours en position initiale de phrase (sujet optionnel après).
\begin{center}
\begin{tabularx}{\linewidth}{|X|X|X|}\hline
\rowcolor{popblueL} \textbf{Connecteur} & \textbf{Pinyin} & \textbf{Exemple} \\ \hline
\zh{首先} & shǒuxiān & \zh{首先，父母应该多陪伴孩子。} (D'abord, les parents devraient accompagner plus les enfants.) \\ \hline
\zh{其次} & qícì & \zh{其次，学校应该加强法制教育。} (Ensuite, l'école devrait renforcer l'éducation juridique.) \\ \hline
\zh{最后} & zuìhòu & \zh{最后，社会要提供就业机会。} (Enfin, la société doit offrir des opportunités d'emploi.) \\ \hline
\end{tabularx}
\end{center}
\textbf{Attention :} Ne pas confondre \zh{最后} (zuìhòu, "enfin" dans une liste) et \zh{最后} (zuìhòu, "à la fin" temporel). Contexte différent, même mot.
\end{propriete}

\begin{propriete}[Modal \zh{应该} (yīnggāi) + Verbe]
\textbf{Sens :} Obligation morale, conseil, "devrait". Négation : \zh{不应该} (bù yīnggāi).
\textbf{Structure :} Sujet + \zh{应该} + Verbe (+ Complément).
\begin{center}
\begin{tabularx}{\linewidth}{|X|X|X|}\hline
\rowcolor{popblueL} \textbf{Sujet} & \zh{应该} + \textbf{Verbe} & \textbf{Traduction} \\ \hline
\zh{父母} (fùmǔ) & \zh{应该 关心 孩子} & Les parents devraient se soucier des enfants. \\ \hline
\zh{学校} (xuéxiào) & \zh{应该 教育 学生 遵守 法律} & L'école devrait éduquer les élèves à respecter la loi. \\ \hline
\zh{我们} (wǒmen) & \zh{应该 一起 解决 问题} & Nous devrions résoudre le problème ensemble. \\ \hline
\end{tabularx}
\end{center}
\end{propriete}

\begin{propriete}[Structure de conclusion : \zh{我相信} + \zh{只要……就/会……}]
\textbf{Schéma :} \zh{我相信} (wǒ xiāngxìn, "je crois que") + \zh{只要} (zhǐyào, "pourvu que / tant que") + Condition, \zh{就 / 会} (jiù / huì) + Résultat.
\end{propriete}

\begin{exemplebox}[Modèles de conclusion]
\zh{我相信，只要大家一起努力，这个问题一定会解决。}\\
\emph{Wǒ xiāngxìn, zhǐyào dàjiā yìqǐ nǔlì, zhège wèntí yídìng huì jiějué.}\\
« Je crois que si tout le monde fait des efforts ensemble, ce problème sera sûrement résolu. »\\[6pt]
\zh{我相信，只要家庭和学校配合，青少年犯罪就会减少。}\\
\emph{Wǒ xiāngxìn, zhǐyào jiātíng hé xuéxiào pèihé, qīngshàonián fànzuì jiù huì jiǎnshǎo.}\\
« Je crois que si la famille et l'école coopèrent, la délinquance juvénile diminuera. »
\end{exemplebox}


\courssub{Thème : Traduction français → chinois}

\begin{methode}[Méthode du thème (Français → Chinois)]
\begin{enumerate}
\item \textbf{Analyser} la phrase française : identifier Sujet, Verbe, Compléments, Temps/Aspect, Modalité.
\item \textbf{Transposer} dans l'ordre chinois : Sujet + (Temps) + (Lieu/Manière) + Verbe + (Complément d'objet/Résultat/Direction).
\item \textbf{Insérer} les mots-outils (\zh{了}, \zh{过}, \zh{的}, \zh{得}, \zh{地}, classificateurs).
\item \textbf{Vérifier} : accords (pas d'accord), ordre des mots, caractères, ponctuation chinoise.
\end{enumerate}
\end{methode}

\begin{exercice}[Thème — Phrases sur la délinquance juvénile]
Traduisez en chinois (caractères + pinyin).
\begin{enumerate}
\item La délinquance juvénile est un problème très grave dans la société.
\item Parce que certains parents manquent de temps, ils ne peuvent pas s'occuper de leurs enfants.
\item Nous devrions résoudre ce problème : d'abord, la famille doit éduquer les enfants ; ensuite, l'école doit enseigner la loi.
\item Enfin, la société devrait donner plus d'opportunités aux jeunes.
\item Je crois que si tout le monde fait des efforts, la situation s'améliorera.
\end{enumerate}
\end{exercice}

\begin{corrige}[Thème — Corrigé détaillé]
\begin{enumerate}
\item \zh{青少年犯罪是一个很严重的社会问题。} (Qīngshàonián fànzuì shì yí ge hěn yánzhòng de shèhuì wèntí.) — \textit{Note : "dans la société" $\rightarrow$ attribut \zh{社会} devant \zh{问题}. Classificateur \zh{个}.}
\item \zh{因为一些父母缺少时间，所以不能关心孩子。} (Yīnwèi yìxiē fùmǔ quēshǎo shíjiān, suǒyǐ bù néng guānxīn háizi.) — \textit{Structure \zh{因为...所以...}. \zh{不能} (ne peut pas).}
\item \zh{我们应该解决这个问题：首先，家庭应该教育孩子；其次，学校应该教学生遵守法律。} (Wǒmen yīnggāi jiějué zhège wèntí: shǒuxiān, jiātíng yīnggāi jiàoyù háizi; qícì, xuéxiào yīnggāi jiāo xuésheng zūnshǒu fǎlǜ.) — \textit{Point-virgule chinois \zh{；} pour séparer clauses coordonnées longues. \zh{教} (jiāo) ou \zh{教育} (jiàoyù).}
\item \zh{最后，社会应该给青少年提供更多的机会。} (Zuìhòu, shèhuì yīnggāi gěi qīngshàonián tígōng gèng duō de jīhuì.) — \textit{\zh{给} + Bénéficiaire + \zh{提供} + Objet. \zh{更多的} (plus de).}
\item \zh{我相信，只要大家一起努力，情况一定会好转。} (Wǒ xiāngxìn, zhǐyào dàjiā yìqǐ nǔlì, qíngkuàng yídìng huì hǎozhuǎn.) — \textit{\zh{情况} (situation), \zh{好转} (s'améliorer, se retourner en bien). Variante : \zh{会变好} (huì biàn hǎo).}
\end{enumerate}
\end{corrige}

\courssub{Version : Traduction chinois → français}

\begin{methode}[Méthode de la version (Chinois → Français)]
\begin{enumerate}
\item \textbf{Segmenter} la phrase chinoise : repérer Sujet, Verbe, Compléments, Mots-outils (\zh{了, 过, 的, 得, 地, 把, 被, 被}).
\item \textbf{Traduire} mot à mot (glose) pour vérifier la structure.
\item \textbf{Reformuler} en français naturel, correct, au temps adapté (passé composé, présent, futur... selon \zh{了}/\zh{过}/contexte).
\item \textbf{Soigner} le style : éviter le "chinois de cuisine" (ex: "Parce que lui manquer amour, donc il aller mauvais chemin").
\end{enumerate}
\end{methode}

\begin{exercice}[Version — Texte court]
Traduisez en français naturel le texte suivant :
\zh{现在，青少年犯罪已经成为一个严重的社会问题。因为很多青少年缺少父母的陪伴和关心，所以他们容易走上错误的道路。专家建议：首先，父母应该多花时间陪伴孩子；其次，学校应该加强法制教育，教学生遵守法律；最后，政府和社会应该为青少年创造更多就业机会。我相信，只要全社会共同努力，青少年犯罪问题一定能够得到解决。}
\end{exercice}

\begin{corrige}[Version — Proposition de traduction]
« Aujourd'hui, la délinquance juvénile est déjà devenue un grave problème social. Parce que beaucoup de jeunes manquent de la compagnie et de l'attention de leurs parents, ils ont tendance à s'engager sur la mauvaise voie. Les experts suggèrent : d'abord, les parents devraient passer plus de temps à accompagner leurs enfants ; ensuite, l'école devrait renforcer l'éducation juridique, apprendre aux élèves à respecter la loi ; enfin, le gouvernement et la société devraient créer plus d'opportunités d'emploi pour les jeunes. Je crois que pourvu que toute la société fasse des efforts ensemble, le problème de la délinquance juvénile pourra certainement être résolu. »

\textbf{Notes de traduction :}
\begin{itemize}
\item \zh{已经成为} (yǐjīng chéngwéi) $\rightarrow$ "est déjà devenue" (changement d'état \zh{了} implicite ou aspect accompli).
\item \zh{容易} (róngyì) $\rightarrow$ "ont tendance à / risquent de / il est facile qu'ils".
\item \zh{走上错误的道路} $\rightarrow$ "s'engager sur la mauvaise voie / prendre le mauvais chemin" (métaphore).
\item \zh{建议} (jiànyì) $\rightarrow$ "suggèrent / recommandent".
\item \zh{花时间} (huā shíjiān) $\rightarrow$ "passer du temps".
\item \zh{法制教育} (fǎzhì jiàoyù) $\rightarrow$ "éducation juridique / éducation à l'État de droit".
\item \zh{创造机会} (chuàngzào jīhuì) $\rightarrow$ "créer des opportunités" (plus précis que \zh{给...机会}).
\item \zh{全社会} (quán shèhuì) $\rightarrow$ "toute la société / l'ensemble de la société".
\item \zh{共同努力} (gòngtóng nǔlì) $\rightarrow$ "faire des efforts ensemble / conjuguer ses efforts".
\item \zh{能够得到解决} (nénggòu dédào jiějué) $\rightarrow$ "pourra être résolu / pourra trouver une solution" (passif \zh{被} omis, structure \zh{得到} + V).
\end{itemize}
\end{corrige}

\coursec{Production guidée et grille d'évaluation}

Après avoir travaillé la \cle{version} (chinois → français) et le \cle{thème} (français → chinois), il est temps de passer à l'étape ultime de cette leçon : la \cle{production écrite autonome}. Vous allez maintenant rédiger votre propre texte argumentatif sur la délinquance juvénile, en mobilisant tout ce que vous avez appris : le vocabulaire du thème, les connecteurs, la structure en cinq étapes et l'écriture correcte des caractères.

\courssub{Le sujet et le plan imposé}

\begin{definition}[Le sujet de production]
\textbf{Sujet :} « La délinquance juvénile : causes et solutions »\\
\zh{青少年犯罪：原因和解决办法}\\
\emph{Qīngshàonián fànzuì : yuányīn hé jiějué bànfǎ.}\\
« La délinquance juvénile : causes et solutions. »\\
Vous devez rédiger un texte de \cle{5 à 8 phrases} en chinois, en respectant exactement le plan imposé ci-dessous.
\end{definition}

\begin{propriete}[Le plan imposé — à respecter à la lettre]
\begin{enumerate}
\item \textbf{Phrase 1 — Le problème :} présenter le fait (la délinquance juvénile est un problème grave). Utiliser \zh{青少年犯罪}, \zh{严重}, \zh{问题}.
\item \textbf{Phrase 2 — La cause :} une phrase avec la structure \zh{因为……所以……} (\emph{yīnwèi... suǒyǐ...}, « parce que... donc... »). Utiliser \zh{缺少}.
\item \textbf{Phrases 3 à 5 — Les solutions :} trois solutions introduites par \zh{首先} (\emph{shǒuxiān}, « d'abord »), \zh{其次} (\emph{qícì}, « ensuite »), \zh{最后} (\emph{zuìhòu}, « enfin »). Utiliser \zh{应该} dans chaque phrase.
\item \textbf{Phrase 6 — La conclusion :} commencer par \zh{我相信} (\emph{wǒ xiāngxìn}, « je crois que »). Utiliser \zh{解决} et la structure \zh{只要……就/会……}.
\end{enumerate}
\textbf{Banque de mots obligatoires} (à employer tous dans le texte) : \zh{青少年} (jeunes, adolescents), \zh{犯罪} (commettre un crime), \zh{严重} (grave), \zh{缺少} (manquer de), \zh{应该} (devoir), \zh{解决} (résoudre).
\end{propriete}

\begin{attention}[Erreurs qui coûtent cher à l'examen]
\begin{itemize}
\item \textbf{Caractères mal écrits :} ne confondez pas \zh{犯} (\emph{fàn}, commettre — clé « chien » 犭) et \zh{范} (\emph{fàn}, modèle — clé « herbe » 艹). On écrit \zh{犯罪}, jamais \zh{范罪}. Attention aussi à \zh{育} (éducation, clé 月) vs \zh{肓} (aveugle, clé 月) ; \zh{供} (fournir, clé 亻) vs \zh{工} (travail).
\item \textbf{Oubli de \zh{是} :} en chinois, on dit \zh{这个问题很严重} (ce problème est très grave) sans \zh{是} devant un adjectif prédicatif, mais on dit \zh{青少年犯罪是一个严重的问题} avec \zh{是} devant un groupe nominal. Ne mélangez pas les deux structures.
\item \textbf{Connecteurs absents :} un texte sans \zh{首先 / 其次 / 最后} perd jusqu'à 5 points. Les connecteurs sont la colonne vertébrale du texte argumentatif.
\item \textbf{Phrases calquées du français :} ne traduisez pas mot à mot « il y a beaucoup de jeunes qui... » par une structure française (\zh{有很多青少年...}). Pensez en chinois : Sujet + Verbe + Complément. Ex: \zh{很多青少年缺少关心。}
\item \textbf{Classificateur manquant :} \zh{一个问题}, \zh{一些青少年}, \zh{三个解决办法}. Ne dites pas \zh{是严重问题}.
\end{itemize}
\end{attention}

\courssub{Observation : un texte modèle complet}

Avant de rédiger, observez ce texte modèle qui respecte parfaitement le plan imposé. Repérez les six mots obligatoires et les connecteurs.

\begin{textebox}[Texte modèle : 青少年犯罪]
现在，青少年犯罪是一个严重的问题。\\
\emph{Xiànzài, qīngshàonián fànzuì shì yí ge hěn yánzhòng de wèntí.}\\
« Aujourd'hui, la délinquance juvénile est un problème très grave. »\\[4pt]
因为一些青少年缺少父母的爱和关心，所以他们走上了错误的道路。\\
\emph{Yīnwèi yìxiē qīngshàonián quēshǎo fùmǔ de ài hé guānxīn, suǒyǐ tāmen zǒushàng le cuòwù de dàolù.}\\
« Parce que certains jeunes manquent d'amour et d'attention de leurs parents, ils prennent un mauvais chemin. »\\[4pt]
我们应该怎么解决这个问题呢？首先，父母应该多陪伴孩子，关心他们的生活。\\
\emph{Wǒmen yīnggāi zěnme jiějué zhège wèntí ne? Shǒuxiān, fùmǔ yīnggāi duō péibàn háizi, guānxīn tāmen de shēnghuó.}\\
« Comment devrions-nous résoudre ce problème ? D'abord, les parents devraient passer plus de temps avec leurs enfants et s'intéresser à leur vie. »\\[4pt]
其次，学校应该教育学生遵守法律。\\
\emph{Qícì, xuéxiào yīnggāi jiàoyù xuésheng zūnshǒu fǎlǜ.}\\
« Ensuite, l'école devrait éduquer les élèves à respecter la loi. »\\[4pt]
最后，社会应该给青少年提供更多的就业机会。\\
\emph{Zuìhòu, shèhuì yīnggāi gěi qīngshàonián tígōng gèng duō de jiùyè jīhuì.}\\
« Enfin, la société devrait offrir plus d'opportunités d'emploi aux jeunes. »\\[4pt]
我相信，只要大家一起努力，这个问题一定会解决。\\
\emph{Wǒ xiāngxìn, zhǐyào dàjiā yìqǐ nǔlì, zhège wèntí yídìng huì jiějué.}\\
« Je crois que si tout le monde travaille ensemble, ce problème sera sûrement résolu. »
\end{textebox}

\textbf{Questions de compréhension :}
\begin{enumerate}
\item Combien de solutions le texte propose-t-il ? Quels connecteurs les introduisent ?
\item Retrouvez les six mots obligatoires dans le texte et soulignez-les (mentalement ou sur papier).
\item Quelle structure exprime la cause ? Recopiez la phrase concernée en caractères.
\item Identifiez la structure grammaticale de la conclusion (\zh{只要...会...}).
\end{enumerate}

\begin{corrige}[Questions de compréhension]
\begin{enumerate}
\item Trois solutions, introduites par \zh{首先} (d'abord), \zh{其次} (ensuite) et \zh{最后} (enfin).
\item Les six mots obligatoires : \zh{青少年} (phrase 1), \zh{犯罪} (phrase 1), \zh{严重} (phrase 1), \zh{缺少} (phrase 2), \zh{应该} (phrases 3, 4, 5), \zh{解决} (phrases 3 et 6).
\item La cause est exprimée par \zh{因为……所以……} : \zh{因为一些青少年缺少父母的爱和关心，所以他们走上了错误的道路。}
\item Structure conditionnelle : \zh{只要……就/会……} (\emph{zhǐyào... jiù/huì...}, « pourvu que... alors... »).
\end{enumerate}
\end{corrige}

\courssub{Méthode : réussir sa production en 20 minutes}

\begin{methode}[La méthode en 5 étapes pour la rédaction individuelle]
\begin{enumerate}
\item \textbf{Lire le sujet et le plan (2 min) :} surlignez les six mots obligatoires et les trois connecteurs. Votre texte doit avoir exactement la structure imposée (6 phrases types).
\item \textbf{Préparer un brouillon d'idées (3 min) :} notez en français ou en chinois une cause et trois solutions simples (famille, école, société — le trio le plus facile et attendu). Ex: Cause = parents occupés/travail. Sol 1 = parents : temps/amour. Sol 2 = école : loi/moral. Sol 3 = société : emploi/centres de loisirs.
\item \textbf{Rédiger phrase par phrase (10 min) :} écrivez d'abord les phrases faciles (problème, conclusion), puis la cause avec \zh{因为……所以……}, puis les trois solutions avec \zh{首先 / 其次 / 最后}. Utilisez \textbf{au moins 3 connecteurs} (les 3 obligatoires + \zh{因为/所以}).
\item \textbf{Vérifier chaque caractère (3 min) :} relisez en comptant les mots obligatoires un par un. Vérifiez \zh{犯} (et non 范), \zh{严重}, \zh{缺少}, \zh{解决}, \zh{育}, \zh{供}. Vérifiez l'ordre des mots : Sujet + Verbe + Complément. Vérifiez la présence de \zh{是} devant nom, absence devant adjectif.
\item \textbf{Recopier proprement (2 min) :} une présentation soignée rapporte 2 points. Caractères bien formés, ponctuation chinoise pleine chasse （。，？；）， pas d'espaces entre caractères.
\end{enumerate}
\end{methode}

\begin{exemplebox}[Conclusion modèle à réutiliser et adapter]
\zh{我相信，只要大家一起努力，这个问题一定会解决。}\\
\emph{Wǒ xiāngxìn, zhǐyào dàjiā yìqǐ nǔlì, zhège wèntí yídìng huì jiějué.}\\[4pt]
\textbf{Variantes selon vos solutions :}
\begin{itemize}
\item \zh{我相信，只要家庭和学校一起努力，青少年犯罪一定会减少。} (Si famille et école font des efforts, la délinquance diminuera.)
\item \zh{我相信，只要社会提供更多机会，青少年就不会走上犯罪道路。} (Si la société offre plus d'opportunités, les jeunes n'iront pas sur la voie du crime.)
\end{itemize}
\textbf{Structure clé :} \zh{只要} (Condition) + Sujet + \zh{就/会} (Conséquence). \zh{就} pour action immédiate/logique, \zh{会} pour possibilité future.
\end{exemplebox}

\courssub{La grille d'évaluation sur 20 points}

\begin{propriete}[Barème type examen — Première A]
La moitié des points (6/20) porte sur la \cle{grammaire et l'ordre des mots}. Conséquence stratégique : \textbf{un texte simple mais correct vaut mieux qu'un texte ambitieux fautif}. N'écrivez que des structures que vous maîtrisez : mieux vaut trois solutions courtes et justes qu'une phrase longue pleine d'erreurs.
\end{propriete}

\begin{center}
\begin{tabularx}{\linewidth}{|X|X|X|}\hline
\rowcolor{popblueL} \textbf{Critère} & \textbf{Ce qui est vérifié} & \textbf{Points} \\ \hline
Exactitude des caractères & Traits corrects, pas de confusion 犯/范/育/供, caractères du thème bien écrits & 4 \\ \hline
Grammaire et ordre des mots & Sujet + Verbe + Complément, 因为...所以..., 应该 + V, place compléments, 是/非 是 & 6 \\ \hline
Connecteurs et structure & Plan respecté (6 phrases), 首先/其次/最后, 我相信 + 只要...会... en conclusion & 5 \\ \hline
Richesse du vocabulaire & Les 6 mots obligatoires + mots du thème (法律, 社会, 机会, 陪伴, 遵守...) & 3 \\ \hline
Présentation & Copie propre, ponctuation chinoise (。，？；), lisibilité, pas d'espaces & 2 \\ \hline
\textbf{Total} & & \textbf{20} \\ \hline
\end{tabularx}
\end{center}

\begin{popfigure}
\begin{tikzpicture}[
  crit/.style={rectangle, rounded corners=3pt, draw=popblue, fill=popblueL, align=center, minimum height=0.9cm, text width=5.2cm, font=\small},
  pts/.style={rectangle, rounded corners=3pt, draw=poporange, fill=poporangeL, align=center, minimum height=0.9cm, text width=1.6cm, font=\small\bfseries},
  head/.style={rectangle, rounded corners=3pt, draw=popdark, fill=popdark, text=white, align=center, minimum height=0.9cm, font=\small\bfseries},
  total/.style={rectangle, rounded corners=3pt, draw=popgreen, fill=popgreenL, align=center, minimum height=0.9cm, text width=5.2cm, font=\small\bfseries},
  totpts/.style={rectangle, rounded corners=3pt, draw=popgreen, fill=popgreenL, align=center, minimum height=0.9cm, text width=1.6cm, font=\small\bfseries}
]
\node[head, text width=5.2cm] (h1) at (0,0) {Critère};
\node[head, text width=1.6cm] (h2) at (3.6,0) {Points};
\node[crit] (c1) at (0,-1.1) {Exactitude des caractères};
\node[pts] (p1) at (3.6,-1.1) {4};
\node[crit] (c2) at (0,-2.2) {Grammaire et ordre des mots};
\node[pts] (p2) at (3.6,-2.2) {6};
\node[crit] (c3) at (0,-3.3) {Connecteurs et structure};
\node[pts] (p3) at (3.6,-3.3) {5};
\node[crit] (c4) at (0,-4.4) {Richesse du vocabulaire};
\node[pts] (p4) at (3.6,-4.4) {3};
\node[crit] (c5) at (0,-5.5) {Présentation};
\node[pts] (p5) at (3.6,-5.5) {2};
\node[total] (c6) at (0,-6.6) {Total};
\node[totpts] (p6) at (3.6,-6.6) {20};
\end{tikzpicture}
\legende{Grille d'évaluation de la production écrite : la grammaire pèse le plus lourd (6 points sur 20).}
\end{popfigure}

\courssub{Carte mentale de révision}

Avant de vous lancer, mémorisez cette carte mentale : elle résume tout ce dont vous avez besoin pour réussir.

\begin{popfigure}
\begin{tikzpicture}[
  root/.style={ellipse, draw=popdark, fill=popdark, text=white, align=center, font=\small\bfseries, minimum width=3.4cm, minimum height=1.2cm},
  br/.style={rectangle, rounded corners=4pt, draw=popblue, fill=popblueL, align=center, font=\small, text width=3.6cm},
  leaf/.style={rectangle, rounded corners=3pt, draw=poppurple, fill=poppurpleL, align=center, font=\footnotesize, text width=3.4cm}
]
\node[root, align=center] (r) at (0,0){青少年犯罪\\ Production écrite};
\node[br] (plan) at (-5.2,2.2) {\textbf{Plan imposé (6 phrases)}};
\node[leaf, align=center] at (-5.2,3.6){1. Problème (严重问题) \\ 2. Cause (因为...所以...缺少) \\ 3. Sol 1 (首先...应该) \\ 4. Sol 2 (其次...应该) \\ 5. Sol 3 (最后...应该) \\ 6. Conclusion (我相信...只要...会)};
\node[br] (conn) at (5.2,2.2) {\textbf{Connecteurs obligatoires}};
\node[leaf] at (5.2,3.6) {因为 / 所以 / 首先 / 其次 / 最后 / 我相信 / 只要 / 会};
\node[br] (vocab) at (-5.2,-2.2) {\textbf{Mots obligatoires (6)}};
\node[leaf] at (-5.2,-3.6) {青少年 / 犯罪 / 严重 / 缺少 / 应该 / 解决};
\node[br] (check) at (5.2,-2.2) {\textbf{Points de vigilance}};
\node[leaf] at (5.2,-3.6) {犯 ≠ 范 / 育 ≠ 肓 / 供 ≠ 工 / S+V+C / 是+Nom / Ø+Adj / Ponctuation 。，；？};
\draw[-{Stealth}, thick, popblue] (r) -- (plan);
\draw[-{Stealth}, thick, popblue] (r) -- (conn);
\draw[-{Stealth}, thick, popblue] (r) -- (vocab);
\draw[-{Stealth}, thick, popblue] (r) -- (check);
\end{tikzpicture}
\legende{Carte mentale de révision : plan, connecteurs, vocabulaire obligatoire et points de vigilance.}
\end{popfigure}

\courssub{Correction collective d'un texte d'élève}

Voici le texte d'un élève de Première A (anonymisé), tel qu'il a été rendu. Il contient \textbf{5 erreurs} majeures. Trouvez-les avant de lire la correction.

\begin{exoresolu}[Correction du texte d'un élève]
\textbf{Texte rendu :}\\
\zh{青少年范罪是很严重的问题。因为一些孩子缺少父母的爱，他们学坏。首先，父母应该关心孩子。其次，学校应该教育学生。最后，社会帮助青少年。我相信这个问题会解决。}\\[4pt]
\textbf{Consigne :} relevez les 5 erreurs et corrigez-les.
\tcblower
\textbf{Solution détaillée :}
\begin{enumerate}
\item \zh{范罪} $\rightarrow$ \zh{犯罪} : confusion de caractères (\zh{范} = modèle ; \zh{犯} = commettre). Perte de points en « exactitude des caractères ».
\item \zh{是很严重的问题} : structure incomplète, il faut un classificateur : \zh{是一个很严重的问题} (ajouter \zh{个}).
\item \zh{因为一些孩子缺少父母的爱，他们学坏} : la structure \zh{因为...所以...} exige \zh{所以} dans la deuxième clause. De plus, \zh{学坏} (apprendre le mal/devenir mauvais) nécessite l'aspect accompli \zh{了} pour indiquer le résultat : \zh{所以他们学坏了。}
\item \zh{社会帮助青少年} : il manque le modal \zh{应该} (obligatoire dans la banque de mots) : \zh{社会应该帮助青少年。} (Ou mieux : \zh{社会应该提供机会} pour utiliser le vocabulaire du thème).
\item Mot obligatoire \zh{解决} absent du corps des solutions (phrases 3-5) ; il n'apparaît qu'en conclusion (\zh{会解决}). Le plan exige son utilisation active dans les solutions. Mieux : phrase 3 \zh{首先，我们应该一起解决这个问题。} ou phrase 5 \zh{最后，社会应该解决就业问题。}
\end{enumerate}
\textbf{Note estimée :} Caractères 3/4 (erreur 范), Grammaire 4/6 (manque 所以, 了, 一个, 应该), Connecteurs 5/5, Vocabulaire 2/3 (解决 mal placé), Présentation 2/2 $\rightarrow$ \textbf{16/20}.
\end{exoresolu}

\courssub{Auto-évaluation et évaluation croisée}

\begin{experience}[Atelier d'évaluation entre pairs (15 minutes)]
\begin{enumerate}
\item \textbf{Rédaction individuelle (20 min) :} chaque élève rédige son texte de 5 à 8 phrases en silence, sans dictionnaire, en suivant la méthode en 5 étapes.
\item \textbf{Auto-évaluation (3 min) :} relisez votre texte avec la grille sous les yeux. Cochez : ai-je mes 6 mots obligatoires ? mes 3 connecteurs d'énumération ? ma conclusion avec \zh{我相信} + \zh{只要...会...} ? Attribuez-vous une note honnête sur 20.
\item \textbf{Évaluation croisée (7 min) :} échangez votre copie avec votre voisin. Corrigez son texte avec un stylo d'une autre couleur : entourez les caractères douteux, soulignez les connecteurs, notez sur 20 avec la grille. Rédigez un commentaire bienveillant : un point fort + un conseil précis.
\item \textbf{Restitution (5 min) :} rendez la copie à son auteur. Comparez l'auto-évaluation et la note du pair : l'écart est-il justifié ? Le professeur corrige ensuite collectivement deux copies au tableau (une bonne, une perfectible).
\end{enumerate}
\end{experience}

\begin{exercice}[À vous d'écrire — Production finale]
Rédigez maintenant votre propre texte « La délinquance juvénile : causes et solutions » (5 à 8 phrases). Contraintes : plan imposé (6 phrases types), 6 mots obligatoires, au moins 4 connecteurs (\zh{因为/所以, 首先, 其次, 最后, 我相信/只要}), caractères soignés, ponctuation chinoise. Vous pouvez évoquer le contexte camerounais avec des mots simples : \zh{在喀麦隆} (\emph{zài Kāmàilóng}, « au Cameroun »), \zh{城市} (\emph{chéngshì}, « ville »), \zh{雅温得} (\emph{Yǎwēndé}), \zh{杜阿拉} (\emph{Dùālá}).
\end{exercice}

\begin{aretenir}[L'essentiel de la production guidée — Leçon 8]
\begin{itemize}
\item Le \cle{plan imposé} est un contrat : 1 problème (\zh{是...问题}), 1 cause (\zh{因为...所以...缺少}), 3 solutions (\zh{首先/其次/最后...应该}), 1 conclusion (\zh{我相信...只要...会...解决}).
\item Les 6 mots obligatoires (\zh{青少年}, \zh{犯罪}, \zh{严重}, \zh{缺少}, \zh{应该}, \zh{解决}) doivent tous apparaître \textbf{au moins une fois}.
\item La grammaire pèse 6 points sur 20 : privilégiez des phrases simples, courtes et correctes (S + V + C). Évitez les subordonnées relatives complexes (\zh{的} phrases longues) si vous n'êtes pas sûr.
\item Relisez systématiquement : \textbf{caractères} (犯≠范, 育≠肓, 供≠工), \textbf{ordre des mots}, \textbf{connecteurs}, \textbf{classificateurs} (\zh{个}, \zh{些}), \textbf{ponctuation chinoise}.
\item Contexte camerounais bienvenu : \zh{在喀麦隆的城市}... mais restez sur du vocabulaire maîtrisé (HSK 1-3).
\end{itemize}
\end{aretenir}

\newpage

\coursec{Activité d'intégration}

\courssub{Objectifs de l'activité}

À la fin de cette activité, tu seras capable de :
\begin{itemize}
\item rédiger un court article argumentatif en chinois (6 à 8 phrases) sur le thème de la délinquance juvénile ;
\item utiliser correctement le lexique du thème : \zh{青少年} (adolescent), \zh{犯罪} (commettre un crime), \zh{原因} (cause), \zh{办法} (solution) ;
\item employer au moins trois connecteurs logiques : \zh{因为……所以……}, \zh{首先 / 其次 / 最后}, \zh{我相信} ;
\item écrire sans faute les caractères clés de la leçon (\zh{犯}, \zh{罪}, \zh{原}, \zh{因}, \zh{办}, \zh{法}) ;
\item corriger l'article d'un camarade à l'aide d'une grille d'évaluation simple.
\end{itemize}

\courssub{Situation}

Le journal chinois de ton lycée, \zh{《我们的声音》} (\emph{Wǒmen de shēngyīn}, « Notre voix »), prépare un numéro spécial intitulé \zh{《青年与社会》} (\emph{Qīngnián yǔ shèhuì}, « Jeunesse et société »). La rédaction, installée au CDI du lycée de Yaoundé, lance un appel à articles. Voici l'annonce publiée au tableau d'affichage :

\begin{textebox}[Appel à articles du journal du lycée]
征稿启事 \\
\emph{Zhēnggǎo qǐshì.} \\
« Appel à articles » \\[0.4em]
同学们好！\\
\emph{Tóngxuémen hǎo!} \\
« Chers camarades, bonjour ! » \\[0.4em]
我们学校的报纸要出一期《青年与社会》。\\
\emph{Wǒmen xuéxiào de bàozhǐ yào chū yì qī 《Qīngnián yǔ shèhuì》.} \\
« Le journal de notre école va publier un numéro « Jeunesse et société ». » \\[0.4em]
请你写一篇文章，题目是《青少年犯罪：原因和办法》。\\
\emph{Qǐng nǐ xiě yì piān wénzhāng, tímù shì 《Qīngshàonián fànzuì: yuányīn hé bànfǎ》.} \\
« Écris un article intitulé « La délinquance juvénile : causes et solutions ». » \\[0.4em]
文章要有六到八句话。\\
\emph{Wénzhāng yào yǒu liù dào bā jù huà.} \\
« L'article doit contenir de six à huit phrases. » \\[0.4em]
最好的文章会贴在教室的墙上！\\
\emph{Zuì hǎo de wénzhāng huì tiē zài jiàoshì de qiáng shang!} \\
« Les meilleurs articles seront affichés au mur de la classe ! »
\end{textebox}

Tu décides de participer. Ton article devra présenter le problème, expliquer une ou deux causes, proposer des solutions et se terminer par une conclusion personnelle.

\courssub{Consignes}

\begin{enumerate}
\item \textbf{Relecture du modèle.} Relis le texte modèle de la leçon et souligne : le titre, la phrase qui présente le problème, les deux causes, les solutions et la conclusion. Combien de phrases compte-t-il ?
\item \textbf{Préparation du lexique.} Recopie sur ton brouillon six mots du thème que tu veux utiliser (par exemple \zh{青少年}, \zh{问题}, \zh{原因}, \zh{家庭}, \zh{学校}, \zh{帮助}) et vérifie leur écriture, trait par trait.
\item \textbf{Plan.} Rédige ton plan en français en quatre points : problème → cause(s) → solution(s) → conclusion. Chaque point correspondra à une ou deux phrases en chinois.
\item \textbf{Rédaction.} Écris ton article de 6 à 8 phrases en respectant les contraintes suivantes :
\begin{itemize}
\item titre obligatoire : \zh{青少年犯罪：原因和办法} ;
\item au moins une phrase avec \zh{因为……所以……} ;
\item les connecteurs \zh{首先}, \zh{其次}, \zh{最后} pour ordonner les idées ;
\item une conclusion commençant par \zh{我相信} ;
\item des phrases courtes (Sujet + Verbe + Complément), sans mot nouveau non vérifié.
\end{itemize}
\item \textbf{Relecture personnelle.} Vérifie : l'ordre des mots, la place de \zh{很} devant les adjectifs, l'écriture des caractères \zh{犯罪} et \zh{原因}, la ponctuation chinoise （， 。 ：）.
\item \textbf{Correction croisée.} Échange ton article avec ton voisin. À l'aide de la grille d'évaluation, note son texte sur 10 et écris un commentaire positif et un conseil d'amélioration en français.
\item \textbf{Lecture à voix haute.} Lis ta version finale à voix haute devant ton binôme : ton voisin vérifie tes tons, surtout \zh{fànzuì} (4\textsuperscript{e} + 4\textsuperscript{e} ton) et \zh{yuányīn} (2\textsuperscript{e} + 1\textsuperscript{er} ton).
\end{enumerate}

\courssub{Ressources à mobiliser}

\textbf{Vocabulaire du thème}

\begin{center}
\begin{tabularx}{\linewidth}{|X|X|X|X|}
\hline
\rowcolor{popblueL} Caractères & Pinyin & Nature & Traduction \\
\hline
青少年 & qīngshàonián & nom & adolescent, jeune \\
\hline
犯罪 & fànzuì & verbe & commettre un crime \\
\hline
问题 & wèntí & nom & problème, question \\
\hline
原因 & yuányīn & nom & cause, raison \\
\hline
办法 & bànfǎ & nom & solution, moyen \\
\hline
家庭 & jiātíng & nom & famille \\
\hline
帮助 & bāngzhù & verbe/nom & aider ; aide \\
\hline
教育 & jiàoyù & nom/verbe & éducation ; éduquer \\
\hline
社会 & shèhuì & nom & société \\
\hline
应该 & yīnggāi & verbe modal & devoir (il faut) \\
\hline
\end{tabularx}
\end{center}

\textbf{Connecteurs et structures}

\begin{center}
\begin{tabularx}{\linewidth}{|X|X|X|}
\hline
\rowcolor{popblueL} Structure & Exemple en caractères + pinyin & Traduction \\
\hline
因为……所以…… & 因为他们没有工作，所以他们去偷东西。\emph{Yīnwèi tāmen méiyǒu gōngzuò, suǒyǐ tāmen qù tōu dōngxi.} & Parce qu'ils n'ont pas de travail, ils volent des choses. \\
\hline
首先 / 其次 / 最后 & 首先，家长要关心孩子。\emph{Shǒuxiān, jiāzhǎng yào guānxīn háizi.} & D'abord, les parents doivent s'occuper de leurs enfants. \\
\hline
应该 + verbe & 学校应该帮助学生。\emph{Xuéxiào yīnggāi bāngzhù xuésheng.} & L'école doit aider les élèves. \\
\hline
我相信 + phrase & 我相信这个问题可以解决。\emph{Wǒ xiāngxìn zhège wèntí kěyǐ jiějué.} & Je crois que ce problème peut être résolu. \\
\hline
越来越 + adjectif & 这个问题越来越严重。\emph{Zhège wèntí yuè lái yuè yánzhòng.} & Ce problème devient de plus en plus grave. \\
\hline
\end{tabularx}
\end{center}

\begin{attention}[Pièges à éviter]
\begin{itemize}
\item N'écris pas \zh{因为} et \zh{所以} dans la même phrase si tu traduis mot à mot depuis le français : en chinois c'est possible et même fréquent, mais n'oublie jamais la virgule chinoise « ， » entre les deux propositions.
\item Attention aux caractères proches : \zh{罪} (zuì, crime — clé 罒 au-dessus de 非) ne se confond pas avec \zh{罚} (fá, punir) ; \zh{原} (yuán, origine) et \zh{愿} (yuàn, souhaiter) sont homophones proches mais différents.
\item L'adjectif seul ne fait pas une phrase : écris \zh{这个问题很严重} et non \zh{这个问题严重} dans un article soigné.
\end{itemize}
\end{attention}

\courssub{Proposition de production et corrigé commenté}

\begin{exoresolu}[Article modèle : 青少年犯罪：原因和办法]
Rédige un article de 6 à 8 phrases en suivant le plan problème → cause → solutions → conclusion, avec au moins trois connecteurs.
\tcblower
\textbf{Production modèle} \\[0.3em]
青少年犯罪：原因和办法 \\
\emph{Qīngshàonián fànzuì: yuányīn hé bànfǎ.} \\[0.4em]
现在，青少年犯罪是一个越来越严重的问题。\\
\emph{Xiànzài, qīngshàonián fànzuì shì yí ge yuè lái yuè yánzhòng de wèntí.} \\
« Aujourd'hui, la délinquance juvénile est un problème de plus en plus grave. » \\[0.4em]
首先，有些青少年犯罪是因为他们的家庭有困难。\\
\emph{Shǒuxiān, yǒuxiē qīngshàonián fànzuì shì yīnwèi tāmen de jiātíng yǒu kùnnan.} \\
« D'abord, certains jeunes deviennent délinquants parce que leur famille a des difficultés. » \\[0.4em]
其次，因为他们没有工作，也没有钱，所以他们去偷东西。\\
\emph{Qícì, yīnwèi tāmen méiyǒu gōngzuò, yě méiyǒu qián, suǒyǐ tāmen qù tōu dōngxi.} \\
« Ensuite, parce qu'ils n'ont ni travail ni argent, ils volent des choses. » \\[0.4em]
我们应该怎么帮助他们呢？\\
\emph{Wǒmen yīnggāi zěnme bāngzhù tāmen ne?} \\
« Comment devrions-nous les aider ? » \\[0.4em]
家长要关心孩子，学校要教育学生。\\
\emph{Jiāzhǎng yào guānxīn háizi, xuéxiào yào jiàoyù xuésheng.} \\
« Les parents doivent s'occuper de leurs enfants, l'école doit éduquer les élèves. » \\[0.4em]
最后，社会也应该给年轻人工作的机会。\\
\emph{Zuìhòu, shèhuì yě yīnggāi gěi niánqīngrén gōngzuò de jīhuì.} \\
« Enfin, la société doit aussi donner aux jeunes des chances de travailler. » \\[0.4em]
我相信，只要大家一起努力，这个问题一定可以解决。\\
\emph{Wǒ xiāngxìn, zhǐyào dàjiā yìqǐ nǔlì, zhège wèntí yídìng kěyǐ jiějué.} \\
« Je crois que si tout le monde fait des efforts ensemble, ce problème pourra sûrement être résolu. » \\[0.6em]
\textbf{Commentaire des choix de langue}
\begin{itemize}
\item \textbf{Phrase 1 (problème)} : la structure \zh{越来越 + adjectif} présente un phénomène qui s'aggrave ; \zh{严重} est précédé de la structure correcte avec \zh{的} avant le nom \zh{问题}.
\item \textbf{Phrases 2 et 3 (causes)} : \zh{首先} et \zh{其次} ordonnent les causes ; la phrase 3 illustre la corrélation complète \zh{因为……所以……} avec la virgule chinoise ; \zh{也} évite la répétition de \zh{没有}.
\item \textbf{Phrase 4 (transition)} : une question avec \zh{呢} relance le lecteur et annonce les solutions — procédé simple et naturel en chinois.
\item \textbf{Phrases 5 et 6 (solutions)} : phrases courtes et parallèles (\zh{家长……学校……社会……}), procédé apprécié dans l'écriture chinoise ; le modal \zh{应该} exprime le devoir moral.
\item \textbf{Phrase 7 (conclusion)} : \zh{我相信} ouvre la conclusion personnelle ; \zh{只要……就……} (ici avec \zh{一定可以}) exprime la condition, structure de niveau HSK 3 valorisante.
\end{itemize}
\end{exoresolu}

\textbf{Grille d'évaluation pour la correction croisée (sur 10)}

\begin{center}
\begin{tabularx}{\linewidth}{|X|X|}
\hline
\rowcolor{popblueL} Critère & Points \\
\hline
Respect du plan problème → cause → solutions → conclusion & 2 \\
\hline
Présence d'au moins 3 connecteurs correctement employés & 2 \\
\hline
Exactitude des structures (ordre des mots, 因为……所以……, 应该) & 2 \\
\hline
Écriture correcte des caractères du thème (犯罪, 原因, 办法) & 2 \\
\hline
Ponctuation chinoise et présentation (titre, 6 à 8 phrases) & 2 \\
\hline
\end{tabularx}
\end{center}

\courssub{Bilan de l'activité}

À l'issue de cette activité, tu as produit un texte argumentatif complet en chinois, compétence centrale de la leçon. Tu as mobilisé le lexique de la délinquance juvénile, les connecteurs logiques \zh{因为……所以……}, \zh{首先 / 其次 / 最后} et \zh{我相信}, ainsi que l'écriture soignée des caractères du thème. La correction croisée t'a appris à repérer les erreurs typiques des francophones : ordre des mots, oubli de \zh{很}, confusion entre caractères proches. La lecture à voix haute a enfin consolidé ta prononciation des tons. Les meilleurs articles, affichés au mur de la classe, serviront de modèles pour la prochaine production écrite. Garde ton article dans ton portfolio : il pourra être réutilisé et amélioré lors des révisions du Module 1.

\newpage

\coursec{Méthodes et exercices résolus}

\coursec{Leçon 8 — Expression écrite : la délinquance juvénile}

\courssub{Modèle de texte et compréhension}

\begin{textebox}[Texte modèle : La délinquance juvénile à Douala]
\zh{在杜阿拉，青少年犯罪越来越多。很多孩子因为家里穷，不能上学。他们白天在街上，晚上也不回家。警察和政府都想帮助这些年轻人。} \\
\emph{Zài Dù'ālā, qīngshàonián fànzuì yuè lái yuè duō. Hěn duō háizi yīnwèi jiā lǐ qióng, bù néng shàngxué. Tāmen báitiān zài jiēshang, wǎnshang yě bù huí jiā. Jǐngchá hé zhèngfǔ dōu xiǎng bāngzhù zhèxiē niánqīngrén.} \\
« À Douala, la délinquance juvénile est de plus en plus fréquente. Beaucoup d'enfants, parce que leur famille est pauvre, ne peuvent pas aller à l'école. Le jour, ils sont dans la rue ; le soir, ils ne rentrent pas non plus à la maison. La police et le gouvernement veulent tous aider ces jeunes. »
\end{textebox}

\begin{methode}[Répondre aux questions de compréhension écrite]
\begin{enumerate}
\item \textbf{Lire d'abord les questions}, puis le texte : on sait ainsi quoi chercher (qui, quoi, où, pourquoi, comment).
\item \textbf{Repérer dans le texte les mots de la question} (ou leurs synonymes) et souligner la phrase qui contient la réponse.
\item \textbf{Répondre par une phrase complète en chinois}, en reprenant la structure de la question : si la question est \zh{为什么…？}, la réponse commence par \zh{因为…} ; si \zh{谁…？}, le sujet de la réponse est la personne trouvée.
\item \textbf{Ne jamais recopier tout le texte :} extraire uniquement l'information demandée, en adaptant les pronoms (\zh{他们}, \zh{这些年轻人}).
\end{enumerate}
\end{methode}

\begin{exoresolu}[Compréhension du texte modèle]
Énoncé : Lis le texte ci-dessus et réponds en chinois aux questions suivantes.
\begin{enumerate}
\item \zh{为什么很多孩子不能上学？}
\item \zh{谁想帮助这些年轻人？}
\item \zh{这些年轻人白天在哪里？晚上做什么？}
\end{enumerate}
\tcblower
\textbf{Analyse du texte (segmentée) :}
\begin{itemize}
\item \zh{在杜阿拉} (À Douala) — Lieu.
\item \zh{青少年犯罪越来越多} (la délinquance juvénile augmente) — Situation, structure \zh{越来越 + adjectif}.
\item \zh{很多孩子因为家里穷，不能上学} (Beaucoup d'enfants ne peuvent pas aller à l'école car leur famille est pauvre) — Cause (\zh{因为}) + Conséquence (\zh{不能}).
\item \zh{他们白天在街上，晚上也不回家} (Le jour ils sont dans la rue, le soir ils ne rentrent pas non plus) — Compléments de temps \zh{白天}/\zh{晚上} placés AVANT le verbe ; \zh{也} + \zh{不} pour « non plus ».
\item \zh{警察和政府都想帮助这些年轻人} (Police et gouvernement veulent tous aider ces jeunes) — Sujet coordonné + \zh{都} (tous les deux) + \zh{想} (vouloir) + verbe.
\end{itemize}

\textbf{Question 1 :} \zh{为什么很多孩子不能上学？}
\emph{Repérage :} \zh{因为家里穷，不能上学}.
\textbf{Réponse :} \zh{因为他们家里很穷，所以不能上学。} \emph{Yīnwèi tāmen jiā lǐ hěn qióng, suǒyǐ bù néng shàngxué.} (Structure \zh{因为…所以…} obligatoire en réponse à \zh{为什么}).

\textbf{Question 2 :} \zh{谁想帮助这些年轻人？}
\emph{Repérage :} \zh{警察和政府都想帮助这些年轻人}.
\textbf{Réponse :} \zh{警察和政府都想帮助这些年轻人。} \emph{Jǐngchá hé zhèngfǔ dōu xiǎng bāngzhù zhèxiē niánqīngrén.} (\zh{都} insiste sur « tous les deux »).

\textbf{Question 3 :} \zh{这些年轻人白天在哪里？晚上做什么？}
\emph{Repérage :} \zh{他们白天在街上，晚上也不回家}.
\textbf{Réponse :} \zh{他们白天在街上，晚上不回家。} \emph{Tāmen báitiān zài jiēshang, wǎnshang bù huí jiā.} (Réponse synthétique, sans \zh{也} car la question ne porte pas sur l'addition).
\end{exoresolu}

\courssub{Lexique du thème et collocations}

\begin{definition}[Vocabulaire de base — La délinquance juvénile \zh{青少年犯罪}]
Mots à connaître absolument (niveau HSK 2-3) pour lire, écrire et parler du thème.
\end{definition}

\begin{center}
\begin{tabularx}{\linewidth}{|X|X|X|X|}
\hline
\rowcolor{popblueL} \textbf{Caractères} & \textbf{Pinyin} & \textbf{Nature} & \textbf{Traduction} \\ \hline
\zh{青少年} & qīngshàonián & n. & adolescent(s), jeune(s) \\ \hline
\zh{犯罪} & fànzuì & v./n. & commettre un crime / crime, délinquance \\ \hline
\zh{偷} & tōu & v. & voler \\ \hline
\zh{小偷} & xiǎotōu & n. & voleur, pickpocket \\ \hline
\zh{抓} & zhuā & v. & attraper, arrêter (par la police) \\ \hline
\zh{警察} & jǐngchá & n. & police, policier \\ \hline
\zh{政府} & zhèngfǔ & n. & gouvernement \\ \hline
\zh{帮助} & bāngzhù & v. & aider, assistance \\ \hline
\zh{学校} & xuéxiào & n. & école \\ \hline
\zh{工作} & gōngzuò & v./n. & travailler / travail, emploi \\ \hline
\zh{机会} & jīhuì & n. & occasion, opportunité \\ \hline
\zh{穷} & qióng & adj. & pauvre (financièrement) \\ \hline
\zh{富} & fù & adj. & riche \\ \hline
\zh{问题} & wèntí & n. & problème \\ \hline
\zh{原因} & yuányīn & n. & cause, raison \\ \hline
\zh{解决} & jiějué & v. & résoudre (un problème) \\ \hline
\zh{教育} & jiàoyù & v./n. & éduquer / éducation \\ \hline
\zh{家庭} & jiātíng & n. & famille, foyer \\ \hline
\zh{社会} & shèhuì & n. & société \\ \hline
\zh{环境} & huánjìng & n. & environnement \\ \hline
\zh{迷路} & mílù & v. & se perdre (fig. : perdre le droit chemin) \\ \hline
\end{tabularx}
\end{center}

\begin{propriete}[Collocations fréquentes (搭配) — À mémoriser par blocs]
\begin{itemize}
\item \zh{青少年犯罪} (qīngshàonián fànzuì) — délinquance juvénile (sujet + verbe/nom).
\item \zh{偷东西} (tōu dōngxi) — voler des choses (verbe + objet générique).
\item \zh{抓小偷} / \zh{抓住小偷} (zhuā xiǎotōu / zhuāzhù xiǎotōu) — arrêter le voleur (résultat \zh{住}).
\item \zh{帮助年轻人} (bāngzhù niánqīngrén) — aider les jeunes.
\item \zh{上学} (shàngxué) — aller à l'école (verbe + objet locatif) ; négation \zh{不上学}.
\item \zh{找工作} (zhǎo gōngzuò) — chercher du travail.
\item \zh{因为…所以…} (yīnwèi… suǒyǐ…) — parce que… donc… (corrélation inséparable).
\item \zh{越来越 + adj.} (yuè lái yuè + adj.) — de plus en plus… (ex: \zh{越来越多} de plus en plus nombreux).
\item \zh{不但…而且…} (búdàn… érqiě…) — non seulement… mais encore… (niveau supérieur, à reconnaître).
\end{itemize}
\end{propriete}

\begin{exemplebox}[Exemples de collocations en phrases]
\begin{itemize}
\item \zh{政府应该解决青少年犯罪的问题。} \emph{Zhèngfǔ yīnggāi jiějué qīngshàonián fànzuì de wèntí.} (Le gouvernement doit résoudre le problème de la délinquance juvénile.)
\item \zh{因为家庭环境不好，所以他偷东西。} \emph{Yīnwèi jiātíng huánjìng bù hǎo, suǒyǐ tā tōu dōngxi.} (Parce que l'environnement familial est mauvais, il vole.)
\item \zh{警察抓住了两个小偷。} \emph{Jǐngchá zhuāzhù le liǎng ge xiǎotōu.} (La police a attrapé deux voleurs.)
\end{itemize}
\end{exemplebox}

\courssub{Écriture correcte des caractères du thème}

\begin{methode}[Maîtriser les caractères clés : ordre des traits, radicaux, pièges]
\begin{enumerate}
\item \textbf{Identifier le radical (clé de dictionnaire)} qui donne l'indice sémantique.
\item \textbf{Respecter l'ordre des traits universel :} haut $\rightarrow$ bas ; gauche $\rightarrow$ droite ; horizontal $\rightarrow$ vertical ; extérieur $\rightarrow$ intérieur $\rightarrow$ fermer la boîte ; trait central $\rightarrow$ côtés.
\item \textbf{Distinguer les caractères proches (homophones ou graphies voisines)} : c'est la source n°1 d'erreurs au Bac.
\item \textbf{Mémoriser par mots composés (词)}, jamais isolément : le sens et la prononciation se fixent dans le contexte.
\item \textbf{S'entraîner :} écrire chaque caractère 5 fois de suite en prononçant le pinyin (ton inclus) à voix haute.
\end{enumerate}
\end{methode}

\begin{center}
\begin{tabularx}{\linewidth}{|X|X|X|X|X|}
\hline
\rowcolor{popblueL} \textbf{Caractère} & \textbf{Pinyin} & \textbf{Radical (Clé)} & \textbf{Nb traits} & \textbf{Caractère proche / Piège} \\ \hline
\zh{犯} & fàn (4) & \zh{犭} (quǎn, chien) & 5 & \zh{饭} (fàn, riz/repas) — clé \zh{饣} (nourriture) \\ \hline
\zh{罪} & zuì (4) & \zh{网} (wǎng, filet) & 13 & \zh{罚} (fá, punir) — clé \zh{罒} (filet) + \zh{刀} \\ \hline
\zh{偷} & tōu (1) & \zh{亻} (rén, homme) & 11 & \zh{输} (shū, perdre/transporter) — clé \zh{车} \\ \hline
\zh{警} & jǐng (3) & \zh{讠} (yán, parole) & 19 & \zh{敬} (jìng, respecter) — même phonétique, sens diff. \\ \hline
\zh{察} & chá (2) & \zh{宀} (mián, toit) & 14 & \zh{祭} (jì, sacrifier) — structure sup. diff. \\ \hline
\zh{帮} & bāng (1) & \zh{巾} (jīn, tissu) & 12 & \zh{邦} (bāng, état/pays) — phonétique commune \\ \hline
\zh{青} & qīng (1) & \zh{青} (qīng, bleu/vert) & 8 & \zh{静} (jìng, calme) — \zh{青} en haut à droite \\ \hline
\zh{少} & shǎo (3) & \zh{小} (xiǎo, petit) & 4 & \zh{小} (xiǎo) — trait vertical central ne dépasse pas \\ \hline
\zh{年} & nián (2) & \zh{干} (gān/gàn) & 6 & \zh{秂} (ancienne forme) — ne pas ajouter de trait \\ \hline
\end{tabularx}
\end{center}

\begin{exoresolu}[Quiz d'écriture et de distinction — Corrigé commenté]
Énoncé :
1) Complète : \zh{青少年＿＿罪} (commettre un crime). Choisis : \zh{犯} ou \zh{饭} ?
2) Donne la clé de \zh{偷} et son sens. Écris l'ordre des 2 premiers traits.
3) Écris « police » en caractères + pinyin + traduction. Décompose \zh{警}.
4) Parmi \zh{罚} et \zh{罪}, lequel signifie « crime » ? Justifie par la clé.
\tcblower
\textbf{1)} \zh{青少年\underline{犯}罪} \emph{qīngshàonián fànzuì}. \zh{犯} (commettre) a la clé \zh{犭} (chien/bête $\rightarrow$ acte sauvage/illégal). \zh{饭} (repas) a la clé \zh{饣} (nourriture). \textbf{Erreur éliminatoire} si confusion.

\textbf{2)} \zh{偷} (tōu) : clé \zh{亻} (homme, action humaine). 1er trait : \zh{/} (pie gauche) ; 2e trait : \zh{|} (pie droite). Puis partie droite \zh{俞} de haut en bas.

\textbf{3)} \zh{警察} \emph{jǐngchá}. \zh{警} : clé \zh{讠} (parole/avertir) + phonétique \zh{敬} (respect/alerte) $\rightarrow$ « alerter par la parole » $\rightarrow$ police. \zh{察} : clé \zh{宀} (toit/maison) + \zh{祭} (examiner) $\rightarrow$ « examiner sous le toit » $\rightarrow$ inspecter.

\textbf{4)} \zh{罪} (zuì) = crime/faute. Clé \zh{网} (filet) : être pris dans le filet de la loi. \zh{罚} (fá) = punition. Clé \zh{罒} (filet) + \zh{刀} (sabre) : sanction par le sabre.
\end{exoresolu}

\begin{experience}[Atelier d'écriture en classe — 15 min]
\begin{enumerate}
\item \textbf{Dictée muette :} le professeur donne le pinyin (ex: \emph{qīngshàonián fànzuì}), les élèves écrivent les caractères sur ardoise.
\item \textbf{Chasse aux intrus :} sur une liste de 10 caractères (\zh{犯, 饭, 偷, 输, 警, 敬, 罚, 罪, 帮, 邦}), barrer celui qui n'appartient pas au thème « délinquance/sécurité » et justifier par la clé.
\item \textbf{Course aux radicaux :} par équipes, classer 15 cartes-caractères sous le bon radical affiché au tableau (\zh{亻, 讠, 犭, 宀, 巾, 饣}).
\end{enumerate}
\end{experience}

\courssub{Structure et connecteurs du texte argumentatif}

\begin{propriete}[Structure type d'un paragraphe argumentatif simple (5-8 phrases)]
\begin{center}
\begin{tabularx}{\linewidth}{|X|X|X|}
\hline
\rowcolor{popblueL} \textbf{Partie} & \textbf{Fonction} & \textbf{Connecteurs / Structures clés} \\ \hline
\zh{现状} (Situation) & Constater le fait, poser le problème & \zh{现在} (maintenant), \zh{目前} (actuellement), \zh{是一个大问题} (est un grand problème), \zh{越来越严重} (de plus en plus grave) \\ \hline
\zh{原因} (Causes) & Expliquer pourquoi & \zh{因为}…\zh{所以}… (parce que… donc…), \zh{主要原因是} (la cause principale est), \zh{由于} (à cause de), \zh{首先}…\zh{其次}… (premièr… ensuite…) \\ \hline
\zh{后果/例子} (Conséquences/Exemples) & Illustrer concrètement & \zh{例如} (par exemple), \zh{有些}…\zh{有些}… (certains… d'autres…), \zh{导致} (entraîner), \zh{结果} (résultat) \\ \hline
\zh{对策/建议} (Solutions) & Proposer des remèdes & \zh{应该} (doit), \zh{必须} (il faut impérativement), \zh{建议} (suggérer), \zh{加强教育} (renforcer l'éducation), \zh{提供机会} (offrir des opportunités) \\ \hline
\zh{结论/观点} (Conclusion/Avis) & Clore avec opinion personnelle & \zh{我觉得} (je pense que), \zh{我认为} (j'estime que), \zh{总之} (en résumé), \zh{只有…才…} (seulement si… alors…) \\ \hline
\end{tabularx}
\end{center}
\end{propriete}

\begin{methode}[Rédiger un court texte argumentatif en 5 étapes (Bac blanc)]
\begin{enumerate}
\item \textbf{Analyser le sujet.} Thème : \cle{la délinquance juvénile} \zh{青少年犯罪}. Type : paragraphe d'opinion. Longueur : 5-8 phrases (1ère A).
\item \textbf{Construire un mini-plan (brouillon) :} 1. Situation (\zh{现在…}) $\rightarrow$ 2. Causes (\zh{因为…所以…}) $\rightarrow$ 3. Solution (\zh{应该…}) $\rightarrow$ 4. Mon avis (\zh{我觉得…}).
\item \textbf{Sélectionner 4-5 connecteurs} avant d'écrire (cocher la liste ci-dessus). Ne pas en mettre dans chaque phrase.
\item \textbf{Rédiger phrases courtes :} Sujet + Temps/Lieu + Verbe + Objet. Une idée = une phrase. Vérifier l'ordre \zh{因为}…\zh{所以}… et la place des compléments.
\item \textbf{Relire (check-list) :} Ordre des mots ? Tons notés ? Caractères corrects (clés) ? Classificateurs (\zh{个, 些, 位}) ? Ponctuation chinoise (。，) ? Conclusion présente ?
\end{enumerate}
\end{methode}

\begin{exemplebox}[Modèle de paragraphe complet (6 phrases) — Niveau Bac 1ère A]
\zh{现在，大城市的青少年犯罪是一个严重的问题。因为很多年轻人没有工作，也没有钱，所以他们去偷东西。还有些孩子因为家里穷，不上学，整天在街上。但是我觉得，家庭和学校应该多关心这些孩子。政府也应该给年轻人创造更多的工作机会。所以，全社会要一起努力，解决这个问题。} \\
\emph{Xiànzài, dà chéngshì de qīngshàonián fànzuì shì yí ge yánzhòng de wèntí. Yīnwèi hěn duō niánqīngrén méiyǒu gōngzuò, yě méiyǒu qián, suǒyǐ tāmen qù tōu dōngxi. Hái yǒu xiē háizi yīnwèi jiā lǐ qióng, bù shàngxué, zhěngtiān zài jiēshang. Dànshì wǒ juéde, jiātíng hé xuéxiào yīnggāi duō guānxīn zhèxiē háizi. Zhèngfǔ yě yīnggāi gěi niánqīngrén chuàngzào gèng duō de gōngzuò jīhuì. Suǒyǐ, quán shèhuì yào yìqǐ nǔlì, jiějué zhè ge wèntí.} \\
« Aujourd'hui, la délinquance juvénile dans les grandes villes est un problème grave. Parce que beaucoup de jeunes n'ont ni travail ni argent, ils vont voler. Certains enfants, parce que leur famille est pauvre, ne vont pas à l'école et passent la journée dans la rue. Mais je pense que la famille et l'école devraient davantage se soucier de ces enfants. Le gouvernement devrait aussi créer plus d'opportunités d'emploi pour les jeunes. Donc, toute la société doit faire des efforts ensemble pour résoudre ce problème. »
\end{exemplebox}

\courssub{Thème : traduction français → chinois}

\begin{methode}[Traduire une phrase française en chinois — 4 étapes obligatoires]
\begin{enumerate}
\item \textbf{Analyser la phrase française :} Sujet ? Verbe (temps/négation) ? Compléments (Temps, Lieu, Manière, Objet) ?
\item \textbf{Reclasser selon l'ordre chinois canonique :} \cle{Sujet + Temps + Lieu + Manière + Verbe + Objet}. (Le français met souvent Lieu/Temps après le verbe).
\item \textbf{Traduire les mots-outils grammaticaux :}
\begin{itemize}
\item Négation passé/accompli : \zh{没(有)} + V (ex: \zh{没去}) — \textbf{jamais \zh{不} pour le passé}.
\item Aspect accompli : \zh{了} après V (ou V+Objet si court).
\item « Deux » devant classifieur : \zh{两} (liǎng), jamais \zh{二} (èr).
\item \zh{一些} (certains/des), \zh{那些} (ceux-là), \zh{这些} (ceux-ci).
\end{itemize}
\item \textbf{Vérifier :} Classificateurs ? Tons ? Caractères (clés) ? Ponctuation ? Sens global ?
\end{enumerate}
\end{methode}

\begin{exoresolu}[Thème : 4 phrases à traduire — Corrigé détaillé]
Énoncé : Traduis en chinois (caractères + pinyin).
\begin{enumerate}
\item Certains jeunes ne vont pas à l'école parce qu'ils sont pauvres.
\item Hier, la police a arrêté deux voleurs près du marché.
\item Nous devons aider ces enfants pour qu'ils ne commettent pas de crimes.
\item À Yaoundé, la délinquance juvénile devient de plus en plus grave.
\end{enumerate}
\tcblower
\textbf{1)} \emph{Analyse :} Sujet = certains jeunes (\zh{一些年轻人}) ; Verbe = ne vont pas (habitude/présent $\rightarrow$ \zh{不}) ; Complément cause = parce qu'ils sont pauvres (\zh{因为他们很穷}).
\textbf{Ordre chinois :} Sujet + Cause + Verbe principal. \zh{因为} clause peut précéder ou suivre la principale. Ici on place la cause après pour varier.
\textbf{Réponse :} \zh{一些年轻人不上学，因为他们很穷。} \emph{Yìxiē niánqīngrén bù shàngxué, yīnwèi tāmen hěn qióng.}
\emph{Note :} \zh{很} obligatoire devant adj. prédicatif \zh{穷}. \zh{不} pour habitude présente.

\textbf{2)} \emph{Analyse :} Temps = Hier (\zh{昨天}) ; Sujet = Police (\zh{警察}) ; Lieu = près du marché (\zh{在市场附近} / \zh{在市场旁边}) ; Verbe accompli = a arrêté (\zh{抓住了} / \zh{抓了}) ; Objet = deux voleurs (\zh{两个小偷}).
\textbf{Ordre :} Sujet + Temps + Lieu + Verbe + Objet.
\textbf{Réponse :} \zh{昨天警察在市场旁边抓住了两个小偷。} \emph{Zuótiān jǐngchá zài shìchǎng pángbiān zhuāzhù le liǎng ge xiǎotōu.}
\emph{Points clés :} \zh{昨天} après sujet ; \zh{在…旁边} avant verbe ; \zh{两} + \zh{个} ; \zh{抓住了} (résultat + accompli).

\textbf{3)} \emph{Analyse :} Sujet = Nous (\zh{我们}) ; Modale obligation = devons (\zh{应该} / \zh{必须}) ; Verbe = aider (\zh{帮助}) ; Objet = ces enfants (\zh{这些孩子}) ; But = pour qu'ils ne commettent pas de crimes (\zh{以免他们犯罪} / \zh{不让他们犯罪}).
\textbf{Réponse :} \zh{我们应该帮助这些孩子，以免他们犯罪。} \emph{Wǒmen yīnggāi bāngzhù zhèxiē háizi, yǐmiǎn tāmen fànzuì.}
\emph{Alternative :} \zh{我们必须帮助这些孩子，不让他们犯罪。} (\zh{不让} = ne pas laisser/faire en sorte que non). \zh{那些} si éloigné, \zh{这些} si proche/présent à l'esprit.

\textbf{4)} \emph{Analyse :} Lieu = À Yaoundé (\zh{在雅温得}) ; Sujet = délinquance juvénile (\zh{青少年犯罪}) ; Verbe = devient (\zh{变得} / \zh{成为}) ; Complément = de plus en plus grave (\zh{越来越严重}).
\textbf{Réponse :} \zh{在雅温得，青少年犯罪变得越来越严重。} \emph{Zài Yǎwēndé, qīngshàonián fànzuì biànde yuè lái yuè yánzhòng.}
\emph{Structure :} \zh{变得} + \zh{越来越} + adj. \zh{严重} (grave/sérieux) mieux que \zh{大} pour « problème grave ».
\end{exoresolu}

\courssub{Version : traduction chinois → français}

\begin{methode}[Traduire un texte chinois en français — 4 étapes]
\begin{enumerate}
\item \textbf{Lecture globale} sans écrire : identifier le sujet, le temps général (passé/présent/futur), le ton.
\item \textbf{Découpage en segments syntaxiques} (Sujet / Temps / Lieu / Verbe / Objet / Particules). Traduire segment par segment.
\item \textbf{Rédaction en français naturel :} Articles (le/la/des), conjugaison (passé composé, imparfait, présent), prépositions correctes (\zh{在} $\rightarrow$ à/dans/sur), ordre SVO français. \textbf{Ne jamais calquer mot à mot.}
\item \textbf{Relecture finale :} Grammaire française correcte ? Fidélité au sens original (ni omission, ni ajout) ? Style fluide ?
\end{enumerate}
\end{methode}

\begin{exoresolu}[Version : texte court — Corrigé commenté]
Énoncé : Traduis en français naturel.
\zh{在雅温得，有一些年轻人不学习，也不工作。他们有时候在街上偷东西。政府想帮助他们，所以建了很多新学校。}
\tcblower
\textbf{Segmentation \& Analyse :}
\begin{enumerate}
\item \zh{在雅温得，有一些年轻人不学习，也不工作。} \\
\emph{Zài Yǎwēndé, yǒu yìxiē niánqīngrén bù xuéxí, yě bù gōngzuò.} \\
\zh{在雅温得} (À Yaoundé) — \zh{有} (il y a / existence) — \zh{一些年轻人} (des jeunes / certains jeunes) — \zh{不学习} (n'étudient pas) — \zh{也不工作} (ne travaillent pas non plus ; \zh{也} + \zh{不} = non plus).
\item \zh{他们有时候在街上偷东西。} \\
\emph{Tāmen yǒushíhou zài jiēshang tōu dōngxi.} \\
\zh{他们} (Ils) — \zh{有时候} (parfois / de temps en temps) — \zh{在街上} (dans la rue) — \zh{偷东西} (volent des choses / volent).
\item \zh{政府想帮助他们，所以建了很多新学校。} \\
\emph{Zhèngfǔ xiǎng bāngzhù tāmen, suǒyǐ jiàn le hěn duō xīn xuéxiào.} \\
\zh{政府} (Le gouvernement) — \zh{想} (veut / souhaite) — \zh{帮助他们} (les aider) — \zh{所以} (donc / c'est pourquoi) — \zh{建了} (a construit / a bâti ; \zh{了} = accompli $\rightarrow$ passé composé) — \zh{很多新学校} (beaucoup de nouvelles écoles).
\end{enumerate}

\textbf{Traduction finale proposée :}
« À Yaoundé, il y a des jeunes qui n'étudient pas et ne travaillent pas non plus. Ils volent parfois des choses dans la rue. Le gouvernement veut les aider, c'est pourquoi il a construit beaucoup de nouvelles écoles. »

\textbf{Points de vigilance :}
\begin{itemize}
\item \zh{有} $\neq$ « avoir » ici mais « il y a » (existence).
\item \zh{也不} $\rightarrow$ « ne … pas non plus » (négation additionnelle).
\item \zh{在街上} (complément de lieu) placé AVANT le verbe en chinois, APRÈS en français.
\item \zh{了} après \zh{建} $\rightarrow$ action achevée dans le passé $\rightarrow$ Passé Composé « a construit ». Pas « construit » (participe) ni « construit » (présent).
\end{itemize}
\end{exoresolu}

\courssub{Production guidée, expression orale et grille d'évaluation}

\begin{experience}[Débat oral guidé : « Causes et solutions » — 20 min]
\textbf{Consigne :} En groupes de 4. Chaque groupe tire une carte « Rôle » et prépare 3 arguments en chinois (notes autorisées). Débat modéré par un élève-président.
\begin{itemize}
\item \textbf{Groupe A (Gouvernement) :} \zh{我们应该建更多学校，给年轻人工作机会。} (Construire écoles, créer emplois).
\item \textbf{Groupe B (Parents/Familles) :} \zh{家庭教育最重要，父母要管孩子。} (Éducation familiale primordiale, parents doivent surveiller).
\item \textbf{Groupe C (Police) :} \zh{法律要严厉，抓到小偷要惩罚。} (Loi sévère, punir les voleurs arrêtés).
\item \textbf{Groupe D (Jeunes/ONG) :} \zh{年轻人需要尊重和心理帮助，不只是惩罚。} (Les jeunes ont besoin de respect et aide psy, pas seulement punition).
\end{itemize}
\textbf{Connecteurs imposés au moins 1 fois par intervention :} \zh{因为…所以…}, \zh{但是}, \zh{另外}, \zh{我觉得}.
\textbf{Évaluation orale :} Prononciation (tons), fluidité, vocabulaire du thème, structure logique, interaction.
\end{experience}

\begin{exercice}[Production écrite individuelle — Devoir / Bac blanc (30 min)]
\textbf{Sujet :} « Dans ta ville, la délinquance juvénile augmente. Écris un paragraphe de 6 à 8 phrases en chinois pour expliquer la situation, donner deux causes et proposer une solution. Termine par ton avis personnel. »
\textbf{Contraintes obligatoires :}
\begin{enumerate}
\item Utiliser au moins 4 connecteurs différents parmi : \zh{现在}, \zh{因为…所以…}, \zh{但是}, \zh{还有}, \zh{应该}, \zh{我觉得}, \zh{因此}, \zh{总之}.
\item Employer au moins 6 mots du vocabulaire de la leçon (surligne-les dans ta copie).
\item Respecter l'ordre des mots chinois (Temps/Lieu avant verbe), les classificateurs, la négation correcte.
\item Écrire les caractères soigneusement (clés vérifiées). Pinyin au-dessus si doute.
\end{enumerate}
\end{exercice}

\begin{corrige}[Exercice — Production écrite : Proposition de corrigé type (Niveau Excellence)]
\textbf{Texte proposé :}
\zh{现在，我们城市的青少年犯罪越来越严重。因为一些年轻人找不到工作，家里又很穷，所以他们偷东西。还有些孩子不上学，整天在街上游荡。但是我觉得，光抓小偷是不够的。政府应该建更多职业学校，给年轻人学技术的机会。家庭也应该多关心孩子的教育。总之，全社会一起努力，这个问题才能解决。} \\
\emph{Xiànzài, wǒmen chéngshì de qīngshàonián fànzuì yuè lái yuè yánzhòng. Yīnwèi yìxiē niánqīngrén zhǎo bu dào gōngzuò, jiā lǐ yòu hěn qióng, suǒyǐ tāmen tōu dōngxi. Hái yǒu xiē háizi bù shàngxué, zhěngtiān zài jiēshang yóudàng. Dànshì wǒ juéde, guāng zhuā xiǎotōu shì bù gòu de. Zhèngfǔ yīnggāi jiàn gèng duō zhíyè xuéxiào, gěi niánqīngrén xué jìshù de jīhuì. Jiātíng yě yīnggāi duō guānxīn háizi de jiàoyù. Zǒngzhī, quán shèhuì yìqǐ nǔlì, zhè ge wèntí cái néng jiějué.}

\textbf{Analyse du corrigé (Grille) :}
\begin{itemize}
\item \textbf{Structure :} Situation (\zh{现在…严重}) $\rightarrow$ Cause 1 (\zh{因为…所以…}) $\rightarrow$ Cause 2 (\zh{还有些…}) $\rightarrow$ Pivot (\zh{但是我觉得}) $\rightarrow$ Solution 1 (\zh{政府应该…}) $\rightarrow$ Solution 2 (\zh{家庭也应该…}) $\rightarrow$ Conclusion (\zh{总之…才能…}). \checkmark
\item \textbf{Connecteurs (6) :} \zh{现在}, \zh{因为…所以…}, \zh{还有}, \zh{但是}, \zh{应该} (x2), \zh{总之}, \zh{才能}. \checkmark
\item \textbf{Vocabulaire thème (8) :} \zh{青少年犯罪}, \zh{严重}, \zh{工作}, \zh{穷}, \zh{偷东西}, \zh{上学}, \zh{小偷}, \zh{政府}, \zh{职业学校}, \zh{机会}, \zh{关心}, \zh{教育}, \zh{社会}, \zh{解决}. \checkmark
\item \textbf{Grammaire :} \zh{找不到} (complément de résultat), \zh{又…又…} (à la fois… et…), \zh{游荡} (errer - vocab+), \zh{光…不够} (seulement… ne suffit pas), \zh{才能} (seulement ainsi peut…). Structures correctes. \checkmark
\item \textbf{Caractères :} \zh{犯}, \zh{罪}, \zh{偷}, \zh{警} (dans \zh{警察} implicite), \zh{帮}, \zh{教}, \zh{育} — clés respectées. \checkmark
\end{itemize}
\end{corrige}

\begin{aretenir}[Grille d'auto-évaluation / Évaluation professeur (sur 20 pts)]
\begin{center}
\begin{tabularx}{\linewidth}{|X|c|X|}
\hline
\rowcolor{popblueL} \textbf{Critère} & \textbf{Points} & \textbf{Indicateurs de réussite} \\ \hline
\textbf{Contenu \& Structure} & 5 & Plan respecté (Situation, Causes, Solutions, Avis) ; 6-8 phrases ; idées claires, logiques. \\ \hline
\textbf{Vocabulaire du thème} & 4 & $\ge$ 6 mots de la liste utilisés correctement ; collocations justes (\zh{偷东西}, \zh{青少年犯罪}, \zh{帮助}…). \\ \hline
\textbf{Grammaire \& Syntaxe} & 5 & Ordre Sujet-Temps-Lieu-Verbe-Objet respecté ; \zh{因为…所以…} correct ; Négation (\zh{不}/\zh{没}) juste ; \zh{了} aspectuel juste ; Classificateurs (\zh{个, 些, 所}) justes ; \zh{两} vs \zh{二}. \\ \hline
\textbf{Caractères \& Écriture} & 3 & Caractères clés (\zh{犯, 罪, 偷, 警, 察, 帮, 青, 少, 年}) bien formés, clés respectées, pas de confusion \zh{犯/饭}, \zh{罪/罚}. Lisible. \\ \hline
\textbf{Pinyin \& Tons} & 3 & Pinyin noté sur les mots nouveaux/difficiles ; tons exacts (1,2,3,4, neutre) ; pas de ton manquant. \\ \hline
\textbf{Expression orale (si évalué)} & — & Prononciation claire, tons audibles, débit naturel, utilisation connecteurs, interaction. \\ \hline
\end{tabularx}
\end{center}
\textbf{Seuil de réussite :} 10/20. \textbf{Excellence :} 16/20 (peu d'erreurs, vocabulaire riche, structure fluide, caractères parfaits).
\end{aretenir}

\courssub{Élément socioculturel : Cameroun — Chine}

\begin{savaistu}[Jeunesse, éducation et prévention : regards croisés]
\textbf{Cameroun :} La délinquance juvénile (\zh{青少年犯罪}) est une préoccupation majeure dans les grandes villes (Yaoundé, Douala, Bafoussam). Facteurs souvent cités : pauvreté urbaine, échec scolaire (\zh{辍学} \emph{chuòxué} = décrochage), chômage des diplômés, rupture familiale, influence de la rue. Réponses : centres de rééducation (\zh{少管所} \emph{shǎoguǎnsuǒ}), programmes « École pour tous », ONG d'encadrement, rôle des chefferies traditionnelles et des églises/mosquées dans l'éducation morale.

\textbf{Chine :} Taux de criminalité juvénile historiquement bas mais en hausse récente (vols, violences, cybercriminalité). Facteurs : pression scolaire extrême (\zh{内卷} \emph{nèijuǎn}), familles monoparentales ou parents migrants (\zh{留守儿童} \emph{liúshǒu értóng} = enfants laissés à la campagne), jeux vidéo/addiction écrans. Réponses : Loi sur la protection des mineurs (2021, révisée), tribunaux pour mineurs spécialisés, « Éducation par le travail » (\zh{劳动教育} \emph{láodòng jiàoyù}) rendue obligatoire à l'école, surveillance numérique, camps de redressement (\zh{戒网瘾学校} \emph{jiè wǎngyǐn xuéxiào}).

\textbf{Point commun :} Les deux pays misent sur l'\textbf{éducation} (\zh{教育}) et la \textbf{famille} (\zh{家庭}) comme premiers remparts, et sur l'\textbf{insertion professionnelle} (\zh{就业} \emph{jiùyè}) comme solution durable. La Chine légifère plus lourdement (responsabilité pénale abaissée à 12 ans pour crimes graves) ; le Cameroun privilégie l'approche communautaire et éducative.

\textbf{Pour ton expression orale/écrite :} Tu peux citer l'exemple des \zh{留守儿童} en Chine ou du \zh{déficit d'encadrement} au Cameroun pour illustrer la cause « absence parentale ».
\end{savaistu}

\begin{attention}[Les 5 erreurs qui coûtent des points au Bac — RAPPEL FINAL]
\begin{enumerate}
\item \textbf{Ordre des mots calqué du français.} \cle{Compléments Temps/Lieu AVANT le verbe.} \zh{他昨天在市场偷东西。} $\checkmark$ \quad \zh{他偷东西昨天在市场。} $\times$
\item \textbf{Confusion caractères proches (Clés !).} \zh{犯} (chien \zh{犭} = crime) $\neq$ \zh{饭} (nourriture \zh{饣} = riz). \zh{罪} (filet \zh{网} = crime) $\neq$ \zh{罚} (filet+sabre \zh{罒}+\zh{刀} = punition). \zh{偷} (homme \zh{亻} = voler) $\neq$ \zh{输} (voiture \zh{车} = perdre). \textbf{1 erreur de clé = 0 pt sur le mot.}
\item \textbf{Négation du passé :} \zh{没(有)} + V (\zh{没去} « n'est pas allé »). \zh{不} + V = présent/habitude/futur (\zh{不去} « ne va pas / n'ira pas »). \textbf{Jamais \zh{不} pour un passé accompli.}
\item \textbf{Classificateurs oubliés ou \zh{二} au lieu de \zh{两}.} \zh{两个小偷} $\checkmark$ \quad \zh{二个小偷} $\times$ \quad \zh{一个问题} $\checkmark$ \quad \zh{一问题} $\times$.
\item \textbf{Tons absents ou faux sur le pinyin.} \zh{mǎi} (acheter) $\neq$ \zh{mài} (vendre). \zh{xuéxí} (étudier) $\neq$ \zh{xuěxì} (n'existe pas). \textbf{Relis ton pinyin syllabe par syllabe avant de rendre la copie.}
\end{enumerate}
\end{attention}

\newpage

\coursec{Exercices}

\courssub{Je vérifie mes connaissances}

\begin{exercice}[QCM : le bon choix]
Pour chaque question, choisis la bonne réponse (a, b ou c).
\begin{enumerate}
\item 犯罪 (fànzuì) signifie :
\begin{itemize}
\item a) la loi
\item b) commettre un crime ; le crime
\item c) la prison
\end{itemize}
\item « 缺少 » (quēshǎo) signifie :
\begin{itemize}
\item a) manquer de
\item b) posséder
\item c) partager
\end{itemize}
\item Quelle phrase est correcte ?
\begin{itemize}
\item a) 我们应该关心青少年。
\item b) 我们应该青少年关心。
\item c) 关心我们应该青少年。
\end{itemize}
\item « 解决 » (jiějué) s'emploie avec :
\begin{itemize}
\item a) 问题 (wèntí)
\item b) 饭 (fàn)
\item c) 天气 (tiānqì)
\end{itemize}
\end{enumerate}
\end{exercice}

\begin{exercice}[Vrai ou faux justifié]
Réponds par VRAI ou FAUX, puis justifie ta réponse en une phrase en français.
\begin{enumerate}
\item 青少年 (qīngshàonián) désigne les personnes âgées.
\item 严重 (yánzhòng) est un adjectif qui signifie « grave, sérieux ».
\item Dans la phrase 家长应该教育孩子，le verbe 教育 signifie « punir ».
\item 法律 (fǎlǜ) signifie « la loi, le droit ».
\end{enumerate}
\end{exercice}

\begin{exercice}[Vocabulaire : complète le tableau]
Recopie et complète le tableau suivant avec la nature et la traduction française de chaque mot.
\begin{center}
\begin{tabularx}{\linewidth}{|X|X|X|X|}
\hline
\rowcolor{popblueL} Caractères & Pinyin & Nature & Traduction \\
\hline
社会 & shèhuì & \ldots & \ldots \\
\hline
错误 & cuòwù & \ldots & \ldots \\
\hline
努力 & nǔlì & \ldots & \ldots \\
\hline
帮助 & bāngzhù & \ldots & \ldots \\
\hline
原因 & yuányīn & \ldots & \ldots \\
\hline
\end{tabularx}
\end{center}
\end{exercice}

\begin{exercice}[Pinyin et caractères : associe]
Associe chaque caractère ou mot (1 à 6) à son pinyin (a à f).
\begin{enumerate}
\item 犯
\item 青
\item 法
\item 缺
\item 严
\item 决
\end{enumerate}
\begin{itemize}
\item a) quē
\item b) fǎ
\item c) jué
\item d) fàn
\item e) yán
\item f) qīng
\end{itemize}
\end{exercice}

\courssub{Je m'entraîne}

\begin{exercice}[Thème : traduis en chinois]
Traduis les phrases suivantes en chinois (caractères simplifiés).
\begin{enumerate}
\item Certains jeunes manquent d'éducation familiale.
\item La délinquance juvénile est un problème grave.
\item Les parents et l'école devraient aider ces jeunes.
\item Nous devons respecter la loi.
\end{enumerate}
\end{exercice}

\begin{exercice}[Discrimination de caractères]
Choisis le bon caractère pour compléter chaque phrase : 犯 / 范 ; 青 / 清 ; 法 / 去.
\begin{enumerate}
\item 青少年犯罪是一个严重的问题。(少年 = jeunesse ; choisis entre 青 et 清)
\item 他犯了错误，应该改正。(choisis entre 犯 et 范)
\item 每个人都应该遵守\_\_律。(choisis entre 法 et 去)
\item 水很\_\_，我可以看见鱼。(choisis entre 清 et 青)
\end{enumerate}
\end{exercice}

\begin{exercice}[Texte à trous : mots-outils]
Complète le texte suivant avec les mots de la banque : 因为、所以、首先、其次、最后、应该.
\begin{textebox}[La délinquance juvénile]
\_\_\_\_\_ 一些青少年缺少家庭教育，\_\_\_\_\_ 他们会犯错误。怎么解决这个问题？\_\_\_\_\_，家长应该多关心孩子；\_\_\_\_\_，学校应该加强教育；\_\_\_\_\_，社会也应该帮助他们。
\end{textebox}
Pinyin : \emph{Yīnwèi yìxiē qīngshàonián quēshǎo jiātíng jiàoyù, suǒyǐ tāmen huì fàn cuòwù. Zěnme jiějué zhège wèntí? Shǒuxiān, jiāzhǎng yīnggāi duō guānxīn háizi; qícì, xuéxiào yīnggāi jiāqiáng jiàoyù; zuìhòu, shèhuì yě yīnggāi bāngzhù tāmen.}
\end{exercice}

\begin{exercice}[Remise en ordre]
Remets dans l'ordre les phrases de ce texte argumentatif mélangé, en repérant les connecteurs. Numérote-les de 1 à 5.
\begin{enumerate}
\item[a)] 最后，社会应该给年轻人更多的机会。
\item[b)] 现在，青少年犯罪是一个严重的问题。
\item[c)] 其次，学校应该加强法律和道德教育。
\item[d)] 怎么解决这个问题呢？
\item[e)] 首先，家长应该多关心孩子，多和他们交流。
\end{enumerate}
\end{exercice}

\courssub{Je me prépare au Bac}

\begin{exercice}[Compréhension de texte et questions de langue]
Lis le texte suivant, puis réponds aux questions.
\begin{textebox}[Aider les jeunes de nos quartiers]
在一些城市，青少年犯罪的问题越来越严重。为什么呢？因为有些年轻人缺少家庭教育，也有些年轻人没有工作，生活很困难。\\
为了解决这个问题，首先，家长应该多关心孩子，多和他们交流；其次，学校应该加强法律教育，让学生知道什么是对的，什么是错的；最后，政府和社会应该给年轻人更多的工作机会。只要大家一起努力，这个问题一定能解决。
\end{textebox}
Pinyin : \emph{Zài yìxiē chéngshì, qīngshàonián fànzuì de wèntí yuèláiyuè yánzhòng. Wèishénme ne? Yīnwèi yǒuxiē niánqīngrén quēshǎo jiātíng jiàoyù, yě yǒuxiē niánqīngrén méiyǒu gōngzuò, shēnghuó hěn kùnnan. Wèile jiějué zhège wèntí, shǒuxiān, jiāzhǎng yīnggāi duō guānxīn háizi, duō hé tāmen jiāoliú; qícì, xuéxiào yīnggāi jiāqiáng fǎlǜ jiàoyù, ràng xuésheng zhīdào shénme shì duì de, shénme shì cuò de; zuìhòu, zhèngfǔ hé shèhuì yīnggāi gěi niánqīngrén gèng duō de gōngzuò jīhuì. Zhǐyào dàjiā yìqǐ nǔlì, zhège wèntí yídìng néng jiějué.}

\textbf{Questions de compréhension} (réponds en français) :
\begin{enumerate}
\item Quel problème est évoqué dans ce texte ?
\item Quelles sont les deux causes mentionnées ?
\item Quelles sont les trois solutions proposées ? À qui revient chaque solution ?
\end{enumerate}
\textbf{Questions de langue} :
\begin{enumerate}
\item Relève dans le texte les trois connecteurs qui organisent les solutions et traduis-les.
\item Explique la structure de la phrase 只要大家一起努力，这个问题一定能解决。(rôle de 只要 \ldots 就/一定 \ldots)
\item Trouve dans le texte un adjectif signifiant « grave » et un verbe signifiant « résoudre ».
\end{enumerate}
\end{exercice}

\begin{exercice}[Thème et version type Bac]
\textbf{Thème} (français $\rightarrow$ chinois). Traduis en chinois simplifié : « Parce que certains jeunes manquent d'éducation familiale, ils commettent des erreurs. L'école et les parents devraient les aider. »

\textbf{Version} (chinois $\rightarrow$ français). Traduis en français le texte suivant :
\begin{textebox}[Version]
怎么解决青少年犯罪的问题？首先，家庭应该给孩子更多的爱；其次，学校应该教学生遵守法律；最后，社会应该帮助有困难的年轻人。我相信，只要我们一起努力，明天一定会更好。
\end{textebox}
Pinyin : \emph{Zěnme jiějué qīngshàonián fànzuì de wèntí? Shǒuxiān, jiātíng yīnggāi gěi háizi gèng duō de ài; qícì, xuéxiào yīnggāi jiāo xuésheng zūnshǒu fǎlǜ; zuìhòu, shèhuì yīnggāi bāngzhù yǒu kùnnan de niánqīngrén. Wǒ xiāngxìn, zhǐyào wǒmen yìqǐ nǔlì, míngtiān yídìng huì gèng hǎo.}
\end{exercice}

\begin{exercice}[Expression écrite et dictée]
\textbf{A. Expression écrite guidée.} Rédige un texte argumentatif de \textbf{60 à 80 caractères chinois} sur le sujet : « Comment aider les jeunes de mon quartier ? » (怎么帮助我们社区的年轻人？)

Consignes obligatoires :
\begin{itemize}
\item Suis le plan : 1) présenter le problème ; 2) donner une cause ; 3) proposer deux ou trois solutions ; 4) conclure avec un souhait.
\item Utilise au moins deux connecteurs parmi : 首先、其次、最后、因为、所以、只要。
\item Utilise au moins cinq mots de la banque : 青少年、关心、教育、法律、帮助、努力、解决、机会。
\end{itemize}

\textbf{B. Dictée.} Écris en caractères les dix mots dictés par le professeur : 青少年、犯罪、严重、缺少、法律、社会、解决、错误、关心、努力. L'exactitude des traits et l'ordre d'écriture seront notés.
\end{exercice}

\newpage

\coursec{Corrigés des exercices}

\coursec{Corrigés des exercices — Leçon 8 : La délinquance juvénile}

\courssub{Exercices de vocabulaire et grammaire (Exercices 1 à 4)}

\begin{propriete}[Rappel : ordre des mots avec les verbes modaux]
En chinois, le verbe modal (应该、能、想、必须…) se place \textbf{avant} le verbe d'action principal.
\zh{Sujet + VerbeModal + Verbe + (Objet)} \\
Exemple : \zh{我们应该遵守法律。} \emph{Wǒmen yīnggāi zūnshǒu fǎlǜ.} (Nous devons respecter la loi.)
\end{propriete}

\begin{corrige}[Exercice 1 — QCM : le bon choix]
\begin{enumerate}
\item \textbf{Réponse b}. \zh{犯罪} (\emph{fànzuì}) signifie « commettre un crime » (verbe) ou « le crime, la délinquance » (nom). « La loi » se dit \zh{法律} (\emph{fǎlǜ}) et « la prison » \zh{监狱} (\emph{jiānyù}).
\item \textbf{Réponse a}. \zh{缺少} (\emph{quēshǎo}) signifie « manquer de » (verbe transitif direct, sans préposition). Exemple : \zh{缺少家庭教育} (\emph{quēshǎo jiātíng jiàoyù}), « manquer d'éducation familiale ». « Posséder » se dit \zh{拥有} (\emph{yōngyǒu}) et « partager » \zh{分享} (\emph{fēnxiǎng}).
\item \textbf{Réponse a}. \zh{我们应该关心青少年。} (\emph{Wǒmen yīnggāi guānxīn qīngshàonián.}) « Nous devons nous préoccuper des adolescents. » L'ordre correct est : Sujet (\zh{我们}) + Verbe modal (\zh{应该}) + Verbe (\zh{关心}) + Complément d'objet (\zh{青少年}). Les propositions b et c violent l'ordre canonique SVO du chinois.
\item \textbf{Réponse a}. La collocation correcte est \zh{解决问题} (\emph{jiějué wèntí}), « résoudre un problème ». \zh{解决} ne s'utilise pas avec « repas » ni « temps qu'il fait ».
\end{enumerate}
\end{corrige}

\begin{corrige}[Exercice 2 — Vrai ou faux justifié]
\begin{enumerate}
\item \textbf{FAUX}. \zh{青少年} (\emph{qīngshàonián}) désigne les adolescents et les jeunes (généralement 12--18 ans), pas les personnes âgées. Les personnes âgées se disent \zh{老年人} (\emph{lǎoniánrén}).
\item \textbf{VRAI}. \zh{严重} (\emph{yánzhòng}) est bien un adjectif signifiant « grave, sérieux ». Exemple : \zh{问题很严重。} (\emph{Wèntí hěn yánzhòng.}) « Le problème est très grave. »
\item \textbf{FAUX}. \zh{教育} (\emph{jiàoyù}) signifie « éduquer, enseigner », et non « punir ». « Punir » se dit \zh{惩罚} (\emph{chéngfá}). La phrase \zh{家长应该教育孩子} signifie « Les parents doivent éduquer leurs enfants ».
\item \textbf{VRAI}. \zh{法律} (\emph{fǎlǜ}) signifie « la loi, le droit ». Collocation importante : \zh{遵守法律} (\emph{zūnshǒu fǎlǜ}), « respecter la loi / se conformer à la loi ».
\end{enumerate}
\end{corrige}

\begin{corrige}[Exercice 3 — Vocabulaire : complète le tableau]
\begin{center}
\begin{tabularx}{\linewidth}{|X|X|X|X|}
\hline
\rowcolor{popblueL} Caractères & Pinyin & Nature & Traduction \\
\hline
社会 & shèhuì & nom & la société \\
\hline
错误 & cuòwù & nom & l'erreur, la faute \\
\hline
努力 & nǔlì & verbe / adjectif & faire des efforts ; appliqué, travailleur \\
\hline
帮助 & bāngzhù & verbe / nom & aider ; l'aide \\
\hline
原因 & yuányīn & nom & la cause, la raison \\
\hline
\end{tabularx}
\end{center}
\textbf{Remarques} :
\begin{itemize}
\item \zh{努力} : verbe (\zh{我们应该努力。} \emph{Wǒmen yīnggāi nǔlì.} « Nous devons faire des efforts ») ou adjectif (\zh{他很努力。} \emph{Tā hěn nǔlì.} « Il est très travailleur »).
\item \zh{帮助} : verbe (\zh{帮助他们} \emph{bāngzhù tāmen}, « les aider ») et nom (\zh{谢谢你的帮助。} \emph{Xièxie nǐ de bāngzhù.} « Merci de ton aide »).
\end{itemize}
\end{corrige}

\begin{corrige}[Exercice 4 — Pinyin et caractères : associe]
\begin{itemize}
\item 1 — d : \zh{犯} = \emph{fàn} (commettre ; criminel). Attention au 4\textsuperscript{e} ton (descendant).
\item 2 — f : \zh{青} = \emph{qīng} (jeune, vert/bleu). Attention à la nasale finale \emph{-ing}.
\item 3 — b : \zh{法} = \emph{fǎ} (loi, droit). 3\textsuperscript{e} ton (bas-courbe).
\item 4 — a : \zh{缺} = \emph{quē} (manquer). Le \emph{u} après \emph{q} se prononce \emph{[y]} (ü).
\item 5 — e : \zh{严} = \emph{yán} (sévère, grave). 2\textsuperscript{e} ton (montant).
\item 6 — c : \zh{决} = \emph{jué} (décider, résoudre). Le \emph{u} après \emph{j} se prononce \emph{[y]} (ü).
\end{itemize}
\textbf{Point de vigilance} : après les consonnes \emph{j, q, x}, la voyelle \emph{u} se prononce toujours \emph{[y]} (ü) ; on écrit \emph{quē} et \emph{jué} sans tréma selon la norme Pinyin, mais la prononciation reste celle de \emph{ü}.
\end{corrige}

\courssub{Exercices de traduction et discrimination (Exercices 5 à 6)}

\begin{attention}[Pièges fréquents pour francophones]
\begin{itemize}
\item \textbf{犯罪 vs 犯错误} : \zh{犯罪} (\emph{fànzuì}) = « commettre un crime/délit » (terme juridique fort). \zh{犯错误} (\emph{fàn cuòwù}) = « faire une erreur, commettre une faute » (terme courant, scolaire, moral). Ne pas confondre.
\item \textbf{遵守 vs 尊重} : \zh{遵守} (\emph{zūnshǒu}) = « respecter, observer (une loi, une règle, un horaire) ». \zh{尊重} (\emph{zūnzhòng}) = « respecter (une personne, un choix, un sentiment) ». On dit \zh{遵守法律}, jamais \zh{尊重法律}.
\item \textbf{缺少} (verbe transitif direct) : pas de préposition après. \zh{缺少钱} (manquer d'argent), non \zh{缺少 *de* 钱}.
\item \textbf{应该} placement : Toujours \textbf{avant} le verbe. \zh{应该去}, non \zh{去应该}.
\end{itemize}
\end{attention}

\begin{corrige}[Exercice 5 — Thème : traduis en chinois]
\begin{enumerate}
\item \zh{一些年轻人缺少家庭教育。} \\
\emph{Yìxiē niánqīngrén quēshǎo jiātíng jiàoyù.} \\
Justification : « certains » = \zh{一些} + nom ; « manquer de » = \zh{缺少} + nom (verbe transitif direct, sans préposition).
\item \zh{青少年犯罪是一个严重的问题。} \\
\emph{Qīngshàonián fànzuì shì yí ge yánzhòng de wèntí.} \\
Justification : \zh{是} + classificateur \zh{个} + adjectif (\zh{严重}) + \zh{的} + nom. Ne pas oublier \zh{的} entre l'adjectif à deux syllabes et le nom.
\item \zh{家长和学校应该帮助这些年轻人。} \\
\emph{Jiāzhǎng hé xuéxiào yīnggāi bāngzhù zhèxiē niánqīngrén.} \\
Justification : « et » entre deux noms = \zh{和} ; « devraient » = \zh{应该} placé avant le verbe \zh{帮助}.
\item \zh{我们应该遵守法律。} \\
\emph{Wǒmen yīnggāi zūnshǒu fǎlǜ.} \\
Justification : « respecter (la loi/règle) » = \zh{遵守} ; on évite \zh{尊重} (\emph{zūnzhòng}) qui s'emploit pour une personne.
\end{enumerate}
\end{corrige}

\begin{corrige}[Exercice 6 — Discrimination de caractères]
\begin{enumerate}
\item \zh{青}少年犯罪是一个严重的问题。 \zh{青} (\emph{qīng}) = jeune (composé de \zh{月} et \zh{八}) ; \zh{清} (\emph{qīng}) = clair, pur (avec clé de l'eau \zh{氵}). C'est \zh{青} qui convient pour « jeunesse ».
\item 他\textbf{犯}了错误，应该改正。 \zh{犯} (\emph{fàn}, clé du chien \zh{犭}) = commettre ; \zh{范} (\emph{fàn}) = modèle (ex: \zh{模范} \emph{mófàn}). La collocation est \zh{犯错误} (\emph{fàn cuòwù}), « commettre une erreur ».
\item 每个人都应该遵守\textbf{法}律。 \zh{法} (\emph{fǎ}, clé de l'eau \zh{氵}) = loi ; \zh{去} (\emph{qù}) = aller. \zh{法律} = « la loi ».
\item 水很\textbf{清}，我可以看见鱼。 \zh{清} (\emph{qīng}, avec la clé de l'eau \zh{氵}) = clair, limpide ; c'est logique : l'eau claire porte la clé de l'eau.
\end{enumerate}
\textbf{Astuce mémoire} : le radical (clé) donne souvent l'indice sémantique : \zh{清} a la clé de l'eau \zh{氵} (eau claire) ; \zh{犯} a la clé du chien \zh{犭} (autrefois associé aux actes violents/animaux) ; \zh{法} contient \zh{氵} (la loi doit être « plane comme l'eau » / impartiale) et \zh{去} (phonétique).
\end{corrige}

\courssub{Exercices de cohérence textuelle (Exercices 7 à 9)}

\begin{propriete}[Connecteurs logiques essentiels pour l'argumentation]
\begin{center}
\begin{tabularx}{\linewidth}{|X|X|X|}
\hline
\rowcolor{popblueL} Fonction & Connecteur chinois (Pinyin) & Français \\
\hline
Cause (antéposée) & \zh{因为} (\emph{yīnwèi}) & Parce que \\
Conséquence (conséquente) & \zh{所以} (\emph{suǒyǐ}) & Donc, c'est pourquoi \\
\hline
Énumération 1 & \zh{首先} (\emph{shǒuxiān}) & D'abord, en premier lieu \\
Énumération 2 & \zh{其次} (\emph{qícì}) & Ensuite, en second lieu \\
Énumération 3 (fin) & \zh{最后} (\emph{zuìhòu}) & Enfin, en dernier lieu \\
\hline
Condition suffisante & \zh{只要} (\emph{zhǐyào}) & Pourvu que, tant que \\
Résultat certain & \zh{就} (\emph{jiù}) / \zh{一定} (\emph{yídìng}) & Alors / certainement \\
\hline
\end{tabularx}
\end{center}
\textbf{Règle d'or} : En chinois, la paire \zh{因为} \ldots \zh{所以} \ldots est souvent \textbf{obligatoire} dans la phrase complexe (on garde les deux), contrairement au français où « parce que » suffit souvent seul.
\end{propriete}

\begin{corrige}[Exercice 7 — Texte à trous : mots-outils]
Texte complété :
\zh{因为 一些青少年缺少家庭教育，所以 他们会犯错误。怎么解决这个问题？首先，家长应该多关心孩子；其次，学校应该加强教育；最后，社会也应该帮助他们。} \\
\emph{Yīnwèi yīxiē qīngshàonián quēshǎo jiātíng jiàoyù, suǒyǐ tāmen huì fàn cuòwù. Zěnme jiějué zhège wèntí? Shǒuxiān, jiāzhǎng yīnggāi duō guānxīn háizi; qícì, xuéxiào yīnggāi jiāqiáng jiàoyù; zuìhòu, shèhuì yě yīnggāi bāngzhù tāmen.}

Traduction : « Parce que certains adolescents manquent d'éducation familiale, ils peuvent commettre des erreurs. Comment résoudre ce problème ? D'abord, les parents doivent se préoccuper davantage de leurs enfants ; ensuite, l'école doit renforcer l'éducation ; enfin, la société doit aussi les aider. »

\textbf{Justifications} :
\begin{itemize}
\item \zh{因为} \ldots \zh{所以} \ldots : paire obligatoire de cause et conséquence ; en chinois, on utilise souvent les DEUX termes.
\item \zh{首先} / \zh{其次} / \zh{最后} : connecteurs d'énumération des solutions (d'abord / ensuite / enfin), placés en tête de phrase suivis d'une virgule.
\item Le mot \zh{应该} de la banque n'est pas utilisé dans les trous : il apparaît déjà dans le texte original ; c'est un distracteur.
\end{itemize}
\end{corrige}

\begin{corrige}[Exercice 8 — Remise en ordre]
Ordre correct : \textbf{b — d — e — c — a}, c'est-à-dire :
\begin{enumerate}
\item \zh{现在，青少年犯罪是一个严重的问题。} (présentation du problème)
\item \zh{怎么解决这个问题呢？} (question de transition, marqueur \zh{呢})
\item \zh{首先，家长应该多关心孩子，多和他们交流。} (solution 1 : \zh{首先})
\item \zh{其次，学校应该加强法律和道德教育。} (solution 2 : \zh{其次})
\item \zh{最后，社会应该给年轻人更多的机会。} (solution 3 : \zh{最后})
\end{enumerate}
\textbf{Méthode} : repérer les connecteurs structurants. \zh{首先} ouvre toujours l'énumération, \zh{其次} la continue, \zh{最后} la ferme. La phrase avec \zh{呢} (\zh{怎么}\ldots\zh{呢？}) est une question rhétorique ou réelle : elle se place logiquement après le constat et avant les solutions.

Traduction du texte reconstitué : « Aujourd'hui, la délinquance juvénile est un problème grave. Comment résoudre ce problème ? D'abord, les parents doivent se préoccuper davantage de leurs enfants et communiquer plus avec eux. Ensuite, l'école doit renforcer l'éducation juridique et morale. Enfin, la société doit donner plus d'opportunités aux jeunes. »
\end{corrige}

\begin{corrige}[Exercice 9 — Compréhension de texte et questions de langue]
\textbf{Questions de compréhension} :
\begin{enumerate}
\item Le texte évoque le problème de la délinquance juvénile (\zh{青少年犯罪}), qui devient de plus en plus grave dans certaines villes.
\item Deux causes sont mentionnées : (1) certains jeunes manquent d'éducation familiale (\zh{缺少家庭教育}) ; (2) d'autres n'ont pas de travail et vivent dans la difficulté (\zh{没有工作，生活很困难}).
\item Les trois solutions : \zh{首先}, les parents doivent se préoccuper davantage de leurs enfants et communiquer avec eux (rôle des \zh{家长}) ; \zh{其次}, l'école doit renforcer l'éducation juridique pour que les élèves distinguent le bien du mal (rôle de l'\zh{学校}) ; \zh{最后}, le gouvernement et la société doivent offrir plus d'opportunités d'emploi aux jeunes (rôle du \zh{政府} et de la \zh{社会}).
\end{enumerate}
\textbf{Questions de langue} :
\begin{enumerate}
\item Les trois connecteurs d'énumération : \zh{首先} (\emph{shǒuxiān}) « d'abord », \zh{其次} (\emph{qícì}) « ensuite », \zh{最后} (\emph{zuìhòu}) « enfin ».
\item Structure \zh{只要} \ldots \zh{就/一定} \ldots : \zh{只要} (\emph{zhǐyào}) introduit la condition suffisante (« pourvu que, tant que »), et la principale exprime le résultat certain (ici renforcé par \zh{一定}, « certainement »). Schéma : \zh{只要} + condition, \zh{Sujet} + \zh{一定/就} + résultat. Traduction : « Pourvu que tout le monde fasse des efforts ensemble, ce problème pourra certainement être résolu. »
\item Adjectif « grave » : \zh{严重} (\emph{yánzhòng}) ; verbe « résoudre » : \zh{解决} (\emph{jiějué}).
\end{enumerate}
\end{corrige}

\courssub{Évaluation type Bac : Thème, Version, Expression écrite (Exercices 10 à 11)}

\begin{methode}[Réussir le thème (Français → Chinois) et la version (Chinois → Français)]
\begin{enumerate}
\item \textbf{Analyser} la phrase française : identifier le sujet, le verbe, les compléments, les temps, les connecteurs logiques.
\item \textbf{Transposer} la structure : appliquer l'ordre chinois (Sujet -- Temps/Aspect -- Verbe Modal -- Verbe -- Objet ; Compléments de lieu/temps \textbf{avant} le verbe ; \zh{的} après adjectif long).
\item \textbf{Vérifier} les collocations (ex: \zh{遵守法律}, \zh{犯错误}, \zh{关心}) et les classificateurs.
\item \textbf{Écrire} les caractères corrects (traits, radicaux) et la ponctuation pleine chasse (\zh{，} \zh{。} \zh{？} \zh{；}).
\item \textbf{Relire} : cohérence du sens, tons du pinyin (mentalement), accord nombre/classificateur.
\end{enumerate}
\end{methode}

\begin{corrige}[Exercice 10 — Thème et version type Bac]
\textbf{Thème — Proposition de traduction} :\\
\zh{因为有些年轻人缺少家庭教育，所以他们会犯错误。学校和家长应该帮助他们。} \\
\emph{Yīnwèi yǒuxiē niánqīngrén quēshǎo jiātíng jiàoyù, suǒyǐ tāmen huì fàn cuòwù. Xuéxiào hé jiāzhǎng yīnggāi bāngzhù tāmen.}

\textbf{Points de vigilance} :
\begin{itemize}
\item Conserver la paire \zh{因为} \ldots \zh{所以} \ldots (structure cause-conséquence).
\item Collocation \zh{犯错误} (\emph{fàn cuòwù}) pour « commettre des erreurs » (et non \zh{犯罪}).
\item \zh{应该} placé \textbf{avant} le verbe \zh{帮助}.
\item \zh{帮助他们} : verbe + objet (pronom \zh{他们} après le verbe).
\end{itemize}

\textbf{Version — Proposition de traduction} :\\
« Comment résoudre le problème de la délinquance juvénile ? D'abord, la famille doit donner plus d'amour aux enfants ; ensuite, l'école doit apprendre aux élèves à respecter la loi ; enfin, la société doit aider les jeunes en difficulté. Je crois que, pourvu que nous fassions des efforts ensemble, demain sera certainement meilleur. »

\textbf{Barème indicatif (sur 10)} :
\begin{itemize}
\item Exactitude du lexique et des collocations (3 pts)
\item Structures grammaticales correctes — ordre des mots, connecteurs, \zh{只要}\ldots\zh{一定} (3 pts)
\item Exactitude des caractères pour le thème / fidélité du sens pour la version (2 pts)
\item Cohérence, ponctuation chinoise pleine chasse, présentation (2 pts)
\end{itemize}
\end{corrige}

\begin{corrige}[Exercice 11 — Expression écrite et dictée]
\textbf{A. Proposition de rédaction modèle} (environ 75 caractères chinois) :
\zh{怎么帮助我们社区的年轻人？} \\
\zh{现在，我们社区有些年轻人没有工作，生活很困难，有的还会犯错误。怎么帮助他们呢？首先，家长应该多关心孩子；其次，学校应该加强法律教育；最后，政府应该给年轻人更多的工作机会。只要大家一起努力，我们的社区一定会更好。} \\
\emph{Xiànzài, wǒmen shèqū yǒuxiē niánqīngrén méiyǒu gōngzuò, shēnghuó hěn kùnnan, yǒude hái huì fàn cuòwù. Zěnme bāngzhù tāmen ne? Shǒuxiān, jiāzhǎng yīnggāi duō guānxīn háizi; qícì, xuéxiào yīnggāi jiāqiáng fǎlǜ jiàoyù; zuìhòu, zhèngfǔ yīnggāi gěi niánqīngrén gèng duō de gōngzuò jīhuì. Zhǐyào dàjiā yìqǐ nǔlì, wǒmen de shèqū yídìng huì gèng hǎo.}

Traduction : « Aujourd'hui, dans notre quartier, certains jeunes n'ont pas de travail, vivent dans la difficulté, et certains peuvent même commettre des erreurs. Comment les aider ? D'abord, les parents doivent se préoccuper davantage de leurs enfants ; ensuite, l'école doit renforcer l'éducation juridique ; enfin, le gouvernement doit offrir plus d'opportunités d'emploi aux jeunes. Pourvu que tout le monde fasse des efforts ensemble, notre quartier sera certainement meilleur. »

\textbf{Barème indicatif (sur 20)} :
\begin{itemize}
\item Respect du plan en 4 étapes (constat, question, 3 solutions, conclusion) : 4 pts
\item Au moins deux connecteurs correctement employés (\zh{首先}/\zh{其次}/\zh{最后}, \zh{因为}/\zh{所以}, \zh{只要}\ldots) : 4 pts
\item Au moins cinq mots de la banque de vocabulaire utilisés correctement : 4 pts
\item Exactitude des caractères, grammaire (ordre SVO, \zh{的}, \zh{应该} placement) : 6 pts
\item Longueur respectée (60--80 caractères chinois) : 2 pts
\end{itemize}

\textbf{Points de vigilance rédactionnels} :
\begin{itemize}
\item Ordre Sujet -- Verbe -- Objet strict.
\item \zh{应该} toujours \textbf{avant} le verbe.
\item \zh{的} obligatoire après adjectif bi-syllabique + nom (\zh{更多的机会}, \zh{严重的问题}).
\item Ne pas traduire mot à mot le français (ex: « donner une opportunité » $\to$ \zh{给机会}, non \zh{给一个机会} sans classificateur si générique, mais \zh{给更多的机会} est correct).
\item Ponctuation chinoise pleine chasse : \zh{，} \zh{。} \zh{；} \zh{？} (pas d'espaces avant/après).
\end{itemize}

\textbf{B. Dictée — Rappels d'écriture des caractères clés (ordre des traits et radicaux)} :
\begin{itemize}
\item \zh{青少年} : \zh{青} (horizontales supérieures, puis \zh{月} en bas) ; \zh{少} (petit \zh{小} au-dessus de \zh{丿}).
\item \zh{犯罪} : \zh{犯} (clé du chien \zh{犭} à gauche, 5 traits) ; \zh{罪} (clé du filet \zh{罒} en haut, \zh{非} en bas).
\item \zh{严重} : \zh{严} (haut \zh{丷}, bas \zh{口} \zh{口} \zh{工}) ; \zh{重} (horizontales, verticale centrale, \zh{里}).
\item \zh{缺少} : \zh{缺} (clé \zh{缶} à gauche, phonétique \zh{戈} \zh{欠} à droite) ; \zh{少} (cf. ci-dessus).
\item \zh{法律} : \zh{法} (clé de l'eau \zh{氵} à gauche, \zh{去} à droite) ; \zh{律} (clé \zh{彳} double personne, \zh{聿}).
\item \zh{社会} : \zh{社} (clé de l'autel \zh{礻} -- 4 traits : point, horizontale, verticale, point -- à ne pas confondre avec \zh{衤} vêtement) ; \zh{会} (clé \zh{曰} sous \zh{云}).
\item \zh{解决} : \zh{解} (clé \zh{角} \zh{刀} \zh{牛}) ; \zh{决} (clé \zh{氵} eau, \zh{夬}).
\item \zh{错误} : \zh{错} (clé du métal \zh{钅} à gauche, \zh{昔} \zh{掌} à droite) ; \zh{误} (clé de la parole \zh{讠}, \zh{吴}).
\item \zh{关心} : \zh{关} (clé \zh{门} porte, \zh{关} dedans) ; \zh{心} (cœur, 3 traits : points et crochet).
\item \zh{努力} : \zh{努} (haut \zh{奴}, bas \zh{力}) ; \zh{力} (crochet final).
\end{itemize}
\textbf{Règles générales d'ordre des traits} : de haut en bas (\zh{三}), de gauche à droite (\zh{川}), l'horizontale avant la verticale (\zh{十}), l'extérieur avant l'intérieur (\zh{口} \zh{日}), le trait central avant les symétriques (\zh{小} \zh{水}).
\end{corrige}

\newpage

\coursec{Fiche bilan}

\begin{popfigure}
\begin{tikzpicture}[mindmap, grow cyclic, every node/.style={concept, align=center, font=\small},
  root concept/.append style={concept color=popblue, font=\small\bfseries},
  level 1/.append style={level distance=3.1cm, sibling angle=51.4},
  level 1 concept/.append style={font=\scriptsize}]
\node[root concept, align=center]{La délinquance\\juvénile\\青少年犯罪}
  child[concept color=poporange]{ node[align=center]{Lexique\\犯罪 青少年\\法律 监狱} }
  child[concept color=popgreen]{ node[align=center]{Causes\\因为 由于\\家庭 学校} }
  child[concept color=poppurple]{ node[align=center]{Conséquences\\所以 因此\\影响 社会} }
  child[concept color=poppink]{ node[align=center]{Solutions\\应该 必须\\教育 帮助} }
  child[concept color=popgold]{ node[align=center]{Connecteurs\\首先 其次\\最后 总之} }
  child[concept color=popblue!60]{ node[align=center]{Structures\\越来越\\不仅…而且…} }
  child[concept color=popgreen!60]{ node[align=center]{Écriture\\radicaux\\讠 氵 忄} };
\end{tikzpicture}
\legende{Carte mentale de la leçon : les sept compétences à mobiliser pour écrire un texte argumentatif sur la délinquance juvénile.}
\end{popfigure}

\begin{bilan}
\textbf{1. Lexique clé du thème}\par
\begin{itemize}
  \item \zh{青少年} qīngshàonián : adolescent, jeunesse ; \zh{犯罪} fànzuì : commettre un crime ; \zh{罪犯} zuìfàn : criminel
  \item \zh{法律} fǎlǜ : loi ; \zh{监狱} jiānyù : prison ; \zh{警察} jǐngchá : police
  \item \zh{原因} yuányīn : cause ; \zh{影响} yǐngxiǎng : influence, effet ; \zh{解决} jiějué : résoudre
  \item \zh{教育} jiàoyù : éducation ; \zh{保护} bǎohù : protéger ; \zh{帮助} bāngzhù : aider
  \item Collocations : \zh{预防犯罪} yùfáng fànzuì « prévenir la délinquance » ; \zh{遵守法律} zūnshǒu fǎlǜ « respecter la loi »
\end{itemize}

\textbf{2. Grammaire à connaître par cœur}\par
\begin{center}
\begin{tabularx}{\linewidth}{|X|X|X|}
\hline
\rowcolor{popblueL} Structure & Exemple & Traduction \\
\hline
越来越 + adj. & 问题越来越严重。Wèntí yuèláiyuè yánzhòng. & Le problème est de plus en plus grave. \\
\hline
因为…所以… & 因为家庭问题，所以他犯了罪。Yīnwèi jiātíng wèntí, suǒyǐ tā fàn le zuì. & À cause de problèmes familiaux, il a commis un délit. \\
\hline
不仅…而且… & 这不仅影响家庭，而且影响社会。Zhè bùjǐn yǐngxiǎng jiātíng, érqiě yǐngxiǎng shèhuì. & Cela n'affecte pas seulement la famille, mais aussi la société. \\
\hline
应该 / 必须 + verbe & 我们应该帮助这些青少年。Wǒmen yīnggāi bāngzhù zhèxiē qīngshàonián. & Nous devrions aider ces adolescents. \\
\hline
为了 + but & 为了解决这个问题… Wèile jiějué zhège wèntí… & Pour résoudre ce problème… \\
\hline
\end{tabularx}
\end{center}

\textbf{3. Connecteurs du texte argumentatif}\par
\zh{首先} shǒuxiān « d'abord » → \zh{其次} qícì « ensuite » → \zh{最后} zuìhòu « enfin » → \zh{总之} zǒngzhī « en conclusion » ; opposition : \zh{但是} dànshì, \zh{可是} kěshì.

\textbf{4. Méthodes express}\par
\begin{enumerate}
  \item \textbf{Thème (français → chinois)} : identifier sujet + verbe ; placer le complément de temps et de lieu AVANT le verbe ; vérifier les tons et le classificateur.
  \item \textbf{Version (chinois → français)} : repérer les connecteurs (因为, 所以, 但是) ; traduire groupe par groupe ; accorder en français.
  \item \textbf{Production écrite} : plan en 3 parties (causes → conséquences → solutions) ; une idée = une phrase courte ; conclure avec 总之.
  \item \textbf{Caractères} : écrire de haut en bas, de gauche à droite, l'horizontal avant le vertical ; repérer les radicaux (讠parole, 氵eau, 忄cœur).
\end{enumerate}
\end{bilan}

\begin{aretenir}[Checklist avant le Bac]
\begin{itemize}
  \item $\square$ Je connais par cœur le lexique : 青少年, 犯罪, 法律, 监狱, 原因, 影响, 解决.
  \item $\square$ Je sais écrire les caractères du thème avec le bon ordre des traits.
  \item $\square$ Je distingue les caractères proches : 犯 / 饭, 法 / 去, 罪 / 非.
  \item $\square$ Je place le complément de temps et de lieu avant le verbe.
  \item $\square$ J'utilise correctement 因为…所以… et 不仅…而且… dans une même phrase.
  \item $\square$ Je sais exprimer la cause (由于), le but (为了) et la conséquence (因此).
  \item $\square$ Mon texte suit le plan : introduction → causes → conséquences → solutions → conclusion avec 总之.
  \item $\square$ Je mets le classificateur correct devant les noms (一个问题, 这些青少年).
  \item $\square$ J'écris le pinyin avec les tons sur les voyelles, sans chiffres.
  \item $\square$ Je relis ma copie : tons, ordre des mots, ponctuation chinoise ， 。 ？
  \item $\square$ En version, je traduis groupe par groupe et j'accorde en français.
  \item $\square$ Je respecte le nombre de phrases demandé et je reste dans le thème.
\end{itemize}
\end{aretenir}

\findecours

\end{document}
